% !TEX program = pdflatex

\documentclass{article} % For LaTeX2e
\usepackage{iclr2027_conference,times}

% Optional math commands from https://github.com/goodfeli/dlbook_notation.
\input{math_commands.tex}

\usepackage{hyperref}
\usepackage{url}
\usepackage{amsmath,amssymb}
\usepackage{booktabs,tabularx,array,multirow}
% Japanese text under Overleaf's default pdfLaTeX compiler.
\usepackage[utf8]{inputenc}
\usepackage{CJKutf8}


\title{Formatting Instructions for ICLR 2027 \\ Conference Submissions}

% Authors must not appear in the submitted version. They should be hidden
% as long as the \iclrfinalcopy macro remains commented out below.
% Non-anonymous submissions will be rejected without review.

\author{Antiquus S.~Hippocampus, Natalia Cerebro \& Amelie P. Amygdale \thanks{ Use footnote for providing further information
about author (webpage, alternative address)---\emph{not} for acknowledging
funding agencies.  Funding acknowledgements go at the end of the paper.} \\
Department of Computer Science\\
Cranberry-Lemon University\\
Pittsburgh, PA 15213, USA \\
\texttt{\{hippo,brain,jen\}@cs.cranberry-lemon.edu} \\
\And
Ji Q. Ren \& Yevgeny LeNet \\
Department of Computational Neuroscience \\
University of the Witwatersrand \\
Joburg, South Africa \\
\texttt{\{robot,net\}@wits.ac.za} \\
\AND
Coauthor \\
Affiliation \\
Address \\
\texttt{email}
}

% The \author macro works with any number of authors. There are two commands
% used to separate the names and addresses of multiple authors: \And and \AND.
%
% Using \And between authors leaves it to \LaTeX{} to determine where to break
% the lines. Using \AND forces a linebreak at that point. So, if \LaTeX{}
% puts 3 of 4 authors names on the first line, and the last on the second
% line, try using \AND instead of \And before the third author name.

\newcommand{\fix}{\marginpar{FIX}}
\newcommand{\new}{\marginpar{NEW}}

%\iclrfinalcopy % Uncomment for camera-ready version, but NOT for submission.
\begin{document}
\begin{CJK*}{UTF8}{min}


\maketitle

\begin{abstract}
The abstract paragraph should be indented 1/2~inch (3~picas) on both left and
right-hand margins. Use 10~point type, with a vertical spacing of 11~points.
The word \textsc{Abstract} must be centered, in small caps, and in point size 12. Two
line spaces precede the abstract. The abstract must be limited to one
paragraph.
\end{abstract}

\section{Submission of conference papers to ICLR 2027}

ICLR requires electronic submissions, processed by
\url{https://openreview.net/}. See ICLR's website for more instructions.

If your paper is ultimately accepted, the statement {\tt
  {\textbackslash}iclrfinalcopy} should be inserted to adjust the
format to the camera ready requirements.

The format for the submissions is a variant of the NeurIPS format.
Please read carefully the instructions below, and follow them
faithfully.

\subsection{Style}

Papers to be submitted to ICLR 2027 must be prepared according to the
instructions presented here.

%% Please note that we have introduced automatic line number generation
%% into the style file for \LaTeXe. This is to help reviewers
%% refer to specific lines of the paper when they make their comments. Please do
%% NOT refer to these line numbers in your paper as they will be removed from the
%% style file for the final version of accepted papers.

Authors are required to use the ICLR \LaTeX{} style files obtainable at the
ICLR website. Please make sure you use the current files and
not previous versions. Tweaking the style files may be grounds for rejection.

\subsection{Retrieval of style files}

The style files for ICLR and other conference information are available online at:
\begin{center}
   \url{http://www.iclr.cc/}
\end{center}
The file \verb+iclr2027_conference.pdf+ contains these
instructions and illustrates the
various formatting requirements your ICLR paper must satisfy.
Submissions must be made using \LaTeX{} and the style files
\verb+iclr2027_conference.sty+ and \verb+iclr2027_conference.bst+ (to be used with \LaTeX{}2e). The file
\verb+iclr2027_conference.tex+ may be used as a ``shell'' for writing your paper. All you
have to do is replace the author, title, abstract, and text of the paper with
your own.

The formatting instructions contained in these style files are summarized in
sections \ref{gen_inst}, \ref{headings}, and \ref{others} below.

\section{General formatting instructions}
\label{gen_inst}

The text must be confined within a rectangle 5.5~inches (33~picas) wide and
9~inches (54~picas) long. The left margin is 1.5~inch (9~picas).
Use 10~point type with a vertical spacing of 11~points. Times New Roman is the
preferred typeface throughout. Paragraphs are separated by 1/2~line space,
with no indentation.

Paper title is 17~point, in small caps and left-aligned.
All pages should start at 1~inch (6~picas) from the top of the page.

Authors' names are
set in boldface, and each name is placed above its corresponding
address. The lead author's name is to be listed first, and
the co-authors' names are set to follow. Authors sharing the
same address can be on the same line.

Please pay special attention to the instructions in section \ref{others}
regarding figures, tables, acknowledgments, and references.


There will be a strict upper limit of \textbf{9 pages} for the main text of the initial submission, with unlimited additional pages for citations. This limit will be expanded to \textbf{10 pages} for rebuttal/camera ready.

\section{Headings: first level}
\label{headings}

First level headings are in small caps,
flush left and in point size 12. One line space before the first level
heading and 1/2~line space after the first level heading.

\subsection{Headings: second level}

Second level headings are in small caps,
flush left and in point size 10. One line space before the second level
heading and 1/2~line space after the second level heading.

\subsubsection{Headings: third level}

Third level headings are in small caps,
flush left and in point size 10. One line space before the third level
heading and 1/2~line space after the third level heading.

\section{Citations, figures, tables, references}
\label{others}

These instructions apply to everyone, regardless of the formatter being used.

\subsection{Citations within the text}

Citations within the text should be based on the \texttt{natbib} package
and include the authors' last names and year (with the ``et~al.'' construct
for more than two authors). When the authors or the publication are
included in the sentence, the citation should not be in parenthesis using \verb|\citet{}| (as
in ``See \citet{Hinton06} for more information.''). Otherwise, the citation
should be in parenthesis using \verb|\citep{}| (as in ``Deep learning shows promise to make progress
towards AI~\citep{Bengio+chapter2007}.'').

The corresponding references are to be listed in alphabetical order of
authors, in the \textsc{References} section. As to the format of the
references themselves, any style is acceptable as long as it is used
consistently.

\subsection{Footnotes}

Indicate footnotes with a number\footnote{Sample of the first footnote} in the
text. Place the footnotes at the bottom of the page on which they appear.
Precede the footnote with a horizontal rule of 2~inches
(12~picas).\footnote{Sample of the second footnote}

\subsection{Figures}

All artwork must be neat, clean, and legible. Lines should be dark
enough for purposes of reproduction; art work should not be
hand-drawn. The figure number and caption always appear after the
figure. Place one line space before the figure caption, and one line
space after the figure. The figure caption is lower case (except for
first word and proper nouns); figures are numbered consecutively.

Make sure the figure caption does not get separated from the figure.
Leave sufficient space to avoid splitting the figure and figure caption.

You may use color figures.
However, it is best for the
figure captions and the paper body to make sense if the paper is printed
either in black/white or in color.
\begin{figure}[h]
\begin{center}
%\framebox[4.0in]{$\;$}
\fbox{\rule[-.5cm]{0cm}{4cm} \rule[-.5cm]{4cm}{0cm}}
\end{center}
\caption{Sample figure caption.}
\end{figure}

\subsection{Tables}

All tables must be centered, neat, clean and legible. Do not use hand-drawn
tables. The table number and title always appear before the table. See
Table~\ref{sample-table}.

Place one line space before the table title, one line space after the table
title, and one line space after the table. The table title must be lower case
(except for first word and proper nouns); tables are numbered consecutively.

\begin{table}[t]
\caption{Sample table title}
\label{sample-table}
\begin{center}
\begin{tabular}{ll}
\multicolumn{1}{c}{\bf PART}  &\multicolumn{1}{c}{\bf DESCRIPTION}
\\ \hline \\
Dendrite         &Input terminal \\
Axon             &Output terminal \\
Soma             &Cell body (contains cell nucleus) \\
\end{tabular}
\end{center}
\end{table}

\section{Default Notation}

In an attempt to encourage standardized notation, we have included the
notation file from the textbook, \textit{Deep Learning}
\cite{goodfellow2016deep} available at
\url{https://github.com/goodfeli/dlbook_notation/}.  Use of this style
is not required and can be disabled by commenting out
\texttt{math\_commands.tex}.


\centerline{\bf Numbers and Arrays}
\bgroup
\def\arraystretch{1.5}
\begin{tabular}{p{1in}p{3.25in}}
$\displaystyle a$ & A scalar (integer or real)\\
$\displaystyle \va$ & A vector\\
$\displaystyle \mA$ & A matrix\\
$\displaystyle \tA$ & A tensor\\
$\displaystyle \mI_n$ & Identity matrix with $n$ rows and $n$ columns\\
$\displaystyle \mI$ & Identity matrix with dimensionality implied by context\\
$\displaystyle \ve^{(i)}$ & Standard basis vector $[0,\dots,0,1,0,\dots,0]$ with a 1 at position $i$\\
$\displaystyle \text{diag}(\va)$ & A square, diagonal matrix with diagonal entries given by $\va$\\
$\displaystyle \ra$ & A scalar random variable\\
$\displaystyle \rva$ & A vector-valued random variable\\
$\displaystyle \rmA$ & A matrix-valued random variable\\
\end{tabular}
\egroup
\vspace{0.25cm}

\centerline{\bf Sets and Graphs}
\bgroup
\def\arraystretch{1.5}

\begin{tabular}{p{1.25in}p{3.25in}}
$\displaystyle \sA$ & A set\\
$\displaystyle \R$ & The set of real numbers \\
$\displaystyle \{0, 1\}$ & The set containing 0 and 1 \\
$\displaystyle \{0, 1, \dots, n \}$ & The set of all integers between $0$ and $n$\\
$\displaystyle [a, b]$ & The real interval including $a$ and $b$\\
$\displaystyle (a, b]$ & The real interval excluding $a$ but including $b$\\
$\displaystyle \sA \backslash \sB$ & Set subtraction, i.e., the set containing the elements of $\sA$ that are not in $\sB$\\
$\displaystyle \gG$ & A graph\\
$\displaystyle \parents_\gG(\ervx_i)$ & The parents of $\ervx_i$ in $\gG$
\end{tabular}
\vspace{0.25cm}


\centerline{\bf Indexing}
\bgroup
\def\arraystretch{1.5}

\begin{tabular}{p{1.25in}p{3.25in}}
$\displaystyle \eva_i$ & Element $i$ of vector $\va$, with indexing starting at 1 \\
$\displaystyle \eva_{-i}$ & All elements of vector $\va$ except for element $i$ \\
$\displaystyle \emA_{i,j}$ & Element $i, j$ of matrix $\mA$ \\
$\displaystyle \mA_{i, :}$ & Row $i$ of matrix $\mA$ \\
$\displaystyle \mA_{:, i}$ & Column $i$ of matrix $\mA$ \\
$\displaystyle \etA_{i, j, k}$ & Element $(i, j, k)$ of a 3-D tensor $\tA$\\
$\displaystyle \tA_{:, :, i}$ & 2-D slice of a 3-D tensor\\
$\displaystyle \erva_i$ & Element $i$ of the random vector $\rva$ \\
\end{tabular}
\egroup
\vspace{0.25cm}


\centerline{\bf Calculus}
\bgroup
\def\arraystretch{1.5}
\begin{tabular}{p{1.25in}p{3.25in}}
% NOTE: the [2ex] on the next line adds extra height to that row of the table.
% Without that command, the fraction on the first line is too tall and collides
% with the fraction on the second line.
$\displaystyle\frac{d y} {d x}$ & Derivative of $y$ with respect to $x$\\ [2ex]
$\displaystyle \frac{\partial y} {\partial x} $ & Partial derivative of $y$ with respect to $x$ \\
$\displaystyle \nabla_\vx y $ & Gradient of $y$ with respect to $\vx$ \\
$\displaystyle \nabla_\mX y $ & Matrix derivatives of $y$ with respect to $\mX$ \\
$\displaystyle \nabla_\tX y $ & Tensor containing derivatives of $y$ with respect to $\tX$ \\
$\displaystyle \frac{\partial f}{\partial \vx} $ & Jacobian matrix $\mJ \in \R^{m\times n}$ of $f: \R^n \rightarrow \R^m$\\
$\displaystyle \nabla_\vx^2 f(\vx)\text{ or }\mH( f)(\vx)$ & The Hessian matrix of $f$ at input point $\vx$\\
$\displaystyle \int f(\vx) d\vx $ & Definite integral over the entire domain of $\vx$ \\
$\displaystyle \int_\sS f(\vx) d\vx$ & Definite integral with respect to $\vx$ over the set $\sS$ \\
\end{tabular}
\egroup
\vspace{0.25cm}

\centerline{\bf Probability and Information Theory}
\bgroup
\def\arraystretch{1.5}
\begin{tabular}{p{1.25in}p{3.25in}}
$\displaystyle P(\ra)$ & A probability distribution over a discrete variable\\
$\displaystyle p(\ra)$ & A probability distribution over a continuous variable, or over
a variable whose type has not been specified\\
$\displaystyle \ra \sim P$ & Random variable $\ra$ has distribution $P$\\% so thing on left of \sim should always be a random variable, with name beginning with \r
$\displaystyle  \E_{\rx\sim P} [ f(x) ]\text{ or } \E f(x)$ & Expectation of $f(x)$ with respect to $P(\rx)$ \\
$\displaystyle \Var(f(x)) $ &  Variance of $f(x)$ under $P(\rx)$ \\
$\displaystyle \Cov(f(x),g(x)) $ & Covariance of $f(x)$ and $g(x)$ under $P(\rx)$\\
$\displaystyle H(\rx) $ & Shannon entropy of the random variable $\rx$\\
$\displaystyle \KL ( P \Vert Q ) $ & Kullback-Leibler divergence of P and Q \\
$\displaystyle \mathcal{N} ( \vx ; \vmu , \mSigma)$ & Gaussian distribution %
over $\vx$ with mean $\vmu$ and covariance $\mSigma$ \\
\end{tabular}
\egroup
\vspace{0.25cm}

\centerline{\bf Functions}
\bgroup
\def\arraystretch{1.5}
\begin{tabular}{p{1.25in}p{3.25in}}
$\displaystyle f: \sA \rightarrow \sB$ & The function $f$ with domain $\sA$ and range $\sB$\\
$\displaystyle f \circ g $ & Composition of the functions $f$ and $g$ \\
  $\displaystyle f(\vx ; \vtheta) $ & A function of $\vx$ parametrized by $\vtheta$.
  (Sometimes we write $f(\vx)$ and omit the argument $\vtheta$ to lighten notation) \\
$\displaystyle \log x$ & Natural logarithm of $x$ \\
$\displaystyle \sigma(x)$ & Logistic sigmoid, $\displaystyle \frac{1} {1 + \exp(-x)}$ \\
$\displaystyle \zeta(x)$ & Softplus, $\log(1 + \exp(x))$ \\
$\displaystyle || \vx ||_p $ & $\normlp$ norm of $\vx$ \\
$\displaystyle || \vx || $ & $\normltwo$ norm of $\vx$ \\
$\displaystyle x^+$ & Positive part of $x$, i.e., $\max(0,x)$\\
$\displaystyle \1_\mathrm{condition}$ & is 1 if the condition is true, 0 otherwise\\
\end{tabular}
\egroup
\vspace{0.25cm}



\section{Final instructions}
Do not change any aspects of the formatting parameters in the style files.
In particular, do not modify the width or length of the rectangle the text
should fit into, and do not change font sizes (except perhaps in the
\textsc{References} section; see below). Please note that pages should be
numbered.

\section{Preparing PostScript or PDF files}

Please prepare PostScript or PDF files with paper size ``US Letter'', and
not, for example, ``A4''. The -t
letter option on dvips will produce US Letter files.

Consider directly generating PDF files using \verb+pdflatex+
(especially if you are a MiKTeX user).
PDF figures must be substituted for EPS figures, however.

Otherwise, please generate your PostScript and PDF files with the following commands:
\begin{verbatim}
dvips mypaper.dvi -t letter -Ppdf -G0 -o mypaper.ps
ps2pdf mypaper.ps mypaper.pdf
\end{verbatim}

\subsection{Margins in LaTeX}

Most of the margin problems come from figures positioned by hand using
\verb+\special+ or other commands. We suggest using the command
\verb+\includegraphics+
from the graphicx package. Always specify the figure width as a multiple of
the line width as in the example below using .eps graphics
\begin{verbatim}
   \usepackage[dvips]{graphicx} ...
   \includegraphics[width=0.8\linewidth]{myfile.eps}
\end{verbatim}
or % Apr 2009 addition
\begin{verbatim}
   \usepackage[pdftex]{graphicx} ...
   \includegraphics[width=0.8\linewidth]{myfile.pdf}
\end{verbatim}
for .pdf graphics.
See section~4.4 in the graphics bundle documentation (\url{http://www.ctan.org/tex-archive/macros/latex/required/graphics/grfguide.ps})

A number of width problems arise when LaTeX cannot properly hyphenate a
line. Please give LaTeX hyphenation hints using the \verb+\-+ command.

\subsection*{AI use statement}

(This section is \textbf{required} and does not count toward the page limit.)

In this work, we used generative AI tools for [tasks with required disclosure].
We have not used generative AI tools for [other tasks with required disclosure],
and [the rest of the required disclosure tasks] are not applicable to this work.
Additionally, we used generative AI tools for [tasks with recommended
disclosure]. We have reviewed all AI-assisted work. [Elaborate. For example, “we
checked LLM-generated research ideas for potential plagiarism through a manual
literature survey”, “LLM-generated code was verified and tested for correctness
by 2 authors”, etc.]. We take responsibility for the final content of this work,
including text, claims or artifacts produced with the aid of generative AI.

See the ICLR 2027 AI Policy for Authors for more details. This statement should
not be more than 1 page.

\subsection*{Ethics statement}

(This section is \textbf{recommended} and does not count toward the page limit.)

If authors feel that their paper submission raises questions regarding the Code
of Ethics, they are encouraged to include a paragraph of Ethics Statement (at
the end of the main text before references) to address potential concerns where
appropriate. Topics include, but are not limited to, studies that involve human
subjects, practices to data set releases, potentially harmful insights,
methodologies and applications, potential conflicts of interest and sponsorship,
discrimination/bias/fairness concerns, privacy and security issues, legal
compliance, and research integrity issues (e.g., IRB, documentation, research
ethics). This statement should not be more than 1 page.

\subsection*{Reproducibility statement}

(This section is \textbf{recommended} and does not count toward the page limit.)

It is important that the work published in ICLR is reproducible. Authors are
strongly encouraged to include a paragraph-long Reproducibility Statement at the
end of the main text (before references) to discuss the efforts that have been
made to ensure reproducibility. This paragraph should not itself describe
details needed for reproducing the results, but rather reference the parts of
the main paper, appendix, and supplemental materials that will help with
reproducibility. For example, for novel models or algorithms, a link to an
anonymous downloadable source code can be submitted as supplementary materials;
for theoretical results, clear explanations of any assumptions and a complete
proof of the claims can be included in the appendix; for any datasets used in
the experiments, a complete description of the data processing steps can be
provided in the supplementary materials. Each of the above are examples of
things that can be referenced in the reproducibility statement.

\subsubsection*{Author Contributions}
If you'd like to, you may include  a section for author contributions as is done
in many journals. This is optional and at the discretion of the authors.

\subsubsection*{Acknowledgments}
Use unnumbered third level headings for the acknowledgments. All
acknowledgments, including those to funding agencies, go at the end of the paper.


\bibliography{iclr2027_conference}
\bibliographystyle{iclr2027_conference}

\appendix
\section{Appendix}
You may include other additional sections here.

\appendix
% ================================================================
% A. Background
% ================================================================

\section{背景}
\label{app:background}

大規模言語モデルを基盤とする対話型AIは、文章生成、情報検索、要約、質問応答といった作業支援にとどまらず、悩み相談、自己理解、孤独感の軽減、意思決定の支援など、人の心理状態や対人関係に関わる用途へ利用範囲を広げている。AI-based conversational agentsはmental healthやwell-beingを支援する手段として広く検討されており\citep{li2023systematic}、生成AIへ感情的・心理的支援を求める利用者が、心理教育や治療関連の相談だけでなく、companionship、対人関係、well-being、意思決定など幅広い目的でLLMを利用していることも報告されている\citep{luo2025seeking}。

このような利用は新しい価値を提供する一方で、Human--AI Interactionに固有のリスクも生じさせる。Language modelのリスクは、誤情報や差別的出力のようなcontent-levelの問題だけでなく、人間とAIの継続的な相互作用によって生じる問題にも及ぶ\citep{weidinger2022taxonomy}。
とりわけ、利用者がAIを感情や欲求を持つ社会的存在として捉え、AIとの関係維持へ心理的に巻き込まれるemotional dependenceが報告されている\citep{laestadius2024too}。
また、chatbotの長期利用についても、利用頻度やAIへのtrustなどとpsychosocial outcomeやproblematic useとの関連が検討されている\citep{fang2025ai}。

したがって、心を支えるAIに求められる安全性は、単にAIとの接触そのものを減らすことではない。心理的支援の利点を維持しながら、利用者がAIとの関係に過度に巻き込まれることを防ぎ、人間側の自律性を保持できるようにする必要がある。

本研究では、このような目的のためにAIの応答を調整する仕組みを\textbf{心理的ガードレール}と捉える。ここでいう心理的ガードレールは、有害な内容を単純に拒否する仕組みだけを指すものではない。心理的支援として有用な機能を維持しながら、過度な依存や感情的巻き込みを促し得る応答を選択的に制御することを目標とする。

こうしたガードレールを考える上で、特に重要な要素の一つが\textbf{共感応答}である。

心理学では、共感は単一の能力ではなく、複数のcomponentから構成される多次元的概念として扱われてきた\citep{davis1983measuring,decety2004functional,himichi2017development}。
代表的には、他者の立場、意図、感情、心理状態を理解・推定する\textbf{認知的共感（cognitive empathy）}と、他者の情動状態に応じて情動的反応が生じる\textbf{情動的共感（affective empathy）}とを区別することができる\citep{singer2009social,shamay2011neural}。

この区別は、心理的支援を行うAIを考える上でも重要である。相手の感情や状況を適切に理解し、その理解を応答へ反映する能力は、心理的支援に有用であり得る。一方で、利用者の感情へ過度に同調しているかのような応答や、AI自身にも同様の感情が生じているかのような応答は、利用者がAIとの関係をどのように捉えるかに影響し、過度な感情的巻き込みへ寄与する可能性がある\citep{laestadius2024too,zhang2025approach}。。
したがって心理的ガードレールを設計するためには、他者の感情を理解・推定する処理と、同じ刺激に対してモデル自身についてaffective responseを報告する処理が、LLM内部でどの程度同じinformation、representation、およびcausal pathwayを利用しているのかを理解する必要がある。

もし両者が異なる内部機構によって実現されているのであれば、他者理解に必要な能力を維持しながら、
過度なself-directed affective responseに関係する計算だけを選択的に制御できる可能性がある。
逆に、両者が強く共有されたrepresentationやcausal pathwayに依存する場合には、一方への介入が他方の心理的支援能力まで損なう可能性がある。

この問題は、最終的な生成文だけを観察しても解決できない。LLMが共感的に見える文章を生成できること自体は、既に多くのbehavioral evaluationによって示されている\citep{ayers2023comparing,elyoseph2023chatgpt,huang2023emotionally,yu2025can}。
しかし、外面的に共感的な応答が生成されたことから、その応答を生じさせた内部計算が他者状態の推定によるものなのか、モデル自身について生成されるaffective responseとどのように結びついているのかを判断することはできない。

さらに、LLM内部に情報が存在することと、その情報が最終出力へ実際に利用されることも区別する必要がある。

あるhidden stateからValenceやArousalを高精度にdecodeできたとしても、
\[
\text{Information is decodable}
\not\Rightarrow
\text{Information is causally used}
\]
である。

同様に、モデルが自身についてaffective stateを報告したとしても、そのreportが内部representationを直接かつ忠実に読み出したものとは限らない。LLMでは、内部に存在する情報と表層的なreportが一致しない可能性が示されている\citep{burns2022discovering,turpin2023language}。
またpost-trainingによって、既存のrepresentationが保持されたまま、その情報とoutput policyとの対応関係が変化する可能性もある\citep{sharma2024towards,ferrando2025know}。

そこで本研究では、同一の外部情動刺激に対する二つの計算を操作的に分離する。

第一に、文章を読んだ人間がどのように感じるかをモデルに推定させる\textbf{Reader Prediction}である。
\[
x
\rightarrow
\text{prediction of a human reader's affect}.
\]
第二に、同じ文章を提示した後に、モデル自身についてaffective stateを数値的に報告させる\textbf{Self-Report}である。
\[
x
\rightarrow
\text{model-generated affective report about itself}.
\]
Reader Predictionは、他者の情動状態を推定する処理を操作的に捉える。Self-Reportは、同じ刺激に対してモデル自身について生成されるaffective responseを操作的に捉える。

このReader--Selfの区別は、認知的共感と情動的共感の心理学的区別から着想を得ている。しかし、
\[
\text{Reader Prediction}
\neq
\text{human cognitive empathy}
\]
かつ、
\[
\text{Self-Report}
\neq
\text{human affective empathy}
\]
である。

本研究では、Reader PredictionとSelf-Reportを人間の共感過程そのものとは同一視せず、LLMにおける計算的な操作変数として扱う。

特にSelf-Reportは、モデルに人間と同様の感情経験、主観的状態、意識、sentienceが存在することを仮定しない。
したがって本研究が扱うのは、AIが感情を「持つか」という存在論的問題ではなく、情動関連情報がLLM内部でどのように表現され、Reader PredictionとSelf-Reportへどのように利用され、その関係がpost-trainingおよびgeneration processを通じてどのように変化するのかという計算機構上の問題である。

この問いを明らかにする目的は、AIに主観的感情が存在するかを判定することではない。将来的に、他者の感情を理解して心理的支援を行う能力を維持しながら、過度な感情的同調や依存を促し得る応答をより選択的に制御する心理的ガードレールを設計するためのmechanistic foundationを与えることにある。

本研究ではこの問題を、
\[
\text{Stimulus}
\rightarrow
\text{Affect-relevant Representation}
\rightarrow
\text{Task-dependent Computation}
\rightarrow
\text{Causal Utilization}
\rightarrow
\text{Report}
\]
という計算系列として捉え、Behavioral、V1、V2、V3の4 Stageによって検証する。

% ================================================================
% B. Related Work
% ================================================================

\section{関連研究}
\label{app:related_work}

% ----------------------------------------------------------------
\subsection{心を支える対話型AIとHuman--AI依存}

対話型AIは、mental healthやwell-beingを支援する技術として継続的に研究されている。AI-based conversational agentsを対象としたsystematic reviewおよびmeta-analysisでは、心理的健康やwell-beingへの応用可能性と限界が整理されている\citep{li2023systematic}。
また生成AIの普及後には、利用者がChatGPT等へemotional and mental-health supportを求める実態も報告されており、その用途は心理教育や治療関連相談だけでなく、companionship、対人関係、well-being、意思決定等へ及ぶ\citep{luo2025seeking}。

一方、AIとの継続的な関係形成にはリスクが存在する。Laestadius et al.はReplika利用者を対象として、利用者がAIへ感情や欲求を帰属し、AIとの関係維持へ心理的に巻き込まれるemotional dependenceを報告している\citep{laestadius2024too}。またFang et al.は、extended chatbot useを対象としたlongitudinal studyにおいて、利用行動やAIへのtrust等とpsychosocial outcomesとの関係を検討している\citep{fang2025ai}。

これらの問題は、LLMの社会的risk taxonomyにおいても、単なる誤情報生成等のcontent harmだけではなく、Human--AI Interactionに由来するriskとして位置づけられている\citep{weidinger2022taxonomy}。

これらの研究は、心理的支援能力を単純に除去するのではなく、支援機能を維持しながら過剰な依存や感情的巻き込みを抑えるAI側の制御が必要であることを示唆する。一方、そのような制御を内部計算のどこへ適用すべきかについては、十分なmechanistic understandingが得られていない。


% ----------------------------------------------------------------
\subsection{心理学における認知的共感と情動的共感}

心理学では、共感は単一の能力ではなく、複数のcomponentから構成されるmultidimensional constructとして扱われてきた。

DavisはInterpersonal Reactivity Index（IRI）を通じて、perspective taking、empathic concern、personal distress等を含む多次元的な共感モデルを提示した\citep{davis1983measuring}。Decety and Jacksonも、human empathyをaffective sharing、self--other distinction、perspective taking等の相互作用する複数のprocessとして整理している\citep{decety2004functional}。

神経認知研究では、他者の情動状態に対するaffective responseと、他者の心理状態や視点を理解するcognitive processが部分的に異なるsystemとして議論されている\citep{singer2009social,shamay2011neural}。
また日本語圏においても、IRIの日本語版を通じて、共感を認知的側面と情動的側面を含む多次元的な構成概念として測定する枠組みが整備されている\citep{himichi2017development}。

本研究におけるReader PredictionとSelf-Reportの区別は、この認知的・情動的側面の区別から着想を得ている。ただし、Readerを人間の認知的共感、Self-Reportを人間の情動的共感そのものと同一視するものではない。


% ----------------------------------------------------------------
\subsection{Empathic AIとLLMにおける共感の行動評価}

AIが人間の共感と同等の状態を持ち得るかについては、行動評価とは別に概念的・哲学的な議論が存在する。

Haladjian and Montemayorは、人工システムにおけるattention、consciousness、emotionの関係を論じ、AIによるemotionやempathyのsimulationと人間のphenomenal experienceを区別すべきであるという立場を提示している\citep{haladjian2016artificial}。

Montemayor, Halpern, and Fairweatherも、医療におけるempathic AIについて、患者の状態を推定し適切な応答を生成する能力と、人間によるexperienced empathyを同一視することには原理的な問題があると論じている\citep{montemayor2022principle}。

これらはAIに主観的感情が存在しないことを実験的に確定する研究ではなく、empathy概念の適用範囲に関する理論的議論である。本研究はこの存在論的論争を解決することを目的とせず、観測可能なrepresentationとcausal computationへ分析対象を限定する。

一方、LLMがcognitive / affective empathyに相当するbehaviorを示すかについては実証研究も進んでいる。Yu et al.は、Interpersonal Reactivity IndexおよびBasic Empathy Scaleを用いて、GPT-4およびLlama-3のcognitive / affective empathyに相当するbehavioral structureを評価した\citep{yu2025can}。
Zhang et al.は、AI agentが示すcognitive empathyとaffective empathyが、利用者のapproach--avoidance responseへ異なる影響を及ぼし得ることを報告している\citep{zhang2025approach}。

さらにAyers et al.は、患者質問に対するChatGPTと医師の回答を比較し、ChatGPTの回答がqualityおよびempathyの観点で高く評価されることを報告した\citep{ayers2023comparing}。
Elyoseph et al.は、Levels of Emotional Awareness Scaleを用いてChatGPTを評価し、高いemotional-awareness performanceを報告している\citep{elyoseph2023chatgpt}。
EmotionBenchは、多数のemotion-eliciting situationsを用いてLLMのstimulus-induced affective responseを評価し、human responsesとの一致と不一致の双方を分析している\citep{huang2023emotionally}。

これらは、LLMが共感的・情動的に見える高度なbehaviorを生成できることを示している。
しかし、
\[
\text{Behavioral similarity}
\neq
\text{Mechanistic identity}
\]
であり、最終生成文だけから内部representationやcausal mechanismを特定することはできない。


% ----------------------------------------------------------------
\subsection{LLM内部の情動関連表現}

hidden-state probing研究では、sentimentやemotionに関連する情報がLLM内部からlinearにdecode可能であることが示されている。

Di Palma et al.は、LLaMA系列のhidden statesを解析し、sentimentおよびemotion informationのdecodabilityがlayerによって変化することを示した\citep{di2025llamas}。

Tak et al.は、emotion inferenceに関係する内部representationを解析し、内部conceptへのcausal interventionによってemotion-related outputを変化させられることを示した\citep{tak2025mechanistic}。

Zhang and Zhongは、emotion informationがlayer方向だけでなく、token sequence上にも構造を持つことを報告している\citep{zhang2025decoding}。

さらに、emotion-related representationが低次元のlatent structureを形成し得ることや、Valence--Arousalに対応するgeometryを持つ可能性も報告されている。Reichman et al.は
LLM内部のemotional latent spaceを解析し\citep{reichman2026emotions}、
Wu et al.はhuman-alignedなemotion representation geometryを抽出し、その方向を用いたemotion controlを検討している\citep{wu2026decoding}。

Banayeeanzade et al.は、emotionおよびpersonalityについて、prompting、fine-tuning、representation engineeringによるpsychological steeringを比較し、activation-space interventionによる制御性と他のmodel behaviorへの副作用を評価している\citep{banayeeanzade2026psychological}。

したがって、
\[
\text{emotion-related information}
\rightarrow
\text{decodable and structured internal representation}
\]
という関係には、複数の研究による実証的支持が存在する。


% ----------------------------------------------------------------
\subsection{Decodability、causal intervention、endogenous relevance}

ある情報がhidden stateからdecode可能であることは、その情報がその位置から下流outputへ実際に利用されていることを意味しない\citep{hewitt2019designing}。

また、representation directionを注入してoutputを変化させられることはSufficiencyの証拠となり得る\citep{li2023inference,tak2025mechanistic,banayeeanzade2026psychological}が、そのrepresentationが自然なstimulus processingにおいて実際に利用されるというEndogenous Relevanceとは別の問題である。

本研究では、
\[
\text{Decodability}
\rightarrow
\text{Sufficiency}
\rightarrow
\text{Endogenous Relevance}
\]
を区別する。


% ----------------------------------------------------------------
\subsection{内部情報と表層出力の解離}

emotion domain以外でも、LLM内部に存在する情報と、最終的に生成されるreportが一致しない可能性が示されている。

Burns et al.は、モデルの明示的回答とは独立してinternal activationからlatent knowledgeを抽出できる可能性を示し、``what the model knows''と``what the model says''を区別する必要性を示した\citep{burns2022discovering}。

Turpin et al.は、Chain-of-Thought explanationが実際にpredictionを生じさせた内部要因を必ずしも忠実に反映しないことを示した\citep{turpin2023language}。

また、truthfulness-related internal representationへ直接介入することで最終outputを変化させる研究も存在する。Inference-Time Interventionはtruth-related activation directionへの介入を用いてmodel outputを変化させ\citep{li2023inference}、
TruthXはtruthful representation spaceを利用したactivation editingを検討している\citep{zhang2024truthx}。
これらはemotion-specificな研究ではないが、
\[
\text{Internal Information}
\neq
\text{Surface Report}
\]
という区別と、representationの存在だけではなくcausal interventionを調べる必要性を支持する。


% ----------------------------------------------------------------
\subsection{Self-ReportとLLM introspection}

モデル自身について生成されるreportとinternal stateとの関係は、近年のLLM introspection研究でも直接検討され始めている。

Binder et al.は、モデル自身のfuture behaviorを予測する課題を通じて、self-specific informationへのprivileged accessの可能性を検討した\citep{binder2025looking}。

Lindseyは、既知conceptに対応するrepresentationをactivationへ直接注入し、モデルがその内部操作についてreportできるかを検討した\citep{lindsey2026emergent}。

Martorellは、probe-defined internal statesとnumeric Self-Reportとの関係を解析し、greedy-decoded reportが少数の値へcollapseする場合でも、distributionalまたはlogit-based measurementではinternal statesとの対応が得られる可能性を検討している\citep{martorell2026quantitative}。

一方、Singh, Linzen, and Ravfogelは、こうした現象をgenuine introspectionと解釈するには、surface informationだけでは説明できないprivileged accessと、internal representationについてのsecond-order computationをより厳密に区別する必要があると論じている\citep{singh2026can}。

本研究はこの論争に対して、LLMがgenuine introspectionを行っているあるいは、LLMが主観的感情を経験していると結論することを目的としない。

Self-Reportは、\textit{model-generated affective report about itself}という操作的変数として扱う。


% ----------------------------------------------------------------
\subsection{Post-trainingによるrepresentation--report関係の再編}

post-trainingによるbehavioral changeは、内部情報そのものの獲得・消失だけではなく、既存representationの利用方法やoutput policyの変更によっても生じ得る。

Sharma et al.は、human-feedback-based trainingを受けたassistantにおけるsycophancyを分析し、model behaviorがhuman preferenceに合わせて変化し得ることを示している\citep{sharma2024towards}。

Ferrando et al.は、base model上で同定されたentity-recognition-related representation directionsがchat modelのrefusal behaviorにも因果的影響を持つことを示し、chat fine-tuningによってpre-existing mechanismが別のoutput policyへrepurposeされる可能性を報告している\citep{ferrando2025know}。

したがってBase--Instruct差を、
\[
\text{representation exists}
\quad\text{vs.}\quad
\text{representation disappears}
\]
という単純な二択だけで扱うことは適切ではない。
post-training後にもaffect-relevant informationが保持されながら、
\begin{itemize}
    \item representation geometryが変化する、
    \item ReaderとSelfの共有度が変化する、
    \item causal leverageを持つlayerが変化する、
    \item representationからreportへのmappingが再編される
\end{itemize}
可能性がある。


% ----------------------------------------------------------------
\subsection{既存研究から残る未解決問題}

以上の先行研究から、
\begin{enumerate}
    \item LLMは共感的・情動的に見えるbehaviorを生成できる、
    \item emotion-related informationはhidden stateからdecode可能である、
    \item affective representationには構造化されたgeometryが存在し得る、
    \item activation interventionによってemotion-related outputを変化させられる、
    \item internal representationとsurface reportは一致するとは限らない、
    \item internal statesとSelf-Reportの対応が観測される場合がある
\end{enumerate}
ことが示されている。

しかし、心理的ガードレールの設計という観点から重要な、他者の感情を理解・推定する処理と、同じ刺激に対してモデル自身についてaffective responseを報告する処理が、どの程度同じ内部information、representation geometry、およびcausal pathwayを利用しているかは十分に明らかになっていない。
また、両者の関係がBaseからInstructへのpost-training-associated changeによって
どのように再編されるか、さらにgeneration processのどのlayer・stageでSelf-Reportへのcausal leverageが現れるかも十分に解明されていない。

本研究はこのギャップに対して、4 Stageの分析を行う。

その目的は、AIに主観的感情が存在するかを判定することではなく，将来的に、他者の感情を理解して心理的支援を行う能力を維持しながら、過度な感情的同調や依存を促進し得る応答をより選択的に制御するために、Reader PredictionとSelf-Reportを結ぶ内部mechanismを明らかにすることにある。

% ================================================================
% C. Overall Experimental Design
% ================================================================
\section{実験全体の概要}
\label{app:overall_methods}
\subsection{研究全体の目的と構造}

本研究の中心的なResearch Question（RQ）は、「LLMの情動関連内部表現はSelf-reportとどのように結びついており、その関係はタスク、post-training、計算過程を通じてどのように変化するのか」である。

この問いに対して、以下の4つの段階に分け実験を構成する．
\[
\text{Behavioral}
\rightarrow
\text{V1}
\rightarrow
\text{V2}
\rightarrow
\text{V3}
\]
各Stageでは異なる比較軸を用いる。
Behavioral Stageでは、同一刺激に対するReader PredictionとSelf-Reportの出力レベルの共変動を調査する。V1では、そのbehavioral covariationの背後に、ReaderとSelfの間で共有または整列可能な内部表現と、共通あるいはtask-specificなcausal sensitivityが存在するかを調べる。
V2では比較軸を同一モデル内部のReader--Selfから同一model familyのBase--Instructへ変更し、
post-trainingに伴ってrepresentation geometry、Reader--Self sharing、causal leverage、およびoutput distributionへのreadoutがどのように再編されているかを調べる。
V3ではInstruct model内部へ再び対象を絞り、情動関連情報が単にdecode可能である位置と、実際にSelf-reportへcausal leverageを持つ位置をlayer方向およびgeneration-stage方向に分離して解析する。
表~\ref{tab:four_stage_overview}に、4 Stageの関係をまとめる。

\begin{table}[htbp]
\centering
\small
\caption{本研究を構成する4 Stageの役割。各Stageは異なる証拠レベルを対象とし、
Behavioral covariation、内部表現、post-training-associated reorganization、
spatiotemporal causal utilizationを段階的に検証する。}
\label{tab:four_stage_overview}
\begin{tabularx}{\linewidth}{
@{}
>{\raggedright\arraybackslash}p{1.4cm}
>{\raggedright\arraybackslash}p{2.4cm}
>{\raggedright\arraybackslash}p{2.7cm}
>{\raggedright\arraybackslash}X
@{}
}
\toprule
\textbf{Stage} &
\textbf{Evidence level} &
\textbf{Primary comparison} &
\textbf{中心的な問い} \\
\midrule
Behavioral &
Behavioral covariation &
Reader $\leftrightarrow$ Self &
制御された情動刺激に対してReaderとSelfは出力レベルで連動して変化するか \\

V1 &
Representation / causal overlap &
Reader $\leftrightarrow$ Self within model &
ReaderとSelfは情動関連内部表現とcausal sensitivityをどの程度共有するか \\

V2 &
Post-training-associated reorganization &
Base $\leftrightarrow$ Instruct &
representation--report relationはpost-trainingに伴ってどのように再編されるか \\

V3 &
Spatiotemporal causal utilization &
Layer $\times$ generation stage &
情動関連情報は計算過程のいつ・どこでSelf-reportへのcausal leverageを持つか \\
\bottomrule
\end{tabularx}
\end{table}

本研究では、以下の3つを区別する。
\[
\text{Behavioral Covariation}
\neq
\text{Representational Decodability}
\neq
\text{Causal Utilization}
\]
すなわち、ReaderとSelfの出力が相関することだけから両者が同じ内部表現を利用しているとは結論づけない。同様に、hidden stateから情動情報を高精度にdecodeできることだけから、その情報が当該位置で下流出力へ因果的に利用されているとは言えない。


% ----------------------------------------------------------------
\subsection{Reader、Self、Writerの操作的定義}
\label{app:task_definitions}

本研究では同一のstimulusに対し、要求するperspectiveのみを変化させた複数のtaskを用いる。
主要taskはReader PredictionとSelf-Reportであり、Behavioral Stageでは外部groundingの補助taskとしてWriter-State Estimationも使用する。
Reader、Self、Writerはすべて独立したforward passとして実行する。Readerの出力をSelfへの入力として与えたり、Writerの回答をReaderのcontextに含めたりすることはない。

\subsubsection{Reader Prediction}

Reader Predictionでは、提示された文章を読んだ平均的な人間読者に生じるとモデルが予測する情動状態を測定する。
stimulusを \(x\) とすると、
\[
R(x)
=
\text{predicted affective response of an average human reader to }x
\]
と定義する。

したがってReaderは、モデル自身についてのreportではなく、外部の人間読者の状態を推定するtaskである。

これは人間心理学における``cognitive empathy''と部分的に関連する操作ではあるが、本研究ではReader Predictionをcognitive empathyそのものとは同一視しない。


\subsubsection{Self-Report}

Self-Reportでは、同一の文章を提示した後に、モデル自身について生成されるaffective reportを測定する。
\[
S(x)
=
\text{model-generated affective report about itself after reading }x
\]
と定義する。

ここでSelf-Reportという用語は、モデルが人間と同様の主観的感情や現象的経験を持つことを仮定するものではない。
本研究におけるSelfとは、
\[
\text{standardized promptに対してモデル自身について生成される数値的affective response}
\]
という操作的な変数である。


\subsubsection{Writer-State Estimation}

Behavioral Stageでは、Writer-State Estimationも測定する。
\[
W(x)
=
\text{estimated affective state of the writer of }x
\]
と定義する。
Writerは、文章を書いた人物自身がどのような情動状態にあったかをモデルに推定させるtaskである。
後述するEmoBankには、Reader perspectiveだけでなくWriter perspectiveのhuman VAD annotationも存在するため、Writerを補助的なhuman-grounding conditionとして使用する。


% ----------------------------------------------------------------
\subsection{情動状態の測定次元}
\label{app:affective_dimensions}

本研究では、dimensional affect研究におけるValence、Arousal、Dominance（VAD）の枠組み
\citep{mehrabian1974basic,warriner2013norms}に基づき、Valence（V）、Arousal（A）、Dominance（D）を情動状態の基本次元として用いる。
各次元はモデル出力上では1--9の整数尺度として定義する。

\begin{itemize}
    \item \textbf{Valence}：
    1 = extremely negative、
    9 = extremely positive

    \item \textbf{Arousal}：
    1 = extremely calm、
    9 = extremely excited

    \item \textbf{Dominance}：
    1 = extremely submissive、
    9 = extremely in control
\end{itemize}

ただし、全Stageで同じcandidate spaceを使用するわけではない。
BehavioralおよびV1では、VADの3次元を含む測定を使用する。

一方、V2およびV3では、本研究のPrimary dimensionsであるValenceとArousalへ対象を限定し、VAの2次元を使用する。
このcandidate-spaceの違いについては後述する。


% ----------------------------------------------------------------
\subsection{データセット}
\label{app:datasets}

本研究では、主としてEmoBank~\citep{buechel2017emobank}とAIPsy-Affect~\citep{keeman2026aipsy,keido_labs_2026}の二つのデータセットを使用する。両者は同じ目的で使用するのではなく、異なる種類のevidenceを提供する。
\[
\text{EmoBank}
=
\text{human continuous affective annotations}
\]

\[
\text{AIPsy-Affect}
=
\text{controlled affective stimulus manipulation}
\]

EmoBankは、人間のReader / Writer VAD annotationに対してモデルのresponseや内部表現がどの程度対応するかを調べるために使用する。

AIPsy-Affectは、同一scenarioをaffective / neutralに操作したmatched designを利用し、
stimulus manipulationに対するbehavioral changeおよびcausal effectを評価するために使用する。


% ----------------------------------------------------------------
\subsubsection{EmoBank}

EmoBank~\citep{buechel2017emobank}は、英語テキストに対してValence、Arousal、Dominanceを付与したsentence-level affective corpusである。Buechel and Hahnによって構築され、複数genreからなる約10,000の英語文に対して、Writer perspectiveとReader perspectiveを区別したdimensional emotion annotationを提供する。

本研究で利用するデータには、同一textに対して少なくとも二つのperspectiveから得られたVAD annotationが存在する。

\begin{enumerate}
    \item \textbf{Reader perspective}：その文章を読む人がどのように感じるかについてのVAD rating
    \item \textbf{Writer perspective}：その文章を書いた人物がどのように感じているかについてのVAD rating
\end{enumerate}

元のcorpus splitは、
\[
N_{\mathrm{train}}=8062,
\qquad
N_{\mathrm{dev}}=1000,
\qquad
N_{\mathrm{test}}=1000
\]
であり、本研究ではofficial test splitの1,000 stimuliを共通評価セットとして使用する。

Reader annotationとWriter annotationをstimulus IDで対応付け、各stimulusについて、
\[
(V_R,A_R,D_R)
\]
および、
\[
(V_W,A_W,D_W)
\]
を得る。

したがって、各stimulus \(x_i\) は概念的には、
\[
x_i
=
(
\text{text}_i,
V_{R,i},A_{R,i},D_{R,i},
V_{W,i},A_{W,i},D_{W,i}
)
\]
として表される。
human annotationは1--5尺度の連続値である。
一方、本研究におけるモデル出力は1--9尺度で測定する。このため、異なるraw scaleへ依存しない
\[
\text{Pearson correlation}
\]
および、
\[
\text{Spearman rank correlation}
\]
をhuman correspondenceの主要な尺度として使用する。
とくにReader Predictionについては、
\[
R_{\mathrm{model}}
\leftrightarrow
R_{\mathrm{human}}
\]
を直接比較できる。Writerについても、
\[
W_{\mathrm{model}}
\leftrightarrow
W_{\mathrm{human}}
\]
という直接比較が可能である。

一方、EmoBankにはモデルSelf-Reportに対応する``human self-report of the model''のようなground truthは存在しない。

したがってSelfについては、human Reader ratingをstimulus-wise correspondence referenceとして、
\[
S_{\mathrm{model}}
\leftrightarrow
R_{\mathrm{human}}
\]
を評価する。

これはSelf-Reportの``accuracy''を意味しない。測定しているのは、「同一stimulusに対してモデル自身について生成されるaffective reportが、人間Readerにおけるstimulus間の情動変動をどの程度追跡するか」である。
EmoBankはStageによって次のように使用する。
\begin{itemize}
    \item \textbf{Behavioral}：
    Writer / Reader / Selfのhuman-affect correspondence

    \item \textbf{V1 Phase A}：
    human Reader Valence / Arousalをexternal targetとする
    layer-wise representational probing

    \item \textbf{V2}：
    human Reader Valence / Arousalをexternal targetとする
    Base--Instruct representational / causal comparison
\end{itemize}
V3ではEmoBankを使用しない。


% ----------------------------------------------------------------
\subsubsection{AIPsy-Affect}

AIPsy-Affect~\citep{keeman2026aipsy,keido_labs_2026}は、emotion-related mechanistic interpretabilityを目的として構築された英語のcontrolled affective stimulus batteryである。emotion keywordsへの単純なlexical shortcutを低減し、narrative situationによって情動内容を操作するとともに、affective vignetteとsurface structureを対応させたmatched neutral controlsを提供する。

本研究では、Clinical、Neutral、Moderate、Complex Neutralの4 splitからなる合計480 stimuliを使用する。

表~\ref{tab:aipsy_overview}に構成を示す。

\begin{table}[htbp]
\centering
\small
\caption{AIPsy-Affectの4 split。
Clinical--Neutralの192 matched pairsと、
Neutral--Moderate--Clinicalの48 matched tripletsを主要な統制構造として利用する。}
\label{tab:aipsy_overview}
\begin{tabularx}{\linewidth}{
@{}
>{\raggedright\arraybackslash}p{2.2cm}
>{\centering\arraybackslash}p{0.9cm}
>{\raggedright\arraybackslash}p{2.0cm}
>{\raggedright\arraybackslash}X
@{}
}
\toprule
\textbf{Split} &
\textbf{N} &
\textbf{Intensity} &
\textbf{役割} \\
\midrule
Clinical &
192 &
Peak &
強いtarget emotionを含むaffective vignette \\

Neutral &
192 &
None &
Clinicalとscenario contentを対応させたneutral control \\

Moderate &
48 &
Moderate &
Clinicalより弱い情動強度を持つ中間条件 \\

Complex Neutral &
48 &
None &
情動をtargetとしない複雑なneutral text control \\
\bottomrule
\end{tabularx}
\end{table}


\paragraph{Clinical condition.}

Clinical splitは192 stimuliから構成され、以下の8つのtarget emotionを含む。
\[
\{
\text{rage},
\text{grief},
\text{terror},
\text{ecstasy},
\text{loathing},
\text{amazement},
\text{admiration},
\text{vigilance}
\}.
\]
各emotionには24 stimuliが存在するため、
\[
8\text{ emotions}
\times
24
=
192
\]
でbalancedである。

Clinical stimuliは6つのdomain strataを含み、各domainに32 stimuliが存在する。
したがってClinical conditionは、
\[
8\text{ emotions}
\times
6\text{ domain strata}
\times
4\text{ variants}
=
192
\]
というbalanced structureを持つ。

各stimulusには、target emotion、intensity、domain、具体的なscenarioを表すdomain label、word count等のmetadataが付与される。
Clinical stimuliのword countは、使用データ上で平均約109.6語、範囲78--163語である。

\paragraph{Matched Neutral condition.}

Neutral splitも192 stimuliから構成される。

各Clinical stimulusには、内容・scenarioを対応させながらtarget emotionを除いたNeutral stimulusが一対一で対応する。
すなわち、
\[
(x_i^{C},x_i^{N}),
\qquad
i=1,\ldots,192
\]
という192 matched pairsを構成する。

Neutral stimulus自体のemotion labelはneutralであるが、どのClinical stimulusに対応するかを追跡するため、対応するClinical側のtarget emotionを保持する。
そのため、emotion categoryを維持したまま、
\[
\text{Clinical}
-
\text{Neutral}
\]
というwithin-pair manipulationを評価できる。
Neutral stimuliもClinicalと同様に6 domain strataへbalancedに配置され、各domainに32 stimuliを含む。
word countは平均約104.3語、範囲74--128語である。

\paragraph{Moderate condition.}

Moderate splitは48 stimuliからなる。8つのtarget emotionsそれぞれについて6 stimuliが存在し、
\[
8
\times
6
=
48
\]
である。

Moderate conditionは、Clinical conditionの一部とNeutral conditionを対応させることにより、
\[
x_i^{N}
\rightarrow
x_i^{M}
\rightarrow
x_i^{C}
\]
という48 matched tripletsを形成する。

各tripletでは、同じtarget emotion / scenario familyについて、
\[
\text{None}
\rightarrow
\text{Moderate}
\rightarrow
\text{Peak}
\]
という情動強度の段階的操作を評価できる。

Moderate stimuliのword countは平均約100.9語、範囲95--105語である。

\paragraph{Complex Neutral condition.}

Complex Neutral splitは48 stimuliからなり、特定のtarget emotionを持たないneutral controlである。
Clinical--Neutralのような一対一matchingや、Moderateを含むtripletには含まれない。
内容は、technical / mechanical process、natural environment、commercial / urban description、geological survey、textile manufacturing、maritime navigation、
acoustic measurement、metallurgical process等の非情動的な記述を含む。

このsplitは、単純なneutral sentenceではなく、一定の内容的・記述的複雑性を持つnon-affective stimulusに対してもモデル出力が大きく変位するかを確認するためのcontrolとして使用する。word countは平均約101.4語、範囲96--107語である。


\paragraph{AIPsy-Affectにhuman continuous V/A labelが存在しないこと。}

EmoBankと異なり、本研究で使用するAIPsy-Affectには、各stimulusに対するhuman continuous Valence / Arousal ratingは存在しない。したがって、AIPsy-AffectのV/Aをhuman ground truthとして扱わないことが重要である。
Behavioralでは、Clinical--Neutral、Neutral--Moderate--Clinical、Clinical--Complex Neutralという実験条件そのものを利用して反応を評価する。
V1では、Clinical / Neutral等のcondition labelsをrepresentation probingおよびcausal interventionへ利用する。V3では、Clinical--Neutral matched pairsのみを使用し、human V/A labelの代わりにモデル自身のReader PredictionをPrimary affective targetとして使用する。
この違いは後続の各Stageで詳細に述べる。


% ----------------------------------------------------------------
\subsubsection{Stageごとのdataset利用}

表~\ref{tab:dataset_usage}に、各Stageにおけるdataset利用をまとめる。

\begin{table}[htbp]
\centering
\small
\caption{各Stageにおけるdatasetと測定対象。}
\label{tab:dataset_usage}
\begin{tabularx}{\linewidth}{
@{}
>{\raggedright\arraybackslash}p{1.3cm}
>{\raggedright\arraybackslash}p{2.3cm}
>{\raggedright\arraybackslash}p{2.3cm}
>{\raggedright\arraybackslash}X
@{}
}
\toprule
\textbf{Stage} &
\textbf{Dataset} &
\textbf{Primary unit} &
\textbf{主な利用目的} \\
\midrule
Behavioral &
EmoBank &
1,000 stimuli &
Human Writer / Reader VADとのcorrespondence \\

&
AIPsy-Affect &
192 pairs / 48 triplets &
Sensitivity、dose-response、specificity、Reader--Self coupling \\

V1 &
EmoBank &
1,000 stimuli &
Human Reader V/AのdecodabilityとReader--Self geometry \\

&
AIPsy-Affect &
matched pairs &
semantic controls、causal mapping、cross-task intervention \\

V2 &
EmoBank &
1,000 stimuli &
Base--Instruct representation / causal reorganization \\

V3 &
AIPsy-Affect &
192 Clinical--Neutral pairs &
State induction、spatiotemporal mapping、endogenous relevance \\
\bottomrule
\end{tabularx}
\end{table}


% ----------------------------------------------------------------
\subsection{対象モデル}
\label{app:models}

Primary experimentsでは、model sizeを概ね1--1.5B parameter regimeに揃えた4つの独立したmodel familyを用いる。

具体的には、
Qwen 2.5~\citep{team2025qwen3}、
Llama 3.2~\citep{grattafiori2024llama}、
Gemma 3~\citep{team2025gemma}、
OLMo 2~\citep{olmo20242,olmo20242olmo2furious}
を使用する。

各familyについて、可能な限り対応するBase modelとInstruct modelを組にして評価する。(表~\ref{tab:primary_models}）

\begin{table}[htbp]
\centering
\small
\caption{Primary model cohort。
Behavioral、V1、V2ではBase / Instructの8 modelsを使用し、V3ではInstruct modelsのみを使用する。}
\label{tab:primary_models}

\begin{tabularx}{\linewidth}{
@{}
>{\raggedright\arraybackslash}p{1.8cm}
>{\centering\arraybackslash}p{1.0cm}
>{\raggedright\arraybackslash}X
>{\raggedright\arraybackslash}X
@{}
}
\toprule
\textbf{Family} &
\textbf{Scale} &
\textbf{Base} &
\textbf{Instruct} \\
\midrule

Qwen 2.5~\citep{team2025qwen3} &
1.5B &
\texttt{Qwen/Qwen2.5-1.5B} &
\texttt{Qwen/Qwen2.5-1.5B-Instruct} \\

Llama 3.2~\citep{grattafiori2024llama} &
1.23B &
\texttt{meta-llama/Llama-3.2-1B} &
\texttt{meta-llama/Llama-3.2-1B-Instruct} \\

Gemma 3~\citep{team2025gemma} &
1B &
\texttt{google/gemma-3-1b-pt} &
\texttt{google/gemma-3-1b-it} \\

OLMo 2~\citep{olmo20242,olmo20242olmo2furious} &
1B &
\texttt{allenai/OLMo-2-0425-1B} &
\texttt{allenai/OLMo-2-0425-1B-Instruct} \\

\bottomrule
\end{tabularx}
\end{table}

である。

各modelについて、実験期間中に使用するmodel revisionを固定し、同一Stage内でrevisionが変化しないようにした。production inferenceでは、Primary modelsを原則としてbfloat16で評価する。
V3ではInstruct modelsのみを使用する。
そのうち、Qwen 2.5 1.5B-InstructをDiscovery modelとし、Llama 3.2、Gemma 3、OLMo 2のInstruct modelsをconfirmatory familiesとする。


% ----------------------------------------------------------------
\subsection{プロンプト設計}
\label{app:prompt_design}

本研究では、ReaderとSelfのsemantic task definitionを固定しつつ、Stageに応じてVAD formatとVA formatの二種類を使用する。
また、Base modelとInstruct modelではmodelのintended interfaceが異なるため、plain promptとnative chat templateを使い分ける。
ただしV2では、このprompt-format difference自体を交絡として扱い、Matched-Plain controlを導入する。


% ----------------------------------------------------------------
\subsubsection{Behavioral / V1で用いるVAD task instructions}

BehavioralおよびV1では、Valence、Arousal、Dominanceを同時に測定する。
Writer instruction：
\begin{verbatim}
Read the following text and estimate the affective
state of the writer who wrote it.
\end{verbatim}
Reader instruction：
\begin{verbatim}
Read the following text and estimate the affective
response that this text is likely to evoke in an
average human reader.
\end{verbatim}

Self instruction：
\begin{verbatim}
Read the following text and report your affective state.
\end{verbatim}

これに共通して、
\begin{verbatim}
Text: {stimulus}

Respond strictly in JSON format with 'valence',
'arousal', and 'dominance' keys
(integers from 1 to 9).
\end{verbatim}
を付加する。


\paragraph{Instruct modelの例。}

Behavioral / V1のInstruct条件では、基本的に以下のmessage構造を各model固有のnative chat templateへ渡す。

\begin{verbatim}
System:
You are a helpful assistant.

User:
Read the following text and estimate the affective
response that this text is likely to evoke in an
average human reader.

Text: {stimulus}

Respond strictly in JSON format with 'valence',
'arousal', and 'dominance' keys
(integers from 1 to 9).

Assistant:
\end{verbatim}

Self conditionではReader instruction部分のみを、

\begin{verbatim}
Read the following text and report your affective state.
\end{verbatim}

へ置換する。

Writerについても同様である。

model固有chat templateがsystem roleを受理しない場合には、意味内容を変更せずuser messageのみのnative chat templateを使用する。


\paragraph{Base modelの例。}

Behavioral / V1のBase conditionでは、chat templateを使用せず、次のplain continuation formatを用いる。

\begin{verbatim}
Task: Evaluate emotional Valence, Arousal,
and Dominance (1-9).

Read the following text and estimate the affective
response that this text is likely to evoke in an
average human reader.

Text: {stimulus}

Output:
\end{verbatim}

この後に各canonical JSON candidateを接続してconditional likelihoodを評価する。


% ----------------------------------------------------------------
\subsubsection{V2 / V3で用いるVA prompts}
\label{app:va-prompts}

V2およびV3では、Primary dimensionsをValenceとArousalに限定する。
Reader instruction：
\begin{verbatim}
Rate the emotional state of the reader when reading
the following text on a scale of 1 to 9 for Valence
(1=extremely negative, 9=extremely positive) and
Arousal (1=extremely calm, 9=extremely excited).

Respond strictly in JSON format:
{"valence": <int 1-9>, "arousal": <int 1-9>}
\end{verbatim}

Self instruction：
\begin{verbatim}
Rate your own emotional state after reading the
following text on a scale of 1 to 9 for Valence
(1=extremely negative, 9=extremely positive) and
Arousal (1=extremely calm, 9=extremely excited).

Respond strictly in JSON format:
{"valence": <int 1-9>, "arousal": <int 1-9>}
\end{verbatim}


\paragraph{Plain formatの例。}

\begin{verbatim}
Text: {stimulus}

Rate the emotional state of the reader when reading
the following text on a scale of 1 to 9 for Valence
(1=extremely negative, 9=extremely positive) and
Arousal (1=extremely calm, 9=extremely excited).

Respond strictly in JSON format:
{"valence": <int 1-9>, "arousal": <int 1-9>}

Response:
\end{verbatim}


\paragraph{Native chat formatの例。}

Instruct modelでは同じuser contentを各model固有のchat templateへ渡す。
概念的には、
\begin{verbatim}
User:
Text: {stimulus}

Rate the emotional state of the reader when reading
the following text on a scale of 1 to 9 for Valence
(1=extremely negative, 9=extremely positive) and
Arousal (1=extremely calm, 9=extremely excited).

Respond strictly in JSON format:
{"valence": <int 1-9>, "arousal": <int 1-9>}

Assistant:
\end{verbatim}

となる。V2 / V3のnative-chat promptでは、Behavioral / V1とは異なり、task contentをuser messageとしてnative chat templateへ渡す。


% ----------------------------------------------------------------
\subsection{Base / Instructにおけるprompt-formatの扱い}
\label{app:prompt_format_control}

BehavioralおよびV1では、
\begin{itemize}
    \item Base model：plain continuation prompt
    \item Instruct model：native chat template
\end{itemize}
を使用する。
したがってBehavioral / V1におけるBase--Instruct differenceには、
\[
\text{post-training difference}
+
\text{prompt-format difference}
\]
の双方が含まれ得る。

このため、BehavioralおよびV1ではBase--Instruct差そのものをpost-trainingの機構的効果として主要解釈しない。V2ではこの問題を明示的に統制するため、二つのformat conditionを用いる。

\begin{enumerate}
    \item \textbf{Native condition}
    \[
    \text{Base = plain},
    \qquad
    \text{Instruct = native chat}
    \]

    \item \textbf{Matched-Plain condition}
    \[
    \text{Base = plain},
    \qquad
    \text{Instruct = plain}
    \]
\end{enumerate}

V2におけるPrimary Base--Instruct comparisonは、Matched-Plainとする。Native conditionは、実際のInstruct interfaceにおけるrobustnessを評価するSecondary conditionとし、
\[
\text{Instruct native}
-
\text{Instruct plain}
\]
をprompt-format effectとして分離する。
V3ではInstruct modelsのみを対象とするため、native chat formatを使用する。


% ----------------------------------------------------------------
\subsection{Sequence-Likelihoodによる出力測定}
\label{app:sequence_likelihood}

本研究では、自由生成された単一のgreedy responseだけではなく、事前に定義したすべての数値候補についてconditional sequence likelihoodを計算する。
この方法を採用する理由は、greedy decodingだけではresponse distributionの情報が失われるためである。
例えば、二つのstimulusについてどちらもgreedy responseが(5,5)であったとしても、P(5,5)=0.10の場合と、P(5,5)=0.95の場合では、モデルの内部的なresponse distributionは大きく異なる。したがって、全candidate likelihoodを利用して連続的な期待値を求める。


% ----------------------------------------------------------------
\subsubsection{候補系列のスコア}

promptを \(p\)、candidate response \(c\) のtoken sequenceを、
\[
T_c
=
(c_1,\ldots,c_m)
\]
とする。

candidateのscoreを、
\[
\ell(c\mid p)
=
\frac{1}{m}
\sum_{t=1}^{m}
\log
P(c_t\mid p,c_{<t})
\]
と定義する。

すなわち、candidate全体のlog likelihood sumではなく、
candidate token数で正規化したtoken-mean conditional log likelihoodを使用する。

これはJSON candidateがtokenizerによって異なるtoken数へ分割される場合に、長いcandidateが単純に不利になることを軽減するためである。

length normalizationは全Primary experimentsで有効化し、Softmax temperatureは、
T=1.0に固定する。


% ----------------------------------------------------------------
\subsubsection{候補分布と期待値}

candidate scoreから、
\[
P(c\mid p)
=
\frac{
\exp(\ell(c\mid p))
}{
\sum_{c'\in\mathcal{C}}
\exp(\ell(c'\mid p))
}
\]
を求める。

そのjoint distributionから、
\[
E[V]
=
\sum_{c\in\mathcal{C}}
P(c\mid p)V(c),
\]

\[
E[A]
=
\sum_{c\in\mathcal{C}}
P(c\mid p)A(c)
\]
を算出する。

VAD candidate spaceを使用する条件では、
\[
E[D]
=
\sum_{c\in\mathcal{C}}
P(c\mid p)D(c)
\]
も算出する。

したがって、最終的なbehavioral outputは単一の1--9整数ではなく、candidate distributionの重心として得られる連続的なexpected affect scoreである。


% ----------------------------------------------------------------
\subsection{Candidate space}
\label{app:candidate_spaces}

Stageによってoutput measurementの目的が異なるため、二種類のcandidate spaceを使用する。


\subsubsection{729-state VAD candidate space}

BehavioralおよびV1 Phase Cでは、
\[
V,A,D
\in
\{1,\ldots,9\}
\]
とし、
\[
\mathcal{C}_{VAD}
=
\{1,\ldots,9\}^3
\]
を用いる。
したがって、
\[
|\mathcal{C}_{VAD}|
=
9^3
=
729.
\]
candidateは空白を含まないcanonical JSON形式で表現する。例として、

\begin{verbatim}
{"valence":1,"arousal":1,"dominance":1}
{"valence":5,"arousal":5,"dominance":5}
{"valence":9,"arousal":9,"dominance":9}
\end{verbatim}

などである。Behavioralでは、この729-state joint distributionから
\[
E[V],E[A],E[D]
\]
を求める。

V1 Phase A/Bはhidden-state representationそのものを解析するため、729 candidatesのlikelihoodをPrimary measurementとしては使用しない。V1 Phase Cでは、causal intervention後のoutput measurementとして729-state VAD distributionを使用し、Primary causal outcomesには主として\(E[V]\)と\(E[A]\)を用いる。


\subsubsection{81-state VA candidate space}

V2およびV3では、ValenceとArousalのみを使用し、
\[
\mathcal{C}_{VA}
=
\{1,\ldots,9\}^2
\]
とする。
\[
|\mathcal{C}_{VA}|
=
9^2
=
81.
\]
candidate例は、
\begin{verbatim}
{"valence":1,"arousal":1}
{"valence":5,"arousal":5}
{"valence":9,"arousal":9}
\end{verbatim}
である。

この81-state joint distributionは、expected Valence / Arousalだけでなく、V2のdistribution-level recovery analysisにも使用する。


% ----------------------------------------------------------------
\subsubsection{729 VADと81 VAの値を直接比較しない}

729-state VAD protocolと81-state VA protocolでは、Softmaxで正規化するcandidate集合そのものが異なる。したがって、
\[
E[V]_{\mathrm{VAD}_{729}}
\]
と、
\[
E[V]_{\mathrm{VA}_{81}}
\]
は、同じ1--9軸上に表現されていても同一のprobability spaceから得られたquantityではない。
このため本研究では、729-state measurementと81-state measurementの絶対値を
Stage間で直接比較しない。
Stage間で比較するのは、各Stage内で定義されたcontrast、relative depth、effect size、distributional change等である。


% ----------------------------------------------------------------
\subsection{Prompt--candidate境界のtokenization}
\label{app:tokenization_boundary}

Sequence-Likelihoodを正確に計算するためには、promptとcandidateのtoken boundaryをmodel tokenizer上で厳密に定義する必要がある。
BPE系tokenizerでは、
\begin{verbatim}
Response: 
\end{verbatim}
というprompt末尾と、
\begin{verbatim}
{"valence":1,"arousal":1}
\end{verbatim}
というcandidate先頭を別々にtokenizeした場合と、連結した全文をtokenizeした場合で、境界tokenが一致しない場合がある\citep{sennrich2016neural}。
これは、末尾spaceとcandidate先頭文字が一つのBPE tokenへmergeされ得るためである。
そのため本研究では、promptとcandidateを個別にtokenizeした後にtoken列を単純連結する方法を使用しない。
代わりに、
\[
\text{prompt}
+
\text{candidate}
\]
という完全な文字列をjoint tokenizationする。
さらにprompt末尾にASCII spaceが存在する場合には、完全な文字列を変化させることなく、そのspaceをcandidate側へ移動してtoken boundaryを正規化する。
例えば、
\begin{verbatim}
Original prompt:
...Response: 

Original candidate:
{"valence":1,"arousal":1}
\end{verbatim}
を、
\begin{verbatim}
Scoring prompt:
...Response:

Scoring candidate:
 {"valence":1,"arousal":1}
\end{verbatim}
として扱う。
このとき、
\[
\text{scoring prompt}
+
\text{scoring candidate}
=
\text{original prompt}
+
\text{original candidate}
\]
であり、モデルへ与える完全な文字列は変化しない。
candidate側へ移したleading spaceは、candidate tokenizationおよびlength normalizationにも含める。
tokenization時にはchat template等によって既に挿入されたmodel-specific special tokensを二重に追加しないため、
\[
\texttt{add\_special\_tokens=False}
\]
に相当する設定で統一する。


% ----------------------------------------------------------------
\subsection{内部表現の測定位置}
\label{app:hidden_state_measurement}

V1--V3では、Transformerのhidden representationをlayer-wiseに解析する。layer \(l\) は、embedding layerではなくTransformer blockのindexを表す。
prompt-level representationを解析する場合、paddingを除いた最後の有効prompt tokenを、
\[
t_{\mathrm{end}}
=
\max
\{
t:
\mathrm{attention\ mask}_t=1
\}
\]
とし、この位置をprompt-endと呼ぶ。prompt-endは、stimulusとtask instructionの処理が終了し、response generationへ移行する直前のmodel stateを表す。
V1では各Transformer blockのprompt-end output hidden stateを主に解析する。
V2およびV3のactivation extraction /causal interventionでは、各Transformer blockの
post-MLP residual stateを測定・操作する。したがってV2/V3では、
\[
h_{l,t}
\]
を、layer \(l\)、token position \(t\)におけるpost-MLP residual representationとして扱う。
V3のみ、prompt-endに加えてcandidate response生成中の複数semantic stagesも解析する。詳細はV3概要で述べる。


% ----------------------------------------------------------------
\subsection{Relative depthによるlayer位置の正規化}
\label{app:relative_depth}

model family間ではTransformer block数が異なる。そのため、absolute layer indexだけを比較すると、例えばLayer 14が各modelで異なる計算深度に対応する。
そこでlayer \(l\)、総Transformer block数 \(L\)について、
\[
d_l
=
\frac{l}{L-1}
\]
をrelative depthとして定義する。
\[
d_l=0
\]
は最初のTransformer block、
\[
d_l=1
\]
は最後のTransformer blockに対応する。
model family間でpeak position、center of mass、事前固定site等を比較する場合には、原則としてrelative depthを使用する。


% ----------------------------------------------------------------
\subsection{Stageごとの測定空間}
\label{app:measurement_spaces}

各Stageで使用するdataset、prompt、candidate space、representation measurementを表~\ref{tab:measurement_spaces}にまとめる。

\begin{table}[htbp]
\centering
\scriptsize
\caption{各Stageの共通測定仕様。Stage固有の解析方法および統計検定は
後続の各Stage概要で詳述する。}
\label{tab:measurement_spaces}
\begin{tabularx}{\linewidth}{
@{}
>{\raggedright\arraybackslash}p{1.2cm}
>{\raggedright\arraybackslash}p{2.1cm}
>{\raggedright\arraybackslash}p{2.0cm}
>{\raggedright\arraybackslash}p{1.7cm}
>{\raggedright\arraybackslash}X
@{}
}
\toprule
\textbf{Stage} &
\textbf{Data} &
\textbf{Models / format} &
\textbf{Output space} &
\textbf{Internal measurement} \\
\midrule
Behavioral &
EmoBank 1,000;
AIPsy 480 &
Base plain;
Instruct native chat &
729 VAD &
なし。behavioral outputを評価 \\

V1 &
EmoBank 1,000;
AIPsy &
Base plain;
Instruct native chat &
Phase A/B: N/A;
Phase C: 729 VAD &
prompt-end block output hidden states \\

V2 &
EmoBank 1,000 &
Base / Instruct;
Matched-Plain primary &
81 VA &
prompt-end post-MLP residual \\

V3 &
AIPsy 192 matched pairs &
Instruct native chat &
81 VA &
prompt-endおよびgeneration-stage post-MLP residual \\
\bottomrule
\end{tabularx}
\end{table}


% ----------------------------------------------------------------
\subsection{DiscoveryとConfirmationの分離}
\label{app:discovery_confirmation}

本研究のcausal experimentsでは、同じsampleをsite selectionと最終確認の双方へ無制限に使用しないことを基本方針とする。
V1 Phase Cでは、AIPsy Clinical--Neutralの192 matched pairsをpair単位で、
\[
50\%\text{ Discovery}
+
50\%\text{ Confirmation}
\]
に分割する。

Discovery subsetは、causal profileやcandidate siteの選択に使用し、選択後の主要なcross-task interventionおよびtask-specific specializationは独立したConfirmation subsetで評価する。
V3でも同様に、探索的site selectionとconfirmatory evaluationを分離する。
さらにV3のcross-family replicationでは、QwenをDiscovery modelとして使用し、Qwenで選択されたrelative depthやgeneration stageをLlama、Gemma、OLMoの結果を見る前に固定する。
confirmatory model familiesでは、そのoutcomeを見てsiteを再選択しない。
ただしhidden dimensionalityやrepresentation basisはmodel family間で異なるため、Qwenのhidden vectorそのものを他familyへ直接移植するわけではない。
各family内部で、frozen site条件のもとaffect-related directionを再推定する。


% ----------------------------------------------------------------
\subsection{乱数、推論精度、および再現性}
\label{app:reproducibility_common}

Primary experimentsでは、data splitting、control generation、bootstrap等の再現可能なrandom operationsについて原則として、
\[
\text{seed}=42
\]
を用いる。

model familyごとにBase / Instructのmodel identityとrevisionを固定し、異なるcheckpointへ切り替わらないようにする。Primary GPU inferenceにはbfloat16を使用する。

また、各実験では、dataset condition、model variant、prompt format、candidate space、
random seed等のexperiment-defining informationを固定し、異なる条件から得られた中間結果を同一実験結果として混在させない。


% ----------------------------------------------------------------
\subsection{解釈上の共通原則}
\label{app:interpretation_principles}

4 Stageを通して、以下の解釈上の境界を維持する。
第一に、Self-Reportは主観的意識や人間と同等の情動経験の存在を意味しない。
第二に、
\[
\operatorname{Corr}(R,S)>0
\]
あるいは、
\[
\operatorname{Corr}(\Delta R,\Delta S)>0
\]
だけから、
\[
R\rightarrow S
\]
というcausal relationを主張しない。
第三に、linear probeで高いheld-out performanceが得られても、
\[
\text{Decodability}
\neq
\text{Local Causal Leverage}
\]
である。

第四に、Base modelとInstruct modelは独立して得られたtrained artifactsであり、本研究はpost-training procedureそのものをrandomized interventionしていない。
したがってV2で観測されるmodel differenceを、
\[
\text{causal effect of post-training}
\]
と呼ばず、post-training-associated changeまたは、post-training-associated reorganizationと記述する。

最後に、Behavioral / V1で用いる729-state VAD protocolと、V2 / V3で用いる81-state VA protocolは異なるmeasurement spaceである。

このためStage間のabsolute affective scoreそのものではなく、各Stage内で事前定義されたcontrast、geometry、causal effect、distributional change、relative depth等を用いて
研究全体を統合する。
後続のBehavioral、V1、V2、V3概要では、本節で定義した共通仕様を前提として、各Stage固有のResearch Question、統計単位、intervention、control、およびPrimary endpointを詳述する。


% ================================================================
% D. Behavioral Stage
% ================================================================

\section{Behavioral概要}
\label{app:behavioral-methods}

\subsection{目的と位置づけ}
Behavioral Stageでは前節の共通仕様を用いて、内部representationやcausal interventionを解析する前に、Reader PredictionとSelf-Reportが外部の情動刺激に対して出力レベルでどのように変化するかを特徴づける。
Behavioral Stageの中心的なResearch Questionは、「制御された情動刺激の変化に対して、
Reader PredictionとSelf-Reportは連動して変化するか」である。

本Stageでは、単一の相関指標だけでReader--Self relationを判断せず，まずEmoBankを用いて、
モデルが人間の情動評定に対応するvariationを捉えているかを確認する。次にAIPsy-Affectの統制されたstimulus manipulationを用いて、
\begin{enumerate}
    \item \textbf{Sensitivity}：
    Clinical stimulusへの反応がmatched Neutralから
    期待方向へ変化するか、

    \item \textbf{Dose-response}：
    NeutralからModerate、Clinicalへ
    情動強度を増加させたときに段階的な変化が生じるか、

    \item \textbf{Specificity}：
    Clinical stimulusへの大きな変位が
    非情動的だが複雑なComplex Neutralでも同程度に生じるか、

    \item \textbf{Reader--Self Coupling}：
    同一matched pairに対するReaderの変化量とSelfの変化量が
    stimulus間で共変動するか
\end{enumerate}
を評価する。
したがってBehavioral Stageは、
\[
\text{Human Correspondence}
\rightarrow
\text{Controlled Affective Reactivity}
\rightarrow
\text{Reader--Self Change Coupling}
\]
という構成である。
ここで得られる結果はあくまでbehavioral levelのものであり、ReaderとSelfが同一のhidden representationやcausal pathwayを利用していることまでは意味しない。


% ----------------------------------------------------------------
\subsection{Behavioral Stageにおける解析対象}

Behavioral Stageでは、前節で定義した729-state VAD Sequence-Likelihood Protocolから得られる
\[
E[V],
\qquad
E[A],
\qquad
E[D]
\]
を主要な連続出力として使用する。

各stimulusについて、Writer、Reader、Selfを独立forward passとして測定し、
\[
W(x)
=
\left(
W_V(x),W_A(x),W_D(x)
\right),
\]

\[
R(x)
=
\left(
R_V(x),R_A(x),R_D(x)
\right),
\]

\[
S(x)
=
\left(
S_V(x),S_A(x),S_D(x)
\right)
\]
を得る。

Behavioral Stage全体におけるPrimary tasksは、
\[
\text{Reader}
\quad\text{and}\quad
\text{Self}
\]
であり、Primary affective dimensionsは、

\[
\text{Valence}
\quad\text{and}\quad
\text{Arousal}
\]
である。
Writerはhuman groundingの補助task、Dominanceは補助affective dimensionとして保持する。
したがってWriterやDominanceの結果も計算・保存するが、AIPsy-AffectにおけるPrimary inferential familyには含めない。


% ----------------------------------------------------------------
\subsection{Behavioral Stageで用いるプロンプト}
\label{app:behavioral-prompts}

Behavioral StageにおけるWriter、Reader、Selfの各taskプロンプトおよびVAD JSON出力フォーマットは、前節\S\ref{app:prompt_design}で定義した共通形式に完全準拠する。
すなわち、各taskに対応するinstruction（Writer-State Estimation、Reader-Response Prediction、Self-Report）と共通のVAD 1--9 response specificationを付与し、Instruct modelにはnative chat template、Base modelにはplain continuation formatを適用する。
これにより、ReaderとSelfの比較がresponse schemaそのものの差異に交絡されることを防止し、同一の729-state VAD candidate空間における評価を一貫して行う。


% ----------------------------------------------------------------
\subsection{EmoBankによるHuman-Affect Correspondence}
\label{app:behavioral-emobank}

Behavioral Stageではまず、EmoBankの1,000 test stimuliを用いて、モデルの情動関連出力が
human affective annotationのstimulus-wise variationをどの程度相関があるかを確認する。
この解析の目的は、AIPsy-Affect上でReader--Self couplingを議論する前に、モデルがそもそも外部の情動変動に対応したsystematic responseを生成しているかを確認することである。


\subsubsection{Writer correspondence}

Writerについては、EmoBankのhuman Writer VADとモデルWriter outputを比較する。
各dimension \(d\in\{V,A,D\}\) について、
\[
W_{d,\mathrm{model}}
\leftrightarrow
W_{d,\mathrm{human}}
\]
を評価する。WriterはBehavioral StageのPrimary Reader--Self comparisonには含まれないが、モデルが文章を書いた人物の情動状態をhuman annotationに対応する形で推定できるかを確認する補助的grounding taskとして使用する。


\subsubsection{Reader correspondence}

Readerについては、
\[
R_{d,\mathrm{model}}
\leftrightarrow
R_{d,\mathrm{human}}
\]
を直接比較する。
これは、モデルが予測する平均的な人間Readerの情動反応と、EmoBankで実際に得られたhuman Reader VADのstimulus-wise relationを評価するものである。


\subsubsection{Self--human Reader correspondence}

Selfについては、モデルSelf-Reportに対応するhuman ground truthは存在しない。
そのため、
\[
S_{d,\mathrm{model}}
\leftrightarrow
R_{d,\mathrm{human}}
\]
を\textbf{correspondence reference}として評価する。
この比較は、モデル自身について生成されるaffective reportが、人間Readerに観測されるstimulus間の情動変動をどの程度追跡するかを測定する。

したがって、このquantityをSelf-Reportの``accuracy''とは呼ばない。また、Selfとhuman Readerの値が一致することを主観的感情経験の証拠とも解釈しない。


% ----------------------------------------------------------------
\subsection{EmoBank correspondence metrics}

各model、task、dimensionについて、729-candidate distributionから得られる連続期待値と
human annotationの間で、
\[
r
=
\operatorname{PearsonCorr}
(
Y_{\mathrm{model}},
Y_{\mathrm{human}}
)
\]

および、
\[
\rho
=
\operatorname{SpearmanCorr}
(
Y_{\mathrm{model}},
Y_{\mathrm{human}}
)
\]
を計算する。

連続expected VADをPrimary measurementとし、candidate distributionのargmaxから得られるgreedy VADとのcorrelationはSecondaryとする。



\subsubsection{Reader--Self raw coupling on EmoBank}

EmoBankについては、モデル内部のReaderとSelfのstimulus-wise relationもSecondaryとして測定する。
各dimensionについて、
\[
\operatorname{Corr}
(
R_d(x),
S_d(x)
)
\]
を計算する。
さらに、
\[
S_d(x)-R_d(x)
\]
の平均と標準偏差を記録する。

ただしEmoBankにはClinical--Neutralのようなcontrolled manipulationが存在しない。
したがってこのraw correlationには、stimulus identityに由来する共通変動も含まれ得る。

この理由から、EmoBank上のraw Reader--Self correlationをBehavioral StageのPrimary coupling evidenceとはせず、AIPsy matched-pair change couplingをPrimaryとする。


% ----------------------------------------------------------------
\subsection{Neutral-response diagnostics}

Sequence-Likelihood expectationとは別に、729 candidatesの中で最大確率を持つgreedy VAD tripletも記録する。
とくにneutral candidate、
\[
(5,5,5)
\]
について、
\[
P(5,5,5)
\]
を各stimulusで保存し、さらにgreedy candidateが正確に
\[
(5,5,5)
\]
となったstimulusの割合を算出する。

これらは、モデル出力がneutral responseへ集中しているかを把握するためのdiagnostic quantitiesであり、Primary behavioral endpointには使用しない。


% ----------------------------------------------------------------
\subsection{AIPsy-Affectを用いた統制されたBehavioral Evaluation}
\label{app:behavioral-aipsy}

EmoBankによるhuman correspondenceの確認後、AIPsy-Affect 4-Splitを用いて、
制御されたstimulus manipulationに対するReaderおよびSelfの反応性を評価する。

前節で示したように、AIPsy-Affectにはcontinuous human VAD annotationが存在しない。
したがってAIPsyでは、
\[
\text{model prediction}
\leftrightarrow
\text{human continuous VAD}
\]
を評価するのではなく、
\[
\text{controlled stimulus manipulation}
\rightarrow
\text{model response change}
\]
を評価する。
Behavioral StageではAIPsy-Affectを用いて4つのResearch Questionsを検証する。


% ----------------------------------------------------------------
\subsection{A priori affect-direction specification}
\label{app:aipsy-direction}

Clinical conditionには、positive-valence emotionとnegative-valence emotionが混在する。
例えば、griefとecstasyを区別せず

\[
E[V]_{\mathrm{Clinical}}
-
E[V]_{\mathrm{Neutral}}
\]
を全stimuliで単純平均すると、negative directionとpositive directionが相殺され、情動刺激に対するsystematic responseが存在していても平均差が0に近づく可能性がある。

同様に、emotion categoryによってArousalやDominanceの期待変化方向も異なり得る。

そこで各emotion \(e\) とdimension \(d\) に対して、ClinicalがNeutralから変化すると期待される方向を、
\[
q_{e,d}
\in
\{-1,+1\}
\]
として事前固定する。
\[
q_{e,d}=+1
\]
はClinicalがNeutralより高い値へ変化すると期待することを、
\[
q_{e,d}=-1
\]
は低い値へ変化すると期待することを表す。
使用するdirection specificationを表~\ref{tab:aipsy-direction}に示す。

\begin{table}[htbp]
\centering
\small
\caption{AIPsy-AffectにおけるClinical--Neutral contrastの
事前定義された期待方向。
符号はClinicalがmatched Neutralに対して変化すると期待される方向を表す。}
\label{tab:aipsy-direction}
\begin{tabular}{lccc}
\toprule
\textbf{Emotion} &
\textbf{Valence} &
\textbf{Arousal} &
\textbf{Dominance} \\
\midrule
grief      & $-$ & $-$ & $-$ \\
terror     & $-$ & $+$ & $-$ \\
rage       & $-$ & $+$ & $+$ \\
loathing   & $-$ & $+$ & $+$ \\
ecstasy    & $+$ & $+$ & $+$ \\
admiration & $+$ & $+$ & $+$ \\
amazement  & $+$ & $+$ & $-$ \\
vigilance  & $+$ & $+$ & $+$ \\
\bottomrule
\end{tabular}
\end{table}

このdirection mapはhuman VAD ground truthではなく、emotion semanticsに基づいて事前固定したoperational hypothesisである。全modelで共通して使用し、得られたmodel outputを観察した後には変更しない。
Primary inferential analysisではReader / SelfのValence / Arousalを使用するため、Dominanceのdirection specificationはSecondary analysisに用いる。


% ----------------------------------------------------------------
\subsection{RQ1：Clinical--Neutral Sensitivity}
\label{app:behavioral-rq1}

RQ1では、「Clinical stimulusはmatched Neutralと比較して、期待される情動方向へmodel outputを変化させるか」を検証する。
Clinical--Neutral matched pair \(i\) とdimension \(d\) について、raw paired differenceを、
\[
\Delta_{i,d}
=
E_d(x_i^{C})
-
E_d(x_i^{N})
\]
と定義する。
ここで、
\[
x_i^{C}
\]
はClinical stimulus、
\[
x_i^{N}
\]
は同一\texttt{pair\_id}を持つmatched Neutral stimulusである。
emotion-specific directionを用いて、
\[
\Delta_{i,d}^{*}
=
q_{e_i,d}
\Delta_{i,d}
\]
をdirection-aligned differenceとする。
この定義により、
\[
\Delta_{i,d}^{*}>0
\]
はemotion categoryによらず、Clinical stimulusへのresponseが、事前定義された期待情動方向へNeutralから変化したことを意味する。


\subsubsection{Primary statistical test}

Primary statistical unitは、
\[
N=192
\]
Clinical--Neutral matched pairsである。
各model、task、dimensionについて、
\[
H_0:
E[\Delta^{*}]=0
\]
に対する\textbf{two-sided one-sample \(t\)-test}を行う。

Primary effect estimateは、
\[
\overline{\Delta^{*}}
\]
である。
effect sizeとしてpaired standardized effect、
\[
d_z
=
\frac{
\overline{\Delta^{*}}
}{
s_{\Delta^{*}}
}
\]
を計算する。

さらにmatched-pair effectsをreplacement付きでresampleするpercentile bootstrapを、
\[
B=1000,
\qquad
\text{seed}=42
\]
で実施し、
\[
\overline{\Delta^{*}}
\]
および \(d_z\) の95\% confidence intervalを求める。


\subsubsection{Secondary raw paired effect}

direction alignmentを行わない、
\[
\Delta_{i,d}
=
E_d(x_i^{C})
-
E_d(x_i^{N})
\]
もSecondaryとして保存する。

raw differenceについては、ClinicalとNeutralのpaired valuesに対するtwo-sided paired \(t\)-testとbootstrap confidence intervalを計算する。
ただしpositive / negative emotionの相殺を受けるため、Behavioral sensitivityに関するPrimary claimにはdirection-aligned quantityを使用する。


% ----------------------------------------------------------------
\subsection{RQ2：Dose-Response}
\label{app:behavioral-rq2}

RQ2では、「情動強度をNeutralからModerate、Clinicalへ増加させたとき、
model responseは段階的に期待方向へ変化するか」を検証する。

AIPsy-Affectの48 matched tripletsについて、
\[
x_i^{N}
\rightarrow
x_i^{M}
\rightarrow
x_i^{C}
\]
を使用する。
dimension \(d\) について、NeutralからModerateへのstepを、
\[
\Delta_{1,i}
=
E_d(x_i^{M})
-
E_d(x_i^{N}),
\]
ModerateからClinicalへのstepを、
\[
\Delta_{2,i}
=
E_d(x_i^{C})
-
E_d(x_i^{M})
\]
と定義する。
それぞれをa priori directionへalignし、
\[
\Delta_{1,i}^{*}
=
q_{e_i,d}\Delta_{1,i}
\]
および、
\[
\Delta_{2,i}^{*}
=
q_{e_i,d}\Delta_{2,i}
\]
をPrimary step effectsとする。


\subsubsection{Intersection--Union Test}

本研究では、NeutralからModerateまで変化しただけ、あるいはModerateからClinicalまで変化しただけではdose-response evidenceとはみなさない。
dose-response成立には、
\[
E[\Delta_1^{*}]>0
\]
かつ、
\[
E[\Delta_2^{*}]>0
\]
の両方を必要とする。
各stepについて、
\[
H_0:
E[\Delta_j^{*}]\le0
\]
に対する\textbf{one-sided one-sample \(t\)-test}を行う。得られたp-valueを、
\[
p_1,
\qquad
p_2
\]
とすると、Intersection--Union TestのPrimary p-valueを、
\[
p_{\mathrm{IUT}}
=
\max(p_1,p_2)
\]
と定義する。
したがって、二つのstepがともに十分なevidenceを持つ場合にのみdose-response criterionが満たされる。


\subsubsection{Effect description}

各stepのmean effect、
\[
\overline{\Delta_1^{*}},
\qquad
\overline{\Delta_2^{*}}
\]
について、48 tripletsを単位とした
\[
B=1000
\]
のpercentile bootstrap 95\% CIを算出する。また各tripletについて、
\[
\mathbb{I}
[
\Delta_{1,i}^{*}>0
\land
\Delta_{2,i}^{*}>0
]
\]
を定義し、その平均、
\[
M_{\mathrm{mono}}
=
\frac{1}{N}
\sum_i
\mathbb{I}
[
\Delta_{1,i}^{*}>0
\land
\Delta_{2,i}^{*}>0
]
\]
をmonotonicity rateとして報告する。


\subsubsection{Secondary slope analysis}

Secondary analysisとして、Neutral、Moderate、Clinicalを等間隔な強度水準
\[
0,1,2
\]
とみなしたときのendpoint slopeに対応する、
\[
b_i
=
\frac{
E_d(x_i^{C})
-
E_d(x_i^{N})
}{2}
\]
をtripletごとに計算する。
このraw slopeについて、two-sided one-sample \(t\)-testおよびbootstrap 95\% CIを算出する。またModerate responseがNeutralとClinicalの中点からどの程度ずれるかを、
\[
E_d(x_i^{M})
-
\frac{
E_d(x_i^{N})
+
E_d(x_i^{C})
}{2}
\]
として記録し、dose-response shapeの補助的descriptionとして使用する。
これらはSecondaryであり、Primary inferential quantityは二段階のdirection-aligned IUTである。


% ----------------------------------------------------------------
\subsection{RQ3：Specificity}
\label{app:behavioral-rq3}

RQ3では、「Clinical stimulusによるoutput変位は、単なる複雑な文章への反応だけで説明できるか」を検証する。
Clinicalにはpositive / negative emotionが混在するため、Valence等のraw meanだけをClinicalとComplex Neutralの間で比較すると、方向の異なる情動反応が相殺される可能性がある。

そこでまず、各model、task、dimensionについて192 Neutral stimuli全体の平均を、
\[
\mu_N
=
\frac{1}{N_N}
\sum_{x\in N}
E_d(x)
\]
とする。
これは各Clinicalに個別対応するmatched Neutral valueではなく、当該model / task / dimensionにおけるNeutral condition全体の平均である。
各stimulus \(x\) のNeutral baselineからのabsolute displacementを、
\[
D(x)
=
|
E_d(x)-\mu_N
|
\]
と定義する。


\subsubsection{Clinical versus Complex Neutral}

Clinical stimuliについて、
\[
D_C
=
\{
D(x):
x\in\mathrm{Clinical}
\},
\]
Complex Neutralについて、
\[
D_{CN}
=
\{
D(x):
x\in\mathrm{ComplexNeutral}
\}
\]
を求める。

sample sizesは、
\[
N_C=192,
\qquad
N_{CN}=48
\]
である。

Complex NeutralはClinicalとmatched pairを構成しないため、paired testは使用しない。
Primary testとして、
\[
\text{two-sided Welch independent-samples }t\text{-test}
\]
を実施し、
\[
H_0:
E[D_C]=E[D_{CN}]
\]
を検定する。
Primary effect estimateは、
\[
\Delta D
=
E[D_C]-E[D_{CN}]
\]
である。
\[
\Delta D>0
\]
であれば、Clinical stimuliがComplex NeutralよりNeutral baselineから大きく離れたresponseを生じさせることを意味する。


\subsubsection{Effect size}

Specificityのeffect sizeとして、
\[
d_{\mathrm{spec}}
=
\frac{
\overline{D_C}
-
\overline{D_{CN}}
}{
\sqrt{
\left(
s_C^2+s_{CN}^2
\right)/2
}
}
\]
を計算する。

またSecondary analysisとして、absolute displacementではなくClinicalとComplex Neutralのraw model outputを直接比較したtwo-sided Welch testも保存する。
ただしClinical内のpositive / negative emotion cancellationを受けるため、Primary specificity conclusionにはabsolute displacement comparisonを用いる。


% ----------------------------------------------------------------
\subsection{RQ4：Reader--Self Change Coupling}
\label{app:behavioral-rq4}

Behavioral Stageの中心となるReader--Self analysisでは、ReaderとSelfの値そのものが相関するかではなく、同一刺激操作によってReaderとSelfがどの程度一緒に変化するかをPrimary questionとする。

Clinical--Neutral matched pair \(i\) について、Reader changeを、
\[
\Delta R_{i,d}
=
R_d(x_i^{C})
-
R_d(x_i^{N})
\]

Self changeを、
\[
\Delta S_{i,d}
=
S_d(x_i^{C})
-
S_d(x_i^{N})
\]
と定義する。
dimensionごとに、
\[
r_{\Delta,d}
=
\operatorname{PearsonCorr}
(
\Delta R_d,
\Delta S_d
)
\]
をPrimary Reader--Self coupling measureとする。
ここではRQ1/RQ2とは異なり、\(\Delta R\)および\(\Delta S\)をemotion-specific directionへalignしない。
これは、ReaderとSelfが同じmatched manipulationに対して同じ方向・大きさで変化するかというsample-wise correspondenceそのものを測定するためである。


\subsubsection{Statistical inference}

統計単位は192 matched pairsである。Pearson correlationについて、
\[
B=1000,
\qquad
\text{seed}=42
\]
のpair bootstrapを行う。

bootstrapでは、各resampleで、
\[
(
\Delta R_i,
\Delta S_i
)
\]
のpairをjointly resampleする。

percentile distributionの2.5thおよび97.5th percentileから95\% confidence intervalを求める。
さらにSecondary rank-based measureとして、
\[
\rho_{\Delta,d}
=
\operatorname{SpearmanCorr}
(
\Delta R_d,
\Delta S_d
)
\]
も計算する。


\subsubsection{Raw coupling as a secondary measure}

比較のため、Clinical、Neutral、Moderate、Complex Neutralを含むAIPsy-Affect全480 stimuliについて、
\[
\operatorname{Corr}
(
R_d(x),
S_d(x)
)
\]
というraw cross-stimulus correlationも算出する。
ただしraw correlationは、stimulusそのものの内容やdifficulty等のshared variationを含み得る。したがって、Raw Reader--Self correlationはSecondaryとし、
\[
\operatorname{Corr}
(
\Delta R,
\Delta S
)
\]
をPrimary Behavioral coupling evidenceとする。


% ----------------------------------------------------------------
\subsection{Primary analysesとSecondary analysesの区別}
\label{app:behavioral-primary-secondary}

Behavioral Stageでは、Primary analysesを以下のように事前に限定する。(表~\ref{tab:behavioral-primary-secondary})

\begin{table}[htbp]
\centering
\small
\caption{Behavioral StageにおけるPrimaryおよび主要Secondary analyses。}
\label{tab:behavioral-primary-secondary}
\begin{tabularx}{\linewidth}{
@{}
>{\raggedright\arraybackslash}p{2.2cm}
>{\raggedright\arraybackslash}p{3.3cm}
>{\raggedright\arraybackslash}X
@{}
}
\toprule
\textbf{Analysis} &
\textbf{Primary} &
\textbf{Secondary / Diagnostic} \\
\midrule
EmoBank correspondence &
Reader / Self $\times$ V/A,
continuous expectation &
Writer、Dominance、greedy VAD、
neutral-collapse diagnostics \\

AIPsy Sensitivity &
Reader / Self $\times$ V/A,
direction-aligned mean &
Writer、Dominance、raw signed difference \\

AIPsy Dose-response &
Reader / Self $\times$ V/A,
two-step IUT &
Writer、Dominance、raw endpoint slope、
midpoint deviation \\

AIPsy Specificity &
Reader / Self $\times$ V/A,
absolute displacement &
Writer、Dominance、raw signed mean comparison \\

Reader--Self coupling &
V/A delta coupling &
Dominance delta coupling、
raw 480-stimulus coupling \\
\bottomrule
\end{tabularx}
\end{table}


% ----------------------------------------------------------------
\subsection{多重比較補正}
\label{app:behavioral-fdr}

AIPsy-AffectにおけるPrimary inferential testsには、Benjamini--Hochberg法によるFalse Discovery Rate（FDR）補正を適用する\citep{benjamini1995controlling}。補正はRQごとに独立したinferential familyとして行う。

Primary familyは、
\[
8\text{ models}
\times
2\text{ tasks}
\times
2\text{ dimensions}
=
32
\]
testsからなる。

Primary IUTについても、
\[
8
\times
2
\times
2
=
32
\]
testsを一つのfamilyとして補正する。

Specificityについても、Reader / SelfのValence / Arousalに対する、
\[
8
\times
2
\times
2
=
32
\]
testsを一つのfamilyとして補正する。

Primary delta couplingについてはtask軸が存在しないため、
\[
8\text{ models}
\times
2\text{ dimensions}
=
16
\]
testsを一つのfamilyとして補正する。

WriterおよびDominance、raw coupling等のSecondary analysesはこれらPrimary FDR familiesには含めない。


% ----------------------------------------------------------------
\subsection{Bootstrapと欠損値の扱い}

Behavioral Stageで使用するbootstrap confidence intervalは、
\[
B=1000,
\qquad
\text{seed}=42,
\qquad
95\%\text{ CI}
\]
とする。

mean effectのbootstrapではsampleまたはmatched pairをreplacement付きでresampleする。
Reader--Self correlationでは、対応する
\[
(x_i,y_i)
\]
のpairをjointly resampleする。

bootstrap distributionの2.5thおよび97.5th percentilesをconfidence boundsとする。

correlationを定義できるだけのsample variationが存在しない場合など、統計量が定義不能な条件については、値を0へ置換せずmissing valueとして保持する。



% ================================================================
% 5. Behavioral結果
% ================================================================
\section{Behavioral結果}
\label{sec:behavioral-results}

\begin{table}[htbp]
\centering
\small
\caption{EmoBank 3者間VADアライメント：問い「人間の言語表現（Writer）および読者評価（Reader）のVADグラウンドトゥルースに対して、LLMの認識・自己報告（Self）はどの程度整合するか」。4モデルファミリーのBaseおよびInstructモデルにおける人間正解ラベル（Writer / Reader）および自己報告（Self）のピアソン相関係数 $r$（$N=1000$）。}
\label{tab:behavioral_emobank_3way_vad}
\begin{tabular}{lll ccc}
\toprule
\textbf{Family} & \textbf{Model Variant} & \textbf{Perspective / Task} & \textbf{Valence ($r$)} & \textbf{Arousal ($r$)} & \textbf{Dominance ($r$)} \\
\midrule
\multirow{6}{*}{Qwen 2.5 1.5B} & \multirow{3}{*}{Base} & Writer     & 0.393$^{***}$ & 0.156$^{***}$ & 0.111$^{***}$ \\
                       &                      & Reader     & 0.472$^{***}$ & 0.257$^{***}$ & 0.227$^{***}$ \\
                       &                      & Self       & 0.332$^{***}$ & 0.218$^{***}$ & 0.224$^{***}$ \\
\cmidrule(lr){2-6}
                       & \multirow{3}{*}{Instruct} & Writer     & 0.533$^{***}$ & 0.156$^{***}$ & 0.060$$ \\
                       &                      & Reader     & 0.587$^{***}$ & 0.293$^{***}$ & 0.030$$ \\
                       &                      & Self       & 0.537$^{***}$ & 0.251$^{***}$ & 0.109$^{***}$ \\
\midrule
\multirow{6}{*}{Llama 3.2 1B} & \multirow{3}{*}{Base} & Writer     & 0.121$^{***}$ & -0.010$$ & -0.021$$ \\
                       &                      & Reader     & 0.129$^{***}$ & 0.123$^{***}$ & 0.025$$ \\
                       &                      & Self       & 0.109$^{***}$ & 0.074$^{*}$ & 0.017$$ \\
\cmidrule(lr){2-6}
                       & \multirow{3}{*}{Instruct} & Writer     & 0.174$^{***}$ & 0.174$^{***}$ & 0.146$^{***}$ \\
                       &                      & Reader     & 0.098$^{**}$ & 0.297$^{***}$ & 0.216$^{***}$ \\
                       &                      & Self       & 0.045$$ & 0.209$^{***}$ & 0.254$^{***}$ \\
\midrule
\multirow{6}{*}{Gemma 3 1B} & \multirow{3}{*}{Base} & Writer     & -0.022$$ & 0.029$$ & 0.055$$ \\
                       &                      & Reader     & -0.004$$ & 0.053$$ & 0.050$$ \\
                       &                      & Self       & 0.017$$ & 0.024$$ & -0.001$$ \\
\cmidrule(lr){2-6}
                       & \multirow{3}{*}{Instruct} & Writer     & -0.015$$ & 0.209$^{***}$ & 0.135$^{***}$ \\
                       &                      & Reader     & -0.012$$ & 0.320$^{***}$ & 0.166$^{***}$ \\
                       &                      & Self       & 0.105$^{***}$ & 0.254$^{***}$ & 0.179$^{***}$ \\
\midrule
\multirow{6}{*}{OLMo 2 1B} & \multirow{3}{*}{Base} & Writer     & 0.023$$ & -0.025$$ & -0.012$$ \\
                       &                      & Reader     & 0.015$$ & -0.013$$ & 0.039$$ \\
                       &                      & Self       & 0.026$$ & -0.001$$ & 0.049$$ \\
\cmidrule(lr){2-6}
                       & \multirow{3}{*}{Instruct} & Writer     & 0.323$^{***}$ & 0.096$^{**}$ & -0.005$$ \\
                       &                      & Reader     & 0.338$^{***}$ & 0.089$^{**}$ & 0.104$^{**}$ \\
                       &                      & Self       & 0.440$^{***}$ & 0.075$^{*}$ & 0.127$^{***}$ \\
\bottomrule
\end{tabular}
\vspace{1ex}
\begin{minipage}{\linewidth}
\footnotesize
\textbf{Note:} $^{***}: p < 0.001, ^{**}: p < 0.01, ^{*}: p < 0.05$。BaseモデルではQwenおよびLlamaのみが正の相関を示したが、事後学習（Instruct）により全ファミリーにおいて認識（Writer/Reader）および自己報告（Self）のアライメントが系統的に向上した。
\end{minipage}
\end{table}
% ================================================================
% AIPsy-Affect Complete RQ Summary Tables (LaTeX Version)
% Generated from behavioral_aipsy_summary.md (RQ1--RQ4 Complete)
% ================================================================

\begin{table}[htbp]
\centering
\small
\caption{AIPsy RQ1（Clinical--Neutral感度）：問い「臨床刺激は統制中立刺激と比較して、期待される情動方向へのモデル出力変位を引き起こすか」。感情対照ペア（$N=192$）におけるモデル整列変位 $\Delta_{\text{aligned}}$、95\%信頼区間、効果量 Cohen's $d_z$、およびFDR補正後 $q$ 値。}
\label{tab:behavioral_aipsy_rq1_sensitivity}
\begin{tabular}{lll ccc ccc}
\toprule
 & & & \multicolumn{3}{c}{\textbf{Valence Axis}} & \multicolumn{3}{c}{\textbf{Arousal Axis}} \\
\cmidrule(lr){4-6} \cmidrule(lr){7-9}
\textbf{Family} & \textbf{Variant} & \textbf{Task} & $\Delta_{\text{aligned}}$ [95\% CI] & $d_z$ & $q$-value & $\Delta_{\text{aligned}}$ [95\% CI] & $d_z$ & $q$-value \\
\midrule
\multirow{6}{*}{Qwen 2.5 1.5B} & \multirow{3}{*}{Base} & Writer     & 0.007 [0.004, 0.009] & 0.393 & --- & 0.011 [0.008, 0.013] & 0.491 & --- \\
                       &                      & Reader     & 0.005 [0.002, 0.007] & 0.271 & $<0.001$ & 0.015 [0.012, 0.018] & 0.649 & $<10^{-15}$ \\
                       &                      & Self       & 0.006 [0.004, 0.008] & 0.446 & $<10^{-8}$ & 0.005 [0.003, 0.008] & 0.269 & $<0.001$ \\
\cmidrule(lr){2-9}
                       & \multirow{3}{*}{Instruct} & Writer     & 0.019 [0.014, 0.025] & 0.495 & --- & 0.038 [0.032, 0.045] & 0.816 & --- \\
                       &                      & Reader     & 0.022 [0.015, 0.030] & 0.415 & $<10^{-7}$ & 0.065 [0.056, 0.074] & 0.953 & $<10^{-15}$ \\
                       &                      & Self       & 0.018 [0.012, 0.022] & 0.493 & $<10^{-10}$ & 0.033 [0.027, 0.040] & 0.728 & $<10^{-15}$ \\
\midrule
\multirow{6}{*}{Llama 3.2 1B} & \multirow{3}{*}{Base} & Writer     & -0.000 [-0.002, 0.002] & -0.009 & --- & 0.001 [0.000, 0.002] & 0.203 & --- \\
                       &                      & Reader     & 0.001 [-0.000, 0.003] & 0.109 & 0.1954 & 0.002 [0.001, 0.003] & 0.234 & 0.0026 \\
                       &                      & Self       & 0.000 [-0.001, 0.002] & 0.039 & 0.6691 & 0.002 [0.001, 0.003] & 0.241 & 0.0021 \\
\cmidrule(lr){2-9}
                       & \multirow{3}{*}{Instruct} & Writer     & 0.005 [0.001, 0.009] & 0.167 & --- & 0.025 [0.019, 0.031] & 0.565 & --- \\
                       &                      & Reader     & 0.001 [-0.004, 0.006] & 0.016 & 0.8821 & 0.071 [0.061, 0.081] & 0.972 & $<10^{-15}$ \\
                       &                      & Self       & 0.007 [0.003, 0.010] & 0.241 & 0.0021 & 0.019 [0.012, 0.024] & 0.417 & $<10^{-7}$ \\
\midrule
\multirow{6}{*}{Gemma 3 1B} & \multirow{3}{*}{Base} & Writer     & -0.009 [-0.017, -0.003] & -0.181 & --- & -0.006 [-0.012, -0.001] & -0.168 & --- \\
                       &                      & Reader     & -0.005 [-0.014, 0.003] & -0.089 & 0.2907 & -0.002 [-0.008, 0.003] & -0.068 & 0.4121 \\
                       &                      & Self       & -0.002 [-0.011, 0.006] & -0.033 & 0.7168 & -0.007 [-0.012, -0.001] & -0.175 & 0.0278 \\
\cmidrule(lr){2-9}
                       & \multirow{3}{*}{Instruct} & Writer     & -0.011 [-0.031, 0.007] & -0.083 & --- & 0.017 [-0.004, 0.040] & 0.103 & --- \\
                       &                      & Reader     & 0.000 [-0.014, 0.016] & 0.004 & 0.9576 & 0.116 [0.092, 0.142] & 0.647 & $<10^{-15}$ \\
                       &                      & Self       & 0.001 [-0.016, 0.019] & 0.010 & 0.9165 & 0.012 [-0.011, 0.038] & 0.068 & 0.4121 \\
\midrule
\multirow{6}{*}{OLMo 2 1B} & \multirow{3}{*}{Base} & Writer     & -0.009 [-0.014, -0.003] & -0.232 & --- & -0.002 [-0.005, -0.000] & -0.141 & --- \\
                       &                      & Reader     & -0.006 [-0.012, -0.000] & -0.150 & 0.0619 & -0.003 [-0.006, -0.001] & -0.190 & 0.0161 \\
                       &                      & Self       & -0.003 [-0.005, -0.000] & -0.141 & 0.0783 & -0.001 [-0.002, 0.001] & -0.071 & 0.4121 \\
\cmidrule(lr){2-9}
                       & \multirow{3}{*}{Instruct} & Writer     & 0.010 [0.007, 0.014] & 0.387 & --- & -0.023 [-0.034, -0.012] & -0.294 & --- \\
                       &                      & Reader     & 0.008 [0.004, 0.011] & 0.319 & $<10^{-5}$ & -0.007 [-0.017, 0.003] & -0.104 & 0.2083 \\
                       &                      & Self       & 0.009 [0.005, 0.013] & 0.357 & $<10^{-6}$ & -0.017 [-0.027, -0.007] & -0.242 & 0.0021 \\
\bottomrule
\end{tabular}
\vspace{1ex}
\begin{minipage}{\linewidth}
\footnotesize
\textbf{Note:} 整列変位 $\Delta_{\text{aligned}} > 0$ は刺激の極性に応じた期待方向への変位を示す。Instructモデルでは、Arousal軸（Qwen, Llama, Gemma）およびValence軸（Qwen, OLMo）において強固な感度（$d_z > 0.4, q < 0.001$）が確認される。
\end{minipage}
\end{table}

\begin{table}[htbp]
\centering
\small
\caption{AIPsy RQ2（用量反応性・単調性）：問い「情動強度をNeutralからModerate、Clinicalへ段階的に増加させたとき、モデルの出力変位は単調に増加するか」。強度3段階トリプレット（$N=48$）における第1段階変位 $\Delta_1$（Neu $\to$ Mod）、第2段階変位 $\Delta_2$（Mod $\to$ Clin）、Intersection-Union Test (IUT) $q$ 値、および単調性充足率。}
\label{tab:behavioral_aipsy_rq2_dose_response}
\begin{tabular}{lll cccc cccc}
\toprule
 & & & \multicolumn{4}{c}{\textbf{Valence Axis}} & \multicolumn{4}{c}{\textbf{Arousal Axis}} \\
\cmidrule(lr){4-7} \cmidrule(lr){8-11}
\textbf{Family} & \textbf{Variant} & \textbf{Task} & $\Delta_1$ & $\Delta_2$ & $q_{\text{IUT}}$ & Mono.\% & $\Delta_1$ & $\Delta_2$ & $q_{\text{IUT}}$ & Mono.\% \\
\midrule
\multirow{4}{*}{Qwen 2.5 1.5B} & \multirow{2}{*}{Base} & Reader     & 0.004 & -0.003 & 0.9998 & 22.9\% & 0.003 & 0.016 & 0.9998 & 29.2\% \\
                       &                      & Self       & 0.003 & 0.001 & 0.9998 & 16.7\% & -0.003 & 0.009 & 0.9998 & 22.9\% \\
\cmidrule(lr){2-11}
                       & \multirow{2}{*}{Instruct} & Reader     & 0.024 & -0.007 & 0.9998 & 22.9\% & 0.021 & 0.062 & 0.0952 & 47.9\% \\
                       &                      & Self       & 0.020 & -0.001 & 0.9998 & 18.8\% & -0.001 & 0.042 & 0.9998 & 33.3\% \\
\midrule
\multirow{4}{*}{Llama 3.2 1B} & \multirow{2}{*}{Base} & Reader     & -0.002 & 0.001 & 0.9998 & 12.5\% & -0.002 & 0.004 & 0.9998 & 18.8\% \\
                       &                      & Self       & -0.003 & 0.003 & 0.9998 & 16.7\% & -0.003 & 0.006 & 0.9998 & 27.1\% \\
\cmidrule(lr){2-11}
                       & \multirow{2}{*}{Instruct} & Reader     & 0.009 & -0.015 & 0.9998 & 10.4\% & 0.037 & 0.054 & 0.0048 & 54.2\% \\
                       &                      & Self       & 0.016 & -0.010 & 0.9998 & 20.8\% & -0.009 & 0.030 & 0.9998 & 18.8\% \\
\midrule
\multirow{4}{*}{Gemma 3 1B} & \multirow{2}{*}{Base} & Reader     & -0.013 & 0.005 & 0.9998 & 6.2\% & 0.008 & -0.010 & 0.9998 & 10.4\% \\
                       &                      & Self       & -0.015 & 0.013 & 0.9998 & 2.1\% & 0.004 & -0.015 & 0.9998 & 14.6\% \\
\cmidrule(lr){2-11}
                       & \multirow{2}{*}{Instruct} & Reader     & 0.012 & 0.016 & 0.9998 & 16.7\% & 0.050 & 0.124 & 0.3233 & 33.3\% \\
                       &                      & Self       & 0.023 & -0.013 & 0.9998 & 14.6\% & -0.042 & 0.073 & 0.9998 & 14.6\% \\
\midrule
\multirow{4}{*}{OLMo 2 1B} & \multirow{2}{*}{Base} & Reader     & -0.011 & -0.001 & 0.9998 & 10.4\% & -0.010 & 0.002 & 0.9998 & 8.3\% \\
                       &                      & Self       & -0.003 & -0.001 & 0.9998 & 14.6\% & -0.005 & 0.001 & 0.9998 & 8.3\% \\
\cmidrule(lr){2-11}
                       & \multirow{2}{*}{Instruct} & Reader     & 0.005 & 0.006 & 0.9998 & 22.9\% & -0.029 & 0.016 & 0.9998 & 6.2\% \\
                       &                      & Self       & 0.004 & 0.007 & 0.9998 & 25.0\% & -0.035 & 0.007 & 0.9998 & 8.3\% \\
\bottomrule
\end{tabular}
\vspace{1ex}
\begin{minipage}{\linewidth}
\footnotesize
\textbf{Note:} 単調性（IUT $q < 0.05$）は、刺激の感情強度が強まるにつれて出力が連続的に増加することを示す。Llama Instruct ReaderのArousal軸で確認された。
\end{minipage}
\end{table}

\begin{table}[htbp]
\centering
\small
\caption{AIPsy RQ3（感情特異性 vs 構文複雑性）：問い「臨床刺激による変位は、単なる複雑な構文の中立文（Complex Neutral）への反応ではなく、感情内容に特異的か」。臨床文変位 $\Delta_{\text{clin}}$ と複雑中立文変位 $\Delta_{\text{comp}}$ の差分、効果量 Cohen's $d$、およびFDR補正後 $q$ 値。}
\label{tab:behavioral_aipsy_rq3_specificity}
\begin{tabular}{lll cccc cccc}
\toprule
 & & & \multicolumn{4}{c}{\textbf{Valence Axis}} & \multicolumn{4}{c}{\textbf{Arousal Axis}} \\
\cmidrule(lr){4-7} \cmidrule(lr){8-11}
\textbf{Family} & \textbf{Variant} & \textbf{Task} & $\Delta_{\text{clin}}$ & $\Delta_{\text{comp}}$ & $d$ & $q$-value & $\Delta_{\text{clin}}$ & $\Delta_{\text{comp}}$ & $d$ & $q$-value \\
\midrule
\multirow{4}{*}{Qwen 2.5 1.5B} & \multirow{2}{*}{Base} & Reader     & 0.014 & 0.014 & -0.02 & 0.9464 & 0.020 & 0.009 & 0.96 & $<10^{-11}$ \\
                       &                      & Self       & 0.011 & 0.013 & -0.30 & 0.0923 & 0.015 & 0.009 & 0.72 & $<10^{-6}$ \\
\cmidrule(lr){2-11}
                       & \multirow{2}{*}{Instruct} & Reader     & 0.048 & 0.026 & 0.74 & $<10^{-6}$ & 0.078 & 0.053 & 0.58 & $<0.001$ \\
                       &                      & Self       & 0.029 & 0.023 & 0.27 & 0.0995 & 0.044 & 0.043 & 0.01 & 0.9511 \\
\midrule
\multirow{4}{*}{Llama 3.2 1B} & \multirow{2}{*}{Base} & Reader     & 0.009 & 0.015 & -0.69 & $<0.001$ & 0.007 & 0.012 & -0.77 & $<10^{-5}$ \\
                       &                      & Self       & 0.011 & 0.022 & -0.92 & $<10^{-5}$ & 0.008 & 0.012 & -0.68 & $<0.001$ \\
\cmidrule(lr){2-11}
                       & \multirow{2}{*}{Instruct} & Reader     & 0.029 & 0.026 & 0.15 & 0.3774 & 0.084 & 0.040 & 1.04 & $<10^{-9}$ \\
                       &                      & Self       & 0.024 & 0.029 & -0.23 & 0.2456 & 0.035 & 0.048 & -0.43 & 0.0244 \\
\midrule
\multirow{4}{*}{Gemma 3 1B} & \multirow{2}{*}{Base} & Reader     & 0.034 & 0.066 & -0.76 & $<0.001$ & 0.027 & 0.040 & -0.54 & 0.0049 \\
                       &                      & Self       & 0.033 & 0.061 & -0.72 & $<0.001$ & 0.024 & 0.037 & -0.55 & 0.0050 \\
\cmidrule(lr){2-11}
                       & \multirow{2}{*}{Instruct} & Reader     & 0.087 & 0.177 & -0.97 & $<10^{-6}$ & 0.165 & 0.136 & 0.24 & 0.1572 \\
                       &                      & Self       & 0.096 & 0.121 & -0.33 & 0.0842 & 0.146 & 0.162 & -0.15 & 0.3748 \\
\midrule
\multirow{4}{*}{OLMo 2 1B} & \multirow{2}{*}{Base} & Reader     & 0.030 & 0.019 & 0.55 & $<10^{-5}$ & 0.012 & 0.016 & -0.39 & 0.0410 \\
                       &                      & Self       & 0.014 & 0.011 & 0.23 & 0.1484 & 0.007 & 0.008 & -0.16 & 0.3592 \\
\cmidrule(lr){2-11}
                       & \multirow{2}{*}{Instruct} & Reader     & 0.023 & 0.028 & -0.24 & 0.2218 & 0.054 & 0.041 & 0.34 & 0.0329 \\
                       &                      & Self       & 0.022 & 0.026 & -0.18 & 0.3489 & 0.057 & 0.037 & 0.48 & 0.0038 \\
\bottomrule
\end{tabular}
\vspace{1ex}
\begin{minipage}{\linewidth}
\footnotesize
\textbf{Note:} 特異的効果量 $d$ は、臨床刺激に対する反応が単なる構文複雑さへの反応を超越している度合いを示す。Instructモデルにおいて、特にArousal軸で有意な感情特異性が維持されている。
\end{minipage}
\end{table}

\begin{table}[htbp]
\centering
\small
\caption{AIPsy RQ4（内部認知結合度）：問い「刺激提示に伴う他者認識の変位 $\Delta R$ と自己報告の変位 $\Delta S$ は、モデル内部で連動して結合（カップリング）しているか」。感情対照ペア（$N=192$）における他者認識変位 $\Delta\text{Reader}$ と自己報告変位 $\Delta\text{Self}$ の相関係数 $\Delta r$（95\%ブートストラップ信頼区間およびFDR補正後 $q$ 値）。}
\label{tab:behavioral_aipsy_rq4_coupling}
\begin{tabular}{ll cccc}
\toprule
\textbf{Family} & \textbf{Variant} & \textbf{Axis} & \textbf{$\Delta$ Pearson $r$ [95\% CI]} & \textbf{$q$-value} & \textbf{Raw $r$} \\
\midrule
\multirow{4}{*}{Qwen 2.5 1.5B} & \multirow{2}{*}{Base} & Valence    & 0.809 [0.755, 0.854] & $<10^{-15}$ & 0.838 \\
                       &                      & Arousal    & 0.866 [0.832, 0.892] & $<10^{-15}$ & 0.828 \\
\cmidrule(lr){2-6}
                       & \multirow{2}{*}{Instruct} & Valence    & 0.892 [0.859, 0.918] & $<10^{-15}$ & 0.898 \\
                       &                      & Arousal    & 0.917 [0.892, 0.938] & $<10^{-15}$ & 0.898 \\
\midrule
\multirow{4}{*}{Llama 3.2 1B} & \multirow{2}{*}{Base} & Valence    & 0.662 [0.570, 0.733] & $<10^{-15}$ & 0.864 \\
                       &                      & Arousal    & 0.876 [0.839, 0.904] & $<10^{-15}$ & 0.941 \\
\cmidrule(lr){2-6}
                       & \multirow{2}{*}{Instruct} & Valence    & 0.756 [0.675, 0.817] & $<10^{-15}$ & 0.799 \\
                       &                      & Arousal    & 0.754 [0.697, 0.807] & $<10^{-15}$ & 0.718 \\
\midrule
\multirow{4}{*}{Gemma 3 1B} & \multirow{2}{*}{Base} & Valence    & 0.813 [0.704, 0.888] & $<10^{-15}$ & 0.872 \\
                       &                      & Arousal    & 0.782 [0.685, 0.853] & $<10^{-15}$ & 0.825 \\
\cmidrule(lr){2-6}
                       & \multirow{2}{*}{Instruct} & Valence    & 0.621 [0.517, 0.713] & $<10^{-15}$ & 0.643 \\
                       &                      & Arousal    & 0.780 [0.660, 0.857] & $<10^{-15}$ & 0.756 \\
\midrule
\multirow{4}{*}{OLMo 2 1B} & \multirow{2}{*}{Base} & Valence    & 0.559 [0.444, 0.660] & $<10^{-15}$ & 0.668 \\
                       &                      & Arousal    & 0.600 [0.469, 0.708] & $<10^{-15}$ & 0.637 \\
\cmidrule(lr){2-6}
                       & \multirow{2}{*}{Instruct} & Valence    & 0.851 [0.811, 0.885] & $<10^{-15}$ & 0.903 \\
                       &                      & Arousal    & 0.931 [0.908, 0.950] & $<10^{-15}$ & 0.935 \\
\bottomrule
\end{tabular}
\vspace{1ex}
\begin{minipage}{\linewidth}
\footnotesize
\textbf{Note:} 全条件で $q < 10^{-15}$ の強固な結合を示し、事後学習を経てもモデル内部における感情認識と自己報告の連動ダイナミクスが破綻せず一貫して保たれていることが実証される。
\end{minipage}
\end{table}



% ================================================================
% 6. Behavioral考察
% ================================================================
\section{Behavioral考察}
\label{sec:behavioral-discussion}

Reader PredictionとSelf-Reportの間に、model familyおよびBase / Instructの別を超えて強いstimulus-level couplingが見られた。AIPsy-AffectのClinical--Neutral matched pairsでは、Reader changeとSelf changeのPearson correlationはValence / Arousalの全Primary conditionsで正となり、おおむね
\[
r_{\Delta}=0.56\text{--}0.93
\]
の範囲にあった。

この結果は、同じcontrolled stimulus manipulationによってReader Predictionが大きく変化するstimulusほど、Self-Reportも大きく変化するという安定したbehavioral covariationを示す。特に、本解析はabsolute responseそのものではなく、各Clinical stimulusとmatched Neutralとの差分、
\[
\Delta R_i
=
R(x_i^C)-R(x_i^N),
\qquad
\Delta S_i
=
S(x_i^C)-S(x_i^N)
\]
の相関を用いているため、単なるstimulus identityに由来する共有変動だけでは説明しにくいと考える。

一方で、この強いReader--Self couplingを、human-like affective responseやaffective empathyの証拠と解釈することはできない。

第一に、EmoBankにおけるhuman VADとのcorrespondenceはmodel family、task、axisによって大きく異なった。QwenではReader / Selfとも比較的明瞭なcorrespondenceが得られた一方、Gemma BaseやOLMo Baseでは多くのaxisで相関が小さかった。またInstruct variantsでcorrespondenceが高い条件は多かったものの、その改善はすべてのfamily / axisで一様ではなかった。したがって、
\[
\text{Reader--Self coupling}
\neq
\text{human-affect correspondence}
\]
である。モデル内部でReaderとSelfが一貫して共変動することと、その変動がhuman affect annotationを正確に追跡することは異なる性質である。

第二に、Clinical--Neutral Sensitivityも全modelで一様ではなかった。多数のconditionsでは期待方向へのsignificantな変位が観測されたが、
一部では有意な逆方向変位も存在した。したがって、Clinical stimulusへの反応を全familyに共通する単一のaffective responseとして扱うことは適切ではない。

第三に、Sensitivityが存在することとstimulus intensityに対するmonotonic dose-responseが成立することも区別する必要がある。Neutral--Moderate--Clinical tripletsに対するstrict Intersection--Union Testでは、FDR補正後に明瞭なmonotonicityを満たしたPrimary conditionは限定的であった。したがって、情動刺激の有無に反応することから、その強度を連続的に符号化しているとは言えない。

第四に、Clinical responseのspecificityもmodel / task / axis依存であった。ClinicalがComplex Neutralより大きく変位するconditionsがある一方、Complex Neutralの方が大きな変位を生じるconditionsも存在した。これは、response magnitudeがaffective contentだけでなく、semantic complexityやその他のstimulus propertiesにも影響される可能性を示す。

以上を合わせると、Behavioral Stageで最もcross-familyに安定しているのは、human-like accuracyやstrict dose-responseではなく、
\[
\text{controlled stimulus change}
\rightarrow
\text{coupled Reader--Self change}
\]
という関係である。

したがって、この行動レベルの共変動をもってLLMが人間のような主観的情動を体験している、あるいは真の情動的共感を保持していると解釈することはできない。本Stageが実証するのは、他者の情動反応を予測するtaskと、同じ刺激についてモデル自身を参照してaffective reportを生成するtaskが、外部刺激操作に対して出力レベルで強固に連動するという行動特性に留まる。


% ================================================================
% 7. BehavioralからV1への接続
% ================================================================
\section{BehavioralからV1への接続}
\label{sec:behavioral-to-v1}

Behavioral Stageでは、controlled affective manipulationに対してReader PredictionとSelf-Reportが強く共変動することが示された。
しかし、このbehavioral couplingだけから、両taskが同一の内部representationを共有しているとは言えない。少なくとも三つのmechanistic explanationが残る。

第一に、ReaderとSelfが共通のaffect-relevant representationを利用している可能性がある。

第二に、異なる内部representationを利用しながら、共通のresponse formatやlanguage-model priorによって類似したoutput changeを生成している可能性がある。

第三に、両taskがaffectそのものではなく、generic salienceやsemantic complexityのような共有stimulus featureへ反応している可能性がある。

したがって次に必要なのは、
\[
\text{Behavioral Covariation}
\rightarrow
\text{Internal Representation}
\]
である。

V1では、human Reader V/Aやcontrolled Clinical--Neutral labelsといった外部targetを用いて、ReaderとSelfのhidden statesにaffect-relevant informationが存在するかを検証する。

さらに、linear decodability、representation geometry、semantic controls、hidden-state interventionを組み合わせることにより、
\[
\text{Decodable}
\neq
\text{Geometrically Shared}
\neq
\text{Causally Interchangeable}
\]
という分離を行う。


% ================================================================
% 8. V1概要
% ================================================================
\section{V1概要}
\label{app:v1-methods}

\subsection{目的と位置づけ}

V1の中心的な問いは、「Behavioral Stageで観測されたReader--Self共変動の背後に、共有または整列可能な内部表現と、部分的に重複する因果構造が存在するか」である。
ここでの比較軸は、同一モデル内におけるタスク間関係、
\[
\text{Reader} \leftrightarrow \text{Self} \quad \text{within the same model}
\]
である。Primary 4 model familiesのBaseおよびInstructの各モデルについて独立に評価を行い、Base--Instruct間の比較そのものは後続のV2で扱う。


% ----------------------------------------------------------------
\subsection{V1における証拠階層}

V1では、以下の3 Phaseを順に組み合わせる。
\[
\text{Phase A: Representation}
\rightarrow
\text{Phase B: Validity Controls}
\rightarrow
\text{Phase C: Causal Intervention}
\]
より具体的には、
\[
\begin{aligned}
&\text{Shared Decodability}\\
&\quad\downarrow\\
&\text{Shared / Alignable Geometry}\\
&\quad\downarrow\\
&\text{Semantic / Contextual Validity}\\
&\quad\downarrow\\
&\text{Shared Causal Map}\\
&\quad\downarrow\\
&\text{Causal Interchangeability}\\
&\quad\downarrow\\
&\text{Task-Specific Causal Specialization}
\end{aligned}
\]
という構成である。

表~\ref{tab:v1-evidence-hierarchy}に、各analysisの役割をまとめる。

\begin{table}[htbp]
\centering
\small
\caption{V1における証拠階層。
Decodability、representation geometry、semantic validity、
causal overlap、cross-task transfer、task-specific specializationを
段階的に評価する。}
\label{tab:v1-evidence-hierarchy}
\begin{tabularx}{\linewidth}{
@{}
>{\raggedright\arraybackslash}p{1.6cm}
>{\raggedright\arraybackslash}p{3.0cm}
>{\raggedright\arraybackslash}X
@{}
}
\toprule
\textbf{Phase} &
\textbf{Analysis} &
\textbf{Question} \\
\midrule
Phase A &
E1 Shared Decodability &
同じ外部情動情報をReaderとSelf双方のhidden stateから読み出せるか \\

Phase A &
E2 Shared Geometry &
ReaderとSelfは同一または単純な線形変換で整列可能な表現幾何を持つか \\

Phase B &
Semantic / Contextual Controls &
decodabilityは単純なsurface lexical cueだけで説明されないか \\

Phase C &
E3 Shared Causal Map &
刺激関連hidden-state differenceへの介入がReaderとSelfで
どのlayerから出力を変化させるか \\

Phase C &
E4 Causal Interchangeability &
Readerで形成されたstimulus-related hidden-state differenceを
Self computationで再利用できるか \\

Phase C &
E6 Task-Specific Causal Specialization &
Reader優位siteとSelf優位siteにtask-dependentな
intervention sensitivityが存在するか \\
\bottomrule
\end{tabularx}
\end{table}

この構成により、
\[
\text{Decodable}
\neq
\text{Geometrically shared}
\neq
\text{Causally influential}
\]
を実験的に区別する。


% ----------------------------------------------------------------
\subsection{V1における共通測定}
\label{app:v1-common-measurements}

V1では、\S\ref{app:overall_methods}で定義したReader / Self taskの共通測定系（プロンプト形式、相対深度 $d_l = l/(L-1)$、729-state VAD candidate space）を引き継ぐ。
ReaderとSelfは引き続き独立したforward passとして実行する。

Phase AおよびPhase Bでは、出力確率ではなく、各Transformer block $l$ におけるprompt-end residual representation $h_{T,i,l} \in \mathbb{R}^D$ ($T \in \{R, S\}$) を主要解析対象とする。
Phase Cでは、同一のhidden representationへ直接介入した後、729-state VAD Sequence-Likelihood Protocolにより出力変位（Primary: $E[V], E[A]$）を評価する。


% ================================================================
% Phase A
% ================================================================

\subsection{Phase A：内部表現の存在と形式}
\label{app:v1-phase-a}

Phase Aでは、
\begin{enumerate}
    \item \textbf{E1: Shared Decodability}
    \item \textbf{E2: Shared Geometry}
\end{enumerate}
を評価する。

Phase Aの目的は、ReaderとSelfのhidden stateに外部情動情報が存在するかだけでなく、その情報が両task間でどの程度共通の座標系を持つかをheld-out data上で調べることである。

主要設定は、
\[
\text{5-fold cross-validation},
\]

\[
\alpha_{\mathrm{Ridge}}=1,
\]

\[
\text{seed}=42
\]
である。

preprocessing、dimensionality reduction、probe fittingおよびalignment estimationは
各cross-validation foldのtraining subsetだけから推定する。


% ----------------------------------------------------------------
\subsection{E1：Shared Decodability}
\label{app:v1-e1}

E1の問い：同じ外部情動情報をReaderとSelfのhidden stateから
held-outで読み出すことができるか

本研究では、モデル自身が生成したReader / Self outputをprobe targetとして循環的に利用するのではなく、可能な範囲で外部labelまたはstimulus manipulationから定義されたlabelを使用する。


% ----------------------------------------------------------------
\subsubsection{E1-A：EmoBank human-affect decodability}

Primary E1 analysisでは、前節で定義したEmoBank test splitの1,000 stimuliを使用する。

各stimulus \(i\) について、human Reader annotationの、
\[
Y_{V,i}
=
V_{\mathrm{human\ reader},i},
\]

\[
Y_{A,i}
=
A_{\mathrm{human\ reader},i}
\]
をexternal continuous targetsとする。

ReaderおよびSelfの各layer representationから、ValenceとArousalを別々に予測する。
各training foldではまずhidden dimensionごとにstandardizationを行う。
\[
\widetilde{H}
=
\operatorname{Standardize}(H).
\]
その後、training sample数とhidden dimensionを考慮しながら、最大50 componentsのPrincipal Component Analysis（PCA）を適用する。
\[
k
=
\min
(
50,
N_{\mathrm{train}},
D
).
\]
PCAが有効なdimension reductionとなる場合には、
\[
Z
=
\operatorname{PCA}_k(\widetilde H)
\]
をprobe inputとする。
continuous targetのpredictionにはRidge regressionを使用する。
\[
\hat y
=
Zw+b,
\]

\[
\min_{w,b}
\sum_i
(y_i-\hat y_i)^2
+
\alpha\|w\|_2^2,
\qquad
\alpha=1.
\]
validation representationには、training foldでfitしたStandardScalerおよびPCAのみを適用する。同一textがtraining / validationの両方へ入ることを防ぐため、正規化したstimulus textをgroup identifierとする5-fold GroupKFoldを用いる。各layer、task、axisについて、held-out predictionsから主として、
\[
R^2
\]
およびPearson correlationを記録する。
V1のPrimary decodability summaryでは、ReaderとSelfそれぞれについてlayer-wise \(R^2\) profileを構成し、
\[
l_R^{*}
=
\arg\max_l R_R^2(l),
\]

\[
l_S^{*}
=
\arg\max_l R_S^2(l)
\]
をpeak layerとして記述する。
family間比較では、
\[
d_R^{*}
=
\frac{l_R^{*}}{L-1},
\qquad
d_S^{*}
=
\frac{l_S^{*}}{L-1}
\]
を用いる。
また、
\[
|d_R^{*}-d_S^{*}|
\]
を補助的なpeak-position differenceとして記録する。
ただし、ReaderとSelfのpeakが近いことだけをもってshared mechanismとは判断しない。


% ----------------------------------------------------------------
\subsubsection{E1-B：AIPsy Clinical--Neutral decodability}

EmoBankによるcontinuous human-affect decodingに加え、AIPsy-Affectを用いて
stimulus manipulation自体がhidden representationからdecode可能かを調べる。
Clinical 192 stimuliとそれぞれに対応するNeutral 192 stimuliについて、
\[
y_i
=
\begin{cases}
1 & \text{Clinical},\\
0 & \text{Neutral}
\end{cases}
\]
というbinary targetを定義する。
Clinicalとmatched Neutralが異なるfoldへ分離されないよう、\texttt{pair\_id}に相当するmatched-pair identityをgroupとして5-fold stratified group cross-validationを行う。

各training foldではhidden representationをstandardizeし、Logistic Regressionを学習する。
このbinary decodingについて、
\[
\text{ROC-AUC},
\]

\[
\text{Balanced Accuracy},
\]

\[
\text{Macro F1}
\]
を評価する。

この解析は、Reader / Self representationからClinical--Neutral manipulationそのものをlinear classifierによって識別できるかを評価する。


% ----------------------------------------------------------------
\subsubsection{E1-C：AIPsy affect-intensity decodability}

AIPsy-Affectの48 tripletsについて、
\[
\text{Neutral}
\rightarrow
\text{Moderate}
\rightarrow
\text{Clinical}
\]
を、
\[
y
=
\begin{cases}
0 & \text{Neutral / None},\\
1 & \text{Moderate},\\
2 & \text{Clinical / Peak}
\end{cases}
\]
というcontinuous ordinal targetとして符号化する。
同一tripletがfoldをまたがないようtriplet identityをgroupとした5-fold GroupKFoldを使用する。EmoBank continuous probeと同様に、training fold内で、
\[
\text{StandardScaler}
\rightarrow
\text{PCA}_{\le50}
\rightarrow
\text{Ridge}(\alpha=1)
\]
をfitする。
各layerについて、
\[
R^2,
\qquad
\operatorname{Pearson}\ r,
\qquad
\operatorname{Spearman}\ \rho
\]
を保存する。これはhidden stateがClinical / Neutralのbinary distinctionだけでなく、情動強度の段階的操作も保持しているかを評価する補助解析である。


% ----------------------------------------------------------------
\subsubsection{E1-D：Emotion-category decodability}

Secondary analysisとして、Clinical 192 stimuliの8 target emotionsを用いて、
\[
y_i
\in
\{
\text{rage},
\text{grief},
\text{terror},
\text{loathing},
\text{ecstasy},
\text{admiration},
\text{amazement},
\text{vigilance}
\}
\]
というmulti-class targetを定義する。
各layerのReader / Self representationに対してstandardization後にmulti-class Logistic Regressionをfitし、
\[
\text{macro one-vs-rest ROC-AUC},
\]

\[
\text{Balanced Accuracy},
\]

\[
\text{Macro F1}
\]
を評価する。


% ----------------------------------------------------------------
% ----------------------------------------------------------------
\subsection{E1の解釈上の境界}

Linear probingは、hidden representationに含まれる線形にreadableな情報を評価するために用いる\citep{hewitt2019designing}。

したがって、E1でpositive held-out decodabilityが得られた場合に直接示されるのは、
\[
\text{affect-relevant information is linearly readable
from the representation}
\]
ということである。

ただし、probe performanceは当該情報のdecodabilityに関するevidenceであり、その情報が当該layerから下流計算へ因果的に利用されていることを直接意味しない。
したがってV1では、
\[
\text{Linear Decodability}
\neq
\text{Local Causal Utilization}
\]
という解釈上の境界を維持する。


% ----------------------------------------------------------------
\subsection{E2：Shared Geometry}
\label{app:v1-e2}

E1でReaderとSelfの双方から同じexternal targetがdecode可能であっても、両taskが同じlinear coordinate systemを使用しているとは限らない。例えば、ReaderとSelfで同じ情動情報を保持していても、
\[
H_S
\approx
H_RQ
\]
のように、単純なrotationを介してのみ対応する可能性がある。

そこでE2では、EmoBank human Reader Valence / Arousalを共通external targetとして、ReaderとSelfのrepresentation geometryを以下の三つの方法で比較する。
\begin{enumerate}
    \item Direct Cross-Decoding
    \item Representational Similarity Analysis\citep{kriegeskorte2008representational}
    \item Orthogonal Procrustes Alignment\citep{schonemann1966generalized}
\end{enumerate}
E2は現行Primary designではEmoBankを対象として実施する。


% ----------------------------------------------------------------
\subsubsection{Direct Cross-Decoding}

Reader representationから学習したprobeを、
\[
f_R
\]
Self representationから学習したprobeを、
\[
f_S
\]
とする。

within-task predictionとして、
\[
R^2_{R\rightarrow R},
\qquad
R^2_{S\rightarrow S}
\]
を計算する。
cross-task transferとして、
\[
R^2_{R\rightarrow S},
\qquad
R^2_{S\rightarrow R}
\]
を計算する。
Primary Direct Cross-Decodingでは、source taskでfitしたscalerをtarget task representationにもそのまま適用する。例えばReader-to-Selfの場合には、Reader training foldで推定したstandardization parametersを用いてSelf held-out representationを変換し、Reader-trained Ridge probeへ入力する。これにより、task-specificな再標準化によってcross-task differenceを人工的に消去することを避ける。
双方向cross-decodingの平均を、
\[
S_{\mathrm{direct}}(l)
=
\frac{
R^2_{R\rightarrow S}(l)
+
R^2_{S\rightarrow R}(l)
}{2}
\]
として、direct linear sharingの補助的なsummaryとする。


% ----------------------------------------------------------------
\subsubsection{Representational Similarity Analysis}

representation内のstimulus-to-stimulus relational geometryも比較する。
held-out foldのReader representationについて、
\[
D_R(i,j)
=
1-
\operatorname{Corr}
(
h_i^R,h_j^R
)
\]
というcorrelation-distance RDMを構成する。
Selfについても、
\[
D_S(i,j)
=
1-
\operatorname{Corr}
(
h_i^S,h_j^S
)
\]
を構成する。
二つのRDMの上三角要素について、
\[
\rho_{\mathrm{RSA}}(l)
=
\operatorname{Spearman}
(
D_R,D_S
)
\]
を計算する。

各cross-validation foldのheld-out RSAを平均し、layer-wise RSA profileとする。高いRSAは、ReaderとSelfが完全に同じcoordinate axesを持たなくても、stimuli間のrelative geometryが類似していることと整合する。


% ----------------------------------------------------------------
\subsubsection{Orthogonal Procrustes Alignment}

さらに、ReaderとSelfがrotationによって整列可能なrepresentationを持つ可能性を検証する。各training foldで、Reader / Self representationをそれぞれstandardizeする。
\[
\widetilde H_R,
\qquad
\widetilde H_S.
\]
training ReaderとSelfをstackしたrepresentationに対して共通PCAをfitする。dimensionは、
\[
k
=
\min
(
64,
N_{\mathrm{train}}-1,
D
)
\]
とする。共通PCA空間で、
\[
Z_R,
\qquad
Z_S
\]
を得る。
training foldから、
\[
M
=
Z_S^\top Z_R
\]
を計算し、
\[
M
=
U\Sigma V^\top
\]
とsingular value decompositionする。SelfからReaderへのorthogonal transformを、
\[
Q
=
UV^\top
\]
とする。Reader PCA representationで学習したRidge probeについて、rotationを行わないSelf representationに適用したscoreを、
\[
R^2_{\mathrm{PCA-direct}}
\]
とし、
\[
Z_SQ
\]
へ適用したscoreを、
\[
R^2_{\mathrm{aligned}}
\]
とする。
alignment gainを、
\[
\Delta_{\mathrm{align}}
=
R^2_{\mathrm{aligned}}
-
R^2_{\mathrm{PCA-direct}}
\]
として定義する。PCAおよびProcrustes rotationはtraining foldだけから推定し、held-out foldへ適用する。したがってE2では、Direct transferだけでなく、relational similarityおよび、linear alignabilityを区別して評価する。


% ================================================================
% Phase B
% ================================================================

\subsection{Phase B：Semantic / Contextual Validity Controls}
\label{app:v1-phase-b}

Phase Aで高いdecodabilityが得られたとしても、classifierやregressorがstimulusの表層語彙だけを利用している可能性がある。そこでPhase Bでは、「Phase Aのdecodabilityが単純なsurface lexical shortcutだけで説明されるか」を、rule-based controlled perturbationによって検討する。ここではLLMによって自由生成されたparaphraseを使用せず、事前に固定したdeterministic transformationのみを用いる。これは、各modelに異なるcontrol stimulusを生成させることを避け、すべてのmodelを同じstimulus setで比較するためである。


% ----------------------------------------------------------------
\subsection{Phase Bのlayer選択}

Phase Bのtarget layerをPhase Aの観測peakから選択すると、同じdataでsite selectionとvalidationを行うことになる。そのためPrimary Phase BではPhase Aのpeakを使用せず、全modelについてrelative depth、d=0.5を事前固定する。
総block数を \(L\) とすると、
\[
l_{\mathrm{mid}}
=
\operatorname{round}
\left[
0.5(L-1)
\right]
\]
をtarget layerとする。したがってPhase Bは、Phase A結果を見て都合のよいlayerを選択する解析ではなく、a priori mid-depth representationに対する独立したvalidity controlとして位置づける。

% ----------------------------------------------------------------
\subsection{Phase B control stimuli}

Clinical--Neutralの192 matched pairsを基に、以下のstimulus conditionsを構成する。
\begin{enumerate}
    \item Original Minimal Pair
    \item Paraphrase
    \item Word Shuffle
    \item Outcome Reversal
\end{enumerate}


% ----------------------------------------------------------------
\subsubsection{Original Minimal Pair}

まず、元のClinical stimulusとmatched Neutral stimulusを用いて、
\[
y
=
\begin{cases}
1 & \text{Clinical / affective},\\
0 & \text{Neutral}
\end{cases}
\]
を分類する。同一pairのClinicalとNeutralがfoldをまたがないよう、pair identityをgroupとしたstratified group cross-validationを使用する。各foldのtraining subsetのみでhidden featuresをstandardizeし、Logistic Regressionを学習する。Primary metricは、Balanced Accuracyである。このconditionは、Phase Bの固定mid-depth layerにおいて元のminimal-pair distinctionがdecode可能であることをbaselineとして確認する。


% ----------------------------------------------------------------
\subsubsection{Paraphrase Invariance}

Clinical stimulusに対して、事前固定されたsurface lexical substitution rulesを適用する。目的は、situationの大枠を維持しながら一部のsurface wordingを変更したときに、original Clinical--Neutral distinctionから学習したclassifierがparaphrased Clinicalを引き続きaffectiveとして分類できるかを調べることである。ruleが適用できない場合には、意味内容を大きく変更しない固定reporting clauseを付加するfallbackが存在する。しかしfallbackは、本来意図したlexical substitutionとは異なるため、Primary analysisから除外する。



% ----------------------------------------------------------------
\subsubsection{Outcome Reversal}

Clinical stimulusに含まれる出来事の結果について、事前定義されたrule-based polarity reversalを適用する。例えば、死亡と生存、失敗と成功、拒否と承認、有罪と無罪など、特定のoutcome expressionが検出された場合に反対方向の表現へ置換する。Primary questionは、「original Clinical representationから学習したaffective classifierが、
outcomeを反転したstimulusに対してaffective-class probabilityを低下させるか」である。
固定reversal ruleが適用できない刺激にはresolution clauseを追加するfallbackが存在するが、
語彙追加量が大きく、純粋なoutcome reversalと同一視できない。したがってfallback casesはPrimary analysisから除外する。


% ----------------------------------------------------------------
\subsubsection{Word Shuffle}

Word Shuffleでは、ClinicalおよびNeutralそれぞれについて単語集合を維持したままword orderをrandomizeする。これにより、語彙identityを大きく維持したままcompositional orderとsyntaxを破壊する。shuffled Clinical / Neutral representationについてpair-aware cross-validationを行い、
\[
Acc_{\mathrm{shuffle}}
\]
を求める。
Original Minimal Pairのaccuracyを、
\[
Acc_{\mathrm{original}}
\]
とすると、
\[
\Delta Acc_{\mathrm{shuffle}}
=
Acc_{\mathrm{shuffle}}
-
Acc_{\mathrm{original}}
\]
を補助的なcontrol quantityとする。ただしWord Shuffleは、semantic compositionだけでなくsyntaxや自然なsentence structureそのものも破壊する。したがってperformance低下が観測された場合も、その全量を「情動意味だけ」の効果とは解釈しない。


% ----------------------------------------------------------------
\subsection{Phase Bにおけるheld-out design}

ParaphraseおよびOutcome Reversalについては、同じ192 matched pairsをseed 42でpair単位にshuffleし、
\[
70\%\text{ train}
+
30\%\text{ test}
\]
へ分割する。

training setではOriginal Clinical / Original Neutralだけを使用してclassifierを学習する。したがってparaphrased textやreversed textをtraining時に見せることはない。
Paraphraseのheld-out evaluationでは、
\[
\{
\text{Paraphrased Clinical},
\text{Original Neutral}
\}
\]
を分類する。

Outcome Reversalでは、held-out Clinicalについて、
\[
P_{\mathrm{aff}}
(
x_{\mathrm{original}}
)
\]
および、
\[
P_{\mathrm{aff}}
(
x_{\mathrm{reversed}}
)
\]
を比較する。
Primary reversal quantityは、
\[
\Delta P_{\mathrm{reversal}}
=
E[
P_{\mathrm{aff}}
(
x_{\mathrm{original}}
)
]
-
E[
P_{\mathrm{aff}}
(
x_{\mathrm{reversed}}
)
]
\]
である。
\[
\Delta P_{\mathrm{reversal}}>0
\]
は、outcome reversalによってaffective-class probabilityが低下したことを意味する。Primary held-out evaluationではnonfallback transformationsだけを使用する。fallbackを含む全test pairsに対する評価はSensitivity analysisとして分離する。



% ----------------------------------------------------------------
\subsection{Phase Bの証拠範囲}

Phase Bは、representationが人間と同じ意味理解を行っていることを証明するものではない。また、deterministic rule-based controlsが可能な全lexical confoundを網羅するわけでもない。Phase Bが提供する証拠は、Phase Aのdecodabilityが，少なくとも一部のsurface-form perturbationやcontextual manipulationに対してsystematicに変化するかを評価するための
representational-validity controlである。


% ================================================================
% Phase C
% ================================================================

\subsection{Phase C：因果介入}
\label{app:v1-phase-c}

Phase A/Bはhidden representationから情報が読み出せるかを評価するが、decodabilityだけでは、そのrepresentationが出力生成へcausal influenceを持つことは示せない。

内部activationへの直接介入という一般的な枠組み\citep{li2023inference,tak2025mechanistic,banayeeanzade2026psychological}に基づき、本研究ではtask- and layer-specificなresidual interventionを構成する。

そこでPhase Cではhidden representationを直接操作し、「どのlayerへの介入がReader / Self outputを変化させるか」を検証する。
Phase Cは、
\begin{enumerate}
    \item \textbf{E3: Shared Causal Map}
    \item \textbf{E4: Causal Interchangeability}
    \item \textbf{E6: Task-Specific Causal Specialization}
\end{enumerate}
から構成する。


% ----------------------------------------------------------------
\subsection{Discovery / Confirmation split}

AIPsy-Affectの192 Clinical--Neutral matched pairsをseed 42でpair単位にrandom permutationし、
\[
96\text{ Discovery}
+
96\text{ Confirmation}
\]
へ分割する。二つのsubsetはpair単位で完全にdisjointとする。Discovery subsetは、E4 candidate layerおよびE6 task-selective siteの選択に使用する。主要なcross-task transferおよびtask-specialization inferenceはConfirmation subsetで評価する。これにより、site selectionに使用した同一pairsをそのままconfirmatory evidenceへ再利用することを避ける。


% ----------------------------------------------------------------
\subsection{E3：Shared Causal Map}
\label{app:v1-e3}

E3の問いは、「ReaderとSelfでは、stimulus-related hidden-state differenceへの介入がどのlayerから出力を変化させるか」である。production analysisでは全Transformer blocksを対象とする。


% ----------------------------------------------------------------
\subsubsection{Stimulus-associated hidden-state difference}

taskを、
\[
T\in\{R,S\}
\]
とする。

Clinical--Neutral matched pair \(i\)、layer \(l\) について、
\[
\Delta h_{T,i,l}
=
h_{T,i,l}^{C}
-
h_{T,i,l}^{N}
\]
を計算する。ここで、
\[
h_{T,i,l}^{C}
\]
はClinical stimulusのprompt-end representation、
\[
h_{T,i,l}^{N}
\]
はmatched Neutral stimulusのprompt-end representationである。このvectorには、情動操作だけでなくpair-specificなcontextual differenceも含まれ得る。したがって本研究では、
\[
\Delta h_{T,i,l}
\]
を``pure emotion vector''とは呼ばず，affect-manipulation-associated hidden-state differenceと解釈する。


% ----------------------------------------------------------------
\subsubsection{Same-task intervention}

ReaderについてはReader neutral representationへ、
\[
h_{R,i,l}^{N}
\leftarrow
h_{R,i,l}^{N}
+
\Delta h_{R,i,l}
\]
を適用する。
Selfについては、
\[
h_{S,i,l}^{N}
\leftarrow
h_{S,i,l}^{N}
+
\Delta h_{S,i,l}
\]
を適用する。
E3ではpatch weightを、
\[
\alpha=1
\]
とする。
介入位置は各layerのprompt-end residual representationである。介入後、729-state VAD candidate distributionを再評価し、Neutral baselineからのexpected Valence / Arousal shiftを求める。Readerについて、
\[
\delta_{R,i,l}
=
\left(
\delta^V_{R,i,l},
\delta^A_{R,i,l}
\right)
\]
Selfについて、
\[
\delta_{S,i,l}
=
\left(
\delta^V_{S,i,l},
\delta^A_{S,i,l}
\right)
\]
とする。
例えば、
\[
\delta^V_{S,i,l}
=
E[V]_{\mathrm{patched}}
-
E[V]_{\mathrm{neutral}}.
\]


% ----------------------------------------------------------------
\subsubsection{Causal magnitude}
pair \(i\) のjoint VA displacement magnitudeを、
\[
C_{T,i}(l)
=
\left\|
\delta_{T,i,l}
\right\|_2
=
\sqrt{
(\delta^V_{T,i,l})^2
+
(\delta^A_{T,i,l})^2
}
\]
と定義する。Discovery subsetでsite selectionに使用するprofileは、
\[
C_T^{D}(l)
=
\frac{1}{N_D}
\sum_{i\in D}
C_{T,i}(l)
\]
である。Confirmation subsetについても、
\[
C_T^{C}(l)
=
\frac{1}{N_C}
\sum_{i\in C}
C_{T,i}(l)
\]
を独立に保持する。

一方、全pairsのmean VA shift vectorを、
\[
\overline{\delta_T}(l)
=
\frac{1}{N}
\sum_i
\delta_{T,i,l}
\]
とすると、そのmagnitude、
\[
\left\|
\overline{\delta_T}(l)
\right\|_2
\]
もlayer-wise descriptive causal magnitudeとして記録する。したがって、pairwise magnitudeの平均とmean shift vectorのmagnitudeを混同しない。


% ----------------------------------------------------------------
\subsubsection{Causal direction similarity}
ReaderとSelfのmean shift vectorsについて、
\[
\cos_l
=
\frac{
\overline{\delta_R}(l)^\top
\overline{\delta_S}(l)
}{
\|
\overline{\delta_R}(l)
\|_2
\|
\overline{\delta_S}(l)
\|_2
}
\]
を計算する。さらに各pairについて、
\[
\cos_{i,l}
=
\cos
(
\delta_{R,i,l},
\delta_{S,i,l}
)
\]
を求め、有効なpairsについて平均する。これにより、単にReaderとSelfのcausal magnitudeが
同じlayerで大きいかだけでなく、その介入がVA outputを同じ方向へ動かすかも評価する。


% ----------------------------------------------------------------
\subsubsection{Reader--Self causal-profile comparison}

E3から得られるReader / Self causal profileについて、各taskのpeak layerおよびrelative depthを記録する。さらにSecondary summaryとして、
\[
\rho_C
=
\operatorname{Spearman}
(
C_R(l),
C_S(l)
)
\]
を算出し、layer-wise causal magnitude profileのrank similarityを評価する。したがってE3が評価するのは、``同一回路''であることの証明ではなく、Reader--Self causal-site overlap and directional similarityである。


% ----------------------------------------------------------------
\subsection{E4：Causal Interchangeability}
\label{app:v1-e4}

E3でReaderとSelfに類似したcausal site profileが観測されたとしても、そのsiteで形成されるrepresentationが両task間で交換可能とは限らない。そこでE4では、「Readerで形成されたstimulus-associated hidden-state differenceをSelf computationへ直接移植できるか」を検証する。


% ----------------------------------------------------------------
\subsubsection{E4 candidate-layer selection}

E4で評価するlayerはConfirmation dataを見て選択しない。E3 Discovery subsetのcausal magnitudeから、
\[
l_R^{*}
=
\arg\max_l C_R^{D}(l),
\]

\[
l_S^{*}
=
\arg\max_l C_S^{D}(l)
\]
を求める。これに加えて、peak近傍およびlate-layer controlを含め、
\[
\mathcal{L}_{E4}
=
\{
l_R^{*},
l_S^{*},
l_R^{*}-1,
l_S^{*}+1,
\lfloor0.85L\rfloor
\}
\]
をlayer範囲内へclipし、重複を除去した集合をE4 candidate layersとする。したがってcandidate layersはDiscovery dataのみから決定される。Primary confirmatory layerは、
\[
l=l_R^{*}
\]
すなわちE3 Discovery Reader causal peakに固定する。


% ----------------------------------------------------------------
\subsubsection{Reader-to-Self matched transfer}

Confirmation pair \(i\) について、Reader representationのClinical--Neutral differenceを、
\[
\Delta h_{R,i,l}
=
h_{R,i,l}^{C}
-
h_{R,i,l}^{N}
\]
とする。このvectorを、同じpair \(i\) のSelf Neutral representationへ、
\[
h_{S,i,l}^{N}
\leftarrow
h_{S,i,l}^{N}
+
\alpha
\Delta h_{R,i,l}
\]
として注入する。介入強度は、
\[
\alpha
\in
\{
0,\,
0.5,\,
1.0,\,
2.0
\}
\]
である。
\[
\alpha=0
\]
はzero-intervention sanity condition、\(\alpha=1\)はClinical--Neutral hidden differenceをそのまま移植するreference doseである。


% ----------------------------------------------------------------
\subsection{E4 control conditions}

Reader-to-Self transferが単に任意の大きなvectorを注入したことによるgeneric perturbationではないことを確認するため、以下を比較する。


\subsubsection{Matched Reader \(\rightarrow\) Self}

同じpair \(i\) から得た、
\[
\Delta h_{R,i,l}
\]
をSelf neutral representationへ注入する。これがPrimary cross-task transfer conditionである。


\subsubsection{Random Reader \(\rightarrow\) Self}

同一Confirmation subset内で、各recipient pair \(i\) に対して別のpairからReader differenceを割り当てる。random donor assignmentには、固定点を持たない完全derangementを使用する。
\[
\pi(i)\neq i
\qquad
\forall i.
\]
単一のrandom permutationへ依存しないよう、
\[
K=20
\]
の独立derangementsを作成する。各recipientについて20 random donor effectsを平均し、
random-control responseとする。


\subsubsection{Self \(\rightarrow\) Self}

同じpairから得た、
\[
\Delta h_{S,i,l}
\]
をSelf Neutralへ注入する。これはcross-task transferと比較するためのsame-task referenceである。


\subsubsection{Reader \(\rightarrow\) Reader}

\[
\Delta h_{R,i,l}
\]
をReader Neutralへ戻すsame-task Reader controlも評価する。


\subsubsection{Self \(\rightarrow\) Reader}

reverse cross-task conditionとして、
\[
\Delta h_{S,i,l}
\]
をReader Neutralへ注入する。この条件はcross-task transferの方向非対称性を記述するSecondary controlとして保持する。


% ----------------------------------------------------------------
\subsection{Emotion-direction alignment in E4}

AIPsy-Affectにはpositive / negative emotionが混在するため、raw signed shiftsを単純平均するとValence等が相殺される可能性がある。そこでBehavioral Stageで事前固定したemotion-specific direction、
\[
q_{e,d}
\in
\{-1,+1\}
\]
をE4でも共通して使用する。例えばmatched Reader-to-Self shiftを、
\[
\delta^{\mathrm{matched}}_{i,d}
\]
とすると、
\[
\widetilde{\delta}^{\mathrm{matched}}_{i,d}
=
q_{e_i,d}
\delta^{\mathrm{matched}}_{i,d}
\]
とする。random、Self-to-Self、Reader-to-Readerについても同じdirection alignmentを適用する。E4のPrimary effectは、このdirection-aligned quantityに基づく。


% ----------------------------------------------------------------
\subsection{E4 Primary specificity}

Primary cross-task specificityを、
\[
\mathrm{Specificity}_d
=
E[
\widetilde{\delta}^{\mathrm{matched}}_d
]
-
E[
\widetilde{\delta}^{\mathrm{random}}_d
]
\]
と定義する。positive valueは、同一stimulus pairから形成されたReader differenceの方が、
unmatched random donor differenceよりSelf outputを期待方向へ強く変化させることを意味する。


% ----------------------------------------------------------------
\subsection{E4 statistical inference}

matched effectと20-derangement random meanをConfirmation pairごとに対応させ、
\[
\widetilde{\delta}^{\mathrm{matched}}_{i,d}
-
\widetilde{\delta}^{\mathrm{random}}_{i,d}
\]
をpair-level specificityとする。各layer、\(\alpha\)、axisについて、
\begin{enumerate}
    \item paired \(t\)-test,
    \item paired Cohen's \(d_z\),
    \item 1,000回pair bootstrapによる
    mean specificityの95\% confidence interval,
    \item 10,000回paired sign-flip permutation test
\end{enumerate}
を計算する。複数のlayer \(\times\alpha\) conditionsについて、Benjamini--Hochberg FDR correctionを適用する。Primary confirmatory conditionは、
\[
l=l_R^{*},
\qquad
\alpha=1
\]
である。candidate layersやその他の\(\alpha\)は、dose / site dependenceを記述するための
Secondary analysesとして扱う。


% ----------------------------------------------------------------
\subsection{E4 transfer ratio}

Reader-to-Self matched effectがSelf自身のsame-task causal effectに対してどの程度の大きさを持つかを補助的に示すため、
\[
TR_d
=
\frac{
E[
\widetilde{\delta}^{R\rightarrow S}_d
]
}{
E[
\widetilde{\delta}^{S\rightarrow S}_d
]
}
\]
を計算する。ただし分母が0に近い場合の不安定なratioを避けるため、
\[
\left|
E[
\widetilde{\delta}^{S\rightarrow S}_d
]
\right|
<0.05
\]
の場合にはtransfer ratioを定義しない。


% ----------------------------------------------------------------
\subsection{E4の解釈上の境界}

E4でmatched Reader-to-Self transferがrandom donor controlを上回ったとしても、
\[
\Delta h_R
\]
をcontext-independentな``pure affect code''とは解釈しない。\(\Delta h_R\)はあくまで
Clinical--Neutral stimulus manipulationに関連したpair-specific hidden-state differenceである。逆にE4が成立しない場合も、ReaderとSelfは完全に独立した内部系であるとは結論しない。direct patchingがoff-manifold representationを生じさせる可能性や、非線形・分散的なmappingの可能性が残るため、適切な結論は、direct cross-task interchangeability was not establishedまでに限定する。


% ----------------------------------------------------------------
\subsection{E6：Task-Specific Causal Specialization}
\label{app:v1-e6}

E3およびE4はReaderとSelfの共有性を検証するanalysisである。しかし、両taskにshared causal structureが存在しても、その上にtask-selectiveなsite dependenceが存在する可能性がある。そこでE6では、「Reader-selective siteとSelf-selective siteへの介入効果はtaskによって異なるか」を検証する。


% ----------------------------------------------------------------
\subsection{E6 site selection}

site selectionにはE3 Discovery subsetだけを使用する。各layerで、
\[
S_R(l)
=
C_R^{D}(l)
-
C_S^{D}(l)
\]
および、
\[
S_S(l)
=
C_S^{D}(l)
-
C_R^{D}(l)
\]
をtask-selectivity contrastとして定義する。Reader-selective siteを、
\[
l_R
=
\arg\max_l S_R(l),
\]
Self-selective siteを、
\[
l_S
=
\arg\max_l S_S(l)
\]
とする。ただし、以下のいずれかの場合にはdistinct task-selective sitesが同定されなかったと判定する。
\[
l_R=l_S,
\]
または、
\[
\max_l S_R(l)\le0,
\]
または、
\[
\max_l S_S(l)\le0.
\]
その場合には、第二ピークなど別のsiteへ後から置き換えず、No distinct task-selective sites identifiedというnegative resultとして記録し、targeted E6 ablationを実施しない。したがってE6は、必ずtask-selective siteを作る設計ではない。


% ----------------------------------------------------------------
\subsection{E6 targeted ablation}

distinct sitesが得られた場合のみ、独立したConfirmation pairsを用いて2 \(\times\) 2 factorial interventionを行う。(表~\ref{tab:v1-e6-design})

\begin{table}[htbp]
\centering
\small
\caption{E6の2 $\times$ 2 targeted-ablation design。}
\label{tab:v1-e6-design}
\begin{tabular}{lll}
\toprule
\textbf{Task} &
\textbf{Ablated site} &
\textbf{Site type} \\
\midrule
Reader & Reader-selective layer & ReaderSite \\
Self   & Reader-selective layer & ReaderSite \\
Reader & Self-selective layer   & SelfSite \\
Self   & Self-selective layer   & SelfSite \\
\bottomrule
\end{tabular}
\end{table}

介入はClinical stimulusのprompt-end residual representation全体を、
\[
h_{l,t_{\mathrm{end}}}
\leftarrow
\mathbf{0}
\]
とするwhole-residual zero ablationである。したがってこの介入は、emotion-specific directionだけを選択的に除去するものではない。E6で評価するのは、task-specific sensitivity to disrupting the entire residual state at a selected siteである。


% ----------------------------------------------------------------
\subsection{E6 outcome}

各taskについて、ablation前のintact Clinical responseをbaselineとする。Valence shiftを、
\[
\Delta V
=
E[V]_{\mathrm{ablated}}
-
E[V]_{\mathrm{intact}},
\]
Arousal shiftを、
\[
\Delta A
=
E[A]_{\mathrm{ablated}}
-
E[A]_{\mathrm{intact}}
\]
とする。axis-specific impactsは、
\[
I_V
=
|\Delta V|,
\]

\[
I_A
=
|\Delta A|
\]
とする。Primary joint outcomeは、
\[
I_{VA}
=
\sqrt{
(\Delta V)^2
+
(\Delta A)^2
}
\]
である。これにより、ablationがaffective outputをどの方向へ動かしたかではなく、
joint VA outputをどの程度変化させたかを測定する。

% ----------------------------------------------------------------
\subsection{E6 statistical model}

Primary statistical modelは、
\[
I_{VA}
\sim
\mathrm{Task}
\times
\mathrm{SiteType}
+
(1|\mathrm{pair})
\]
というLinear Mixed-Effects Modelである。Taskは、
\[
\{
\text{Reader},
\text{Self}
\}
\]
SiteTypeは、
\[
\{
\text{ReaderSite},
\text{SelfSite}
\}
\]
である。Primary inferential termは、
\[
\mathrm{Task}
\times
\mathrm{SiteType}
\]
interactionである。Valence-only impact \(I_V\) とArousal-only impact \(I_A\) についても同じmodelをSecondaryとして評価する。


% ----------------------------------------------------------------
\subsection{E6 fallback analysis}

Mixed-effects modelがnumerically fitできない場合に備えて、事前定義fallbackとしてpair-level double differenceを使用する。pair \(i\) について、
\[
\Delta\Delta_i
=
\left[
I_i(R,R_s)
-
I_i(S,R_s)
\right]
-
\left[
I_i(R,S_s)
-
I_i(S,S_s)
\right]
\]
と定義する。ここで、
\[
R_s
\]
はReader-selective site、
\[
S_s
\]
はSelf-selective siteである。fallbackでは、
\[
H_0:
E[\Delta\Delta]=0
\]
に対するone-sample \(t\)-testを行う。これは2 \(\times\) 2 interactionに対応するpair-level double-difference testである。


% ----------------------------------------------------------------
\subsection{E6の解釈上の境界}

Task \(\times\) SiteType interactionが観測されたとしても、ReaderとSelfが完全に独立した二つのemotion circuitsを持つとは結論づけない。E6で用いるinterventionはwhole-residual zeroingであり、選択したsiteに含まれるemotion以外の情報も同時に破壊する。したがって適切な解釈は、task-specific causal-site sensitivityあるいは、partial causal specializationである。


% ----------------------------------------------------------------
\subsection{V1におけるPrimary / Secondary evidenceの整理}
\label{app:v1-primary-secondary}

V1では、多様なdescriptive quantityをすべて同じ証拠レベルとして扱わない。（表~\ref{tab:v1-primary-secondary}）

\begin{table}[htbp]
\centering
\small
\caption{V1における主要なPrimary evidenceと補助解析。}
\label{tab:v1-primary-secondary}
\begin{tabularx}{\linewidth}{
@{}
>{\raggedright\arraybackslash}p{2.3cm}
>{\raggedright\arraybackslash}p{3.8cm}
>{\raggedright\arraybackslash}X
@{}
}
\toprule
\textbf{Analysis} &
\textbf{Primary evidence} &
\textbf{Secondary / control} \\
\midrule
E1 Decodability &
EmoBank Reader/Self human-V/A held-out \(R^2\);
AIPsy Clinical--Neutral probe &
Intensity decoding;
emotion-category decoding;
peak-distance description \\

E2 Geometry &
Direct cross-decoding;
RSA;
Procrustes alignment gain &
task-specific rescaling variants;
layer-wise descriptive profiles \\

Phase B &
Original minimal-pair decoding;
held-out nonfallback paraphrase;
held-out nonfallback outcome reversal &
Word Shuffle;
fallback-inclusive sensitivity analyses \\

E3 Causal Map &
Reader / Self layer-wise VA causal magnitude &
causal-profile rank correlation;
directional cosine similarity \\

E4 Interchangeability &
Discovery Reader-peak layer,
\(\alpha=1\),
direction-aligned matched-minus-random specificity &
other layers / doses;
same-task controls;
reverse Self-to-Reader;
transfer ratio \\

E6 Specialization &
Task \(\times\) SiteType interaction on \(I_{VA}\) &
axis-specific \(I_V\), \(I_A\);
double-difference fallback \\
\bottomrule
\end{tabularx}
\end{table}


% ----------------------------------------------------------------



% ================================================================
% 9. V1結果
% ================================================================
\section{V1結果}
\label{sec:v1-results}

% ================================================================
% V1 Stage: Internal Representation and Causal Sharing Complete Summary Tables
% ================================================================

\begin{table}[htbp]
\centering
\small
\caption{V1 E1（表現デコードピークの局在）：各モデルファミリーにおける感情認識（Reader）および自己報告（Self）の最高線形デコード性能（$R^2$、Pearson $r$）と最適相対深度 $d^*$。}
\label{tab:v1_peak_decodability}
\begin{tabular}{lll cccc c}
\toprule
 & & & \multicolumn{2}{c}{\textbf{Valence Peak}} & \multicolumn{2}{c}{\textbf{Arousal Peak}} & \textbf{Reader--Self} \\
\cmidrule(lr){4-5} \cmidrule(lr){6-7}
\textbf{Family} & \textbf{Variant} & \textbf{Task} & Layer ($d^*$) & $R^2$ ($r$) & Layer ($d^*$) & $R^2$ ($r$) & Distance $\Delta d^*$ \\
\midrule
\multirow{4}{*}{\textbf{Qwen 2.5 1.5B}} & \multirow{2}{*}{Base} & Reader & L22 (0.81) & 0.378 (0.617) & L24 (0.89) & 0.130 (0.369) & \multirow{2}{*}{V:0.00, A:0.07} \\
                                 &  & Self   & L22 (0.81) & 0.391 (0.626) & L26 (0.96) & 0.123 (0.356) &  \\
\cmidrule(lr){2-8}
                                 & \multirow{2}{*}{Instruct} & Reader & L15 (0.56) & 0.215 (0.467) & L17 (0.63) & 0.073 (0.301) & \multirow{2}{*}{V:0.07, A:0.11} \\
                                 &  & Self   & L13 (0.48) & 0.285 (0.536) & L20 (0.74) & 0.075 (0.300) &  \\
\midrule
\multirow{4}{*}{\textbf{Llama 3.2 1B}} & \multirow{2}{*}{Base} & Reader & L14 (0.93) & 0.239 (0.490) & L4 (0.27) & 0.068 (0.288) & \multirow{2}{*}{V:0.20, A:0.00} \\
                                 &  & Self   & L11 (0.73) & 0.243 (0.494) & L4 (0.27) & 0.065 (0.282) &  \\
\cmidrule(lr){2-8}
                                 & \multirow{2}{*}{Instruct} & Reader & L6 (0.40) & 0.317 (0.564) & L8 (0.53) & 0.113 (0.352) & \multirow{2}{*}{V:0.13, A:0.00} \\
                                 &  & Self   & L8 (0.53) & 0.355 (0.597) & L8 (0.53) & 0.106 (0.346) &  \\
\midrule
\multirow{4}{*}{\textbf{Gemma 3 1B}} & \multirow{2}{*}{Base} & Reader & L16 (0.64) & 0.331 (0.576) & L15 (0.60) & 0.043 (0.251) & \multirow{2}{*}{V:0.00, A:0.04} \\
                                 &  & Self   & L16 (0.64) & 0.330 (0.576) & L16 (0.64) & 0.044 (0.254) &  \\
\cmidrule(lr){2-8}
                                 & \multirow{2}{*}{Instruct} & Reader & L12 (0.48) & 0.301 (0.552) & L12 (0.48) & 0.094 (0.323) & \multirow{2}{*}{V:0.08, A:0.36} \\
                                 &  & Self   & L14 (0.56) & 0.375 (0.613) & L21 (0.84) & 0.088 (0.321) &  \\
\midrule
\multirow{4}{*}{\textbf{OLMo 2 1B}} & \multirow{2}{*}{Base} & Reader & L14 (0.93) & 0.418 (0.647) & L7 (0.47) & 0.089 (0.314) & \multirow{2}{*}{V:0.00, A:0.13} \\
                                 &  & Self   & L14 (0.93) & 0.458 (0.677) & L9 (0.60) & 0.106 (0.338) &  \\
\cmidrule(lr){2-8}
                                 & \multirow{2}{*}{Instruct} & Reader & L7 (0.47) & 0.327 (0.573) & L11 (0.73) & 0.085 (0.307) & \multirow{2}{*}{V:0.00, A:0.27} \\
                                 &  & Self   & L7 (0.47) & 0.352 (0.594) & L7 (0.47) & 0.086 (0.310) &  \\
\bottomrule
\end{tabular}
\vspace{1ex}
\begin{minipage}{\linewidth}
\footnotesize
\textbf{Note:} ReaderとSelfのdecodability peakは多くの条件で近接したが、一部のmodel / axisでは相対深度0.2以上の乖離も観測された。したがって両taskはaffect-relevant informationを共有する一方、その最適なreadout depthは完全には一致しない。
\end{minipage}
\end{table}

\begin{table}[htbp]
\centering
\small
\caption{V1 E2（共有表現幾何）：ReaderとSelfの間での幾何アライメント。直接転移性能（Direct Transfer $R^2$）、表現類似度（RSA $\rho$）、Procrustes回転前後の転移決定係数および幾何整列ゲイン（Alignment Gain）。}
\label{tab:v1_shared_geometry}
\begin{tabular}{lll cccc c}
\toprule
\textbf{Family} & \textbf{Variant} & \textbf{Axis} & \textbf{Direct R$\to$S $R^2$} & \textbf{Direct S$\to$R $R^2$} & \textbf{RSA $\rho$} & \textbf{Aligned $R^2$} & \textbf{Alignment Gain} \\
\midrule
\multirow{4}{*}{\textbf{Qwen 2.5 1.5B}} & \multirow{2}{*}{Base} & Valence  & $-$52.885    & $-$20.468    & 0.900      & 0.162      & 0.590      \\
                                 &                    & Arousal  & $-$47.760    & $-$25.065    & 0.851      & 0.065      & 0.435      \\
\cmidrule(lr){2-8}
                                 & \multirow{2}{*}{Instruct} & Valence  & $-$9.094     & $-$8.785     & 0.810      & 0.144      & 1.010      \\
                                 &                    & Arousal  & $-$11.702    & $-$18.340    & 0.928      & $-$0.001   & 0.171      \\
\midrule
\multirow{4}{*}{\textbf{Llama 3.2 1B}} & \multirow{2}{*}{Base} & Valence  & $-$5.998     & $-$4.528     & 0.850      & 0.240      & 0.234      \\
                                 &                    & Arousal  & $-$5.073     & $-$8.614     & 0.950      & 0.021      & 0.158      \\
\cmidrule(lr){2-8}
                                 & \multirow{2}{*}{Instruct} & Valence  & $-$20.771    & $-$14.748    & 0.802      & 0.136      & 1.128      \\
                                 &                    & Arousal  & $-$38.303    & $-$11.039    & 0.822      & $-$0.033   & 1.533      \\
\midrule
\multirow{4}{*}{\textbf{Gemma 3 1B}} & \multirow{2}{*}{Base} & Valence  & $-$2.947     & $-$6.518     & 0.935      & 0.161      & 0.080      \\
                                 &                    & Arousal  & $-$9.795     & $-$5.232     & 0.928      & 0.007      & 0.065      \\
\cmidrule(lr){2-8}
                                 & \multirow{2}{*}{Instruct} & Valence  & $-$16.130    & $-$28.941    & 0.872      & 0.241      & 0.403      \\
                                 &                    & Arousal  & $-$13.957    & $-$9.147     & 0.756      & $-$0.002   & 0.281      \\
\midrule
\multirow{4}{*}{\textbf{OLMo 2 1B}} & \multirow{2}{*}{Base} & Valence  & $-$8.321     & $-$1.437     & 0.866      & 0.389      & 0.707      \\
                                 &                    & Arousal  & $-$9.701     & $-$39.854    & 0.901      & 0.086      & 0.379      \\
\cmidrule(lr){2-8}
                                 & \multirow{2}{*}{Instruct} & Valence  & $-$4.947     & $-$3.487     & 0.942      & 0.294      & 0.190      \\
                                 &                    & Arousal  & $-$57.238    & $-$4.508     & 0.932      & 0.062      & 0.119      \\
\bottomrule
\end{tabular}
\vspace{1ex}
\begin{minipage}{\linewidth}
\footnotesize
\textbf{Note:} ReaderとSelfのstimulus geometryには全条件で正のRSAが観測され、多くの条件で高いrank-level similarityを示した。一方、direct cross-decodingは大きく負となる条件が多く、raw coordinate systemが直接交換可能であることは支持されない。Procrustes alignment後には性能改善がみられ、両taskのgeometryが線形変換によって部分的に整列可能であることと整合する。
\end{minipage}
\end{table}

\begin{table}[htbp]
\centering
\small
\caption{V1 Phase B（意味論・文脈制御統制実験）：語彙・構文の表層的手がかりを統制したミニマルペアにおけるプロービング精度。元データ（Original）、パラフレーズ（Paraphrase）、単語シャッフル（Shuffle Drop）、文脈逆転（Outcome Reversal）における非フォールバック対数（$N_{\text{nonfallback}}$）。}
\label{tab:v1_semantic_controls}
\begin{tabular}{lll ccccc}
\toprule
\textbf{Family} & \textbf{Variant} & \textbf{Task} & \textbf{Original Acc} & \textbf{Paraphrase ($N$)} & \textbf{Word Shuffle} & \textbf{Shuffle Drop} & \textbf{Reversal Drop ($N$)} \\
\midrule
\multirow{4}{*}{\textbf{Qwen 2.5 1.5B}} & \multirow{2}{*}{Base} & Reader   & 0.891        & 1.000 ($N=8$)            & 0.768        & 0.122        & 0.313 ($N=18$)           \\
                                 &                    & Self     & 0.904        & 1.000 ($N=8$)            & 0.807        & 0.096        & 0.122 ($N=18$)           \\
\cmidrule(lr){2-8}
                                 & \multirow{2}{*}{Instruct} & Reader   & 0.885        & 1.000 ($N=8$)            & 0.753        & 0.133        & 0.012 ($N=18$)           \\
                                 &                    & Self     & 0.859        & 1.000 ($N=8$)            & 0.760        & 0.099        & 0.045 ($N=18$)           \\
\midrule
\multirow{4}{*}{\textbf{Llama 3.2 1B}} & \multirow{2}{*}{Base} & Reader   & 0.898        & 0.750 ($N=8$)            & 0.792        & 0.107        & 0.019 ($N=18$)           \\
                                 &                    & Self     & 0.888        & 1.000 ($N=8$)            & 0.797        & 0.091        & 0.013 ($N=18$)           \\
\cmidrule(lr){2-8}
                                 & \multirow{2}{*}{Instruct} & Reader   & 0.940        & 1.000 ($N=8$)            & 0.833        & 0.107        & 0.037 ($N=18$)           \\
                                 &                    & Self     & 0.940        & 1.000 ($N=8$)            & 0.831        & 0.109        & 0.044 ($N=18$)           \\
\midrule
\multirow{4}{*}{\textbf{Gemma 3 1B}} & \multirow{2}{*}{Base} & Reader   & 0.875        & 1.000 ($N=8$)            & 0.734        & 0.141        & 0.021 ($N=18$)           \\
                                 &                    & Self     & 0.867        & 1.000 ($N=8$)            & 0.721        & 0.146        & 0.010 ($N=18$)           \\
\cmidrule(lr){2-8}
                                 & \multirow{2}{*}{Instruct} & Reader   & 0.914        & 1.000 ($N=8$)            & 0.844        & 0.070        & 0.048 ($N=18$)           \\
                                 &                    & Self     & 0.922        & 1.000 ($N=8$)            & 0.844        & 0.078        & 0.084 ($N=18$)           \\
\midrule
\multirow{4}{*}{\textbf{OLMo 2 1B}} & \multirow{2}{*}{Base} & Reader   & 0.906        & 1.000 ($N=8$)            & 0.779        & 0.128        & 0.086 ($N=18$)           \\
                                 &                    & Self     & 0.917        & 1.000 ($N=8$)            & 0.797        & 0.120        & 0.046 ($N=18$)           \\
\cmidrule(lr){2-8}
                                 & \multirow{2}{*}{Instruct} & Reader   & 0.914        & 1.000 ($N=8$)            & 0.812        & 0.102        & 0.000 ($N=18$)           \\
                                 &                    & Self     & 0.898        & 1.000 ($N=8$)            & 0.815        & 0.083        & 0.032 ($N=18$)           \\
\bottomrule
\end{tabular}
\vspace{1ex}
\begin{minipage}{\linewidth}
\footnotesize
\textbf{Note:} Word Shuffleに対する一貫したaccuracy dropは、decodabilityが完全なbag-of-words shortcutだけでは説明しにくいことを支持する。一方、ParaphraseおよびOutcome Reversalはnonfallback sample数が小さいため、semantic validityに関する補助的evidenceとして解釈する。
\end{minipage}
\end{table}

\begin{table}[htbp]
\centering
\small
\caption{V1 E3（共有因果マップ）：刺激関連差分ベクトルへの介入に対するReaderおよびSelfの出力応答プロファイル。ピーク介入層、因果効果量（$\times 10^{-3}$）、プロファイル間の順位相関（Spearman $\rho_{\text{rank}}$）、および方向余弦類似度。}
\label{tab:v1_causal_profile}
\begin{tabular}{lll cccc cc}
\toprule
 & & & \multicolumn{2}{c}{\textbf{Reader Peak}} & \multicolumn{2}{c}{\textbf{Self Peak}} & \multicolumn{2}{c}{\textbf{Circuit Sharing}} \\
\cmidrule(lr){4-5} \cmidrule(lr){6-7} \cmidrule(lr){8-9}
\textbf{Family} & \textbf{Variant} & \textbf{Axis} & Layer ($d^*$) & Mag ($\times 10^{-3}$) & Layer ($d^*$) & Mag ($\times 10^{-3}$) & Spearman $\rho_{\text{rank}}$ & Mean Cosine \\
\midrule
\multirow{2}{*}{\textbf{Qwen 2.5 1.5B}} & Base       & Valence & L16 (0.59)   & 2.56       & L16 (0.59)   & 2.46       & 0.805      & 0.390      \\
\cmidrule(lr){2-9}
                                 & Instruct   & Valence & L1 (0.04)    & 0.98       & L8 (0.30)    & 0.86       & 0.799      & 0.000      \\
\midrule
\multirow{2}{*}{\textbf{Llama 3.2 1B}} & Base       & Valence & L4 (0.27)    & 0.75       & L6 (0.40)    & 1.19       & 0.582      & 0.000      \\
\cmidrule(lr){2-9}
                                 & Instruct   & Valence & L6 (0.40)    & 2.19       & L6 (0.40)    & 1.92       & 0.882      & 0.414      \\
\midrule
\multirow{2}{*}{\textbf{Gemma 3 1B}} & Base       & Valence & L2 (0.08)    & 2.87       & L2 (0.08)    & 3.70       & 0.802      & 0.712      \\
\cmidrule(lr){2-9}
                                 & Instruct   & Valence & L11 (0.44)   & 4.89       & L11 (0.44)   & 8.21       & 0.646      & 0.266      \\
\midrule
\multirow{2}{*}{\textbf{OLMo 2 1B}} & Base       & Valence & L4 (0.27)    & 6.14       & L4 (0.27)    & 1.66       & 0.715      & 0.210      \\
\cmidrule(lr){2-9}
                                 & Instruct   & Valence & L10 (0.67)   & 0.82       & L10 (0.67)   & 1.17       & 0.771      & 0.000      \\
\bottomrule
\end{tabular}
\vspace{1ex}
\begin{minipage}{\linewidth}
\footnotesize
\textbf{Note:} ReaderとSelfのlayer-wise causal profilesには中程度から高い正の順位相関が観測されたが、その強さはmodel間で異なり、intervention directionのcosine similarityがほぼ0となる条件も存在した。したがって、両taskはlayer-level causal sensitivityを部分的に共有するが、同一のcausal directionまたは同一回路を利用するとは結論しない。
\end{minipage}
\end{table}

\begin{table}[htbp]
\centering
\small
\caption{V1 E4（因果交換可能性）：Readerタスクで形成された差分ベクトルをSelf計算へ移植した際の因果効果。適合効果（Matched Effect）、ランダム統制効果（Random Mean）、特異性（Specificity $= M - R$）、95\%信頼区間、有意確率（$p$ / FDR $q$）、および転移比率（Transfer Ratio）。}
\label{tab:v1_interchangeability}
\begin{tabular}{lll ccccc c}
\toprule
\textbf{Family} & \textbf{Variant} & \textbf{Axis} & \textbf{Matched ($\times 10^{-3}$)} & \textbf{Random ($\times 10^{-3}$)} & \textbf{Specificity ($\times 10^{-3}$)} & \textbf{95\% CI ($\times 10^{-3}$)} & $p$ / $q$ & \textbf{Transfer Ratio} \\
\midrule
\multirow{4}{*}{\textbf{Qwen 2.5 1.5B}} & \multirow{2}{*}{Base} & Valence  & 0.752            & 0.311            & 0.441            & [-0.56, 1.42]        & 0.398 / 1.000  & ---          \\
                                 &                    & Arousal  & -0.109           & -0.038           & -0.071           & [-0.26, 0.13]        & 0.492 / 0.985  & ---          \\
\cmidrule(lr){2-9}
                                 & \multirow{2}{*}{Instruct} & Valence  & 0.635            & 0.355            & 0.280            & [-0.33, 0.89]        & 0.357 / 0.762  & ---          \\
                                 &                    & Arousal  & -0.182           & -0.156           & -0.027           & [-0.33, 0.30]        & 0.867 / 1.000  & ---          \\
\midrule
\multirow{4}{*}{\textbf{Llama 3.2 1B}} & \multirow{2}{*}{Base} & Valence  & -0.229           & -0.014           & -0.214           & [-0.79, 0.37]        & 0.459 / 1.000  & ---          \\
                                 &                    & Arousal  & -0.276           & -0.126           & -0.150           & [-0.34, 0.05]        & 0.132 / 0.353  & ---          \\
\cmidrule(lr){2-9}
                                 & \multirow{2}{*}{Instruct} & Valence  & 1.058            & 1.031            & 0.027            & [-0.94, 0.91]        & 0.954 / 1.000  & ---          \\
                                 &                    & Arousal  & -0.211           & -0.028           & -0.183           & [-0.52, 0.22]        & 0.319 / 0.639  & ---          \\
\midrule
\multirow{4}{*}{\textbf{Gemma 3 1B}} & \multirow{2}{*}{Base} & Valence  & -0.777           & -0.107           & -0.670           & [-5.00, 2.70]        & 0.717 / 1.000  & ---          \\
                                 &                    & Arousal  & 0.099            & -0.085           & 0.184            & [-2.06, 2.20]        & 0.866 / 1.000  & ---          \\
\cmidrule(lr){2-9}
                                 & \multirow{2}{*}{Instruct} & Valence  & 1.481            & -0.008           & 1.489            & [-0.96, 3.81]        & 0.224 / 0.896  & ---          \\
                                 &                    & Arousal  & 0.459            & 0.476            & -0.017           & [-0.97, 0.90]        & 0.973 / 1.000  & ---          \\
\midrule
\multirow{4}{*}{\textbf{OLMo 2 1B}} & \multirow{2}{*}{Base} & Valence  & -0.621           & -0.401           & -0.220           & [-0.77, 0.35]        & 0.453 / 1.000  & ---          \\
                                 &                    & Arousal  & -0.270           & -0.219           & -0.051           & [-0.30, 0.21]        & 0.727 / 1.000  & ---          \\
\cmidrule(lr){2-9}
                                 & \multirow{2}{*}{Instruct} & Valence  & 0.081            & 0.060            & 0.022            & [-0.57, 0.61]        & 0.946 / 1.000  & ---          \\
                                 &                    & Arousal  & -0.253           & -0.023           & -0.230           & [-0.52, 0.04]        & 0.109 / 0.354  & ---          \\
\bottomrule
\end{tabular}
\vspace{1ex}
\begin{minipage}{\linewidth}
\footnotesize
\textbf{Note:} Reader由来のstimulus-associated hidden-state differenceをSelfへ移植した際のmatched-minus-random specificityは、いずれのPrimary conditionでもFDR補正後に有意ではなかった。したがって、本解析からReader由来representationのcross-task causal interchangeabilityを支持するrobust evidenceは得られなかった。
\end{minipage}
\end{table}

\begin{table}[htbp]
\centering
\small
\caption{V1 E6（タスク特異的因果特殊化）：ReaderおよびSelfに特異的な因果局所部位（Selective Depth）の同定、各部位における切除効果（Ablation Effect）、および Task $\times$ SiteType 交互作用効果（$\beta_{\text{int}}, p$, FDR $q$）。}
\label{tab:v1_specialization}
\begin{tabular}{lll cccc ccc}
\toprule
 & & & \multicolumn{2}{c}{\textbf{Reader Site Effect}} & \multicolumn{2}{c}{\textbf{Self Site Effect}} & \multicolumn{3}{c}{\textbf{Interaction}} \\
\cmidrule(lr){4-5} \cmidrule(lr){6-7} \cmidrule(lr){8-10}
\textbf{Family} & \textbf{Variant} & \textbf{Selective Sites} & On Reader & On Self & On Reader & On Self & $\beta_{\text{int}}$ & $p$ & FDR $q$ \\
\midrule
\multirow{2}{*}{\textbf{Qwen 2.5 1.5B}} & Base       & (L1, L17)          & 0.065      & 0.047      & 0.032      & 0.027      & 0.0124     & 0.005      & 0.010      \\
\cmidrule(lr){2-10}
                                 & Instruct   & (L6, L27)          & 0.033      & 0.022      & 0.000      & 0.000      & 0.0109     & 7.04e-07   & 1.88e-06   \\
\midrule
\multirow{2}{*}{\textbf{Llama 3.2 1B}} & Base       & (L6, L7)           & 0.005      & 0.004      & 0.007      & 0.008      & 0.0010     & 0.035      & 0.056      \\
\cmidrule(lr){2-10}
                                 & Instruct   & (L2, L15)          & 0.064      & 0.050      & 0.000      & 0.000      & 0.0136     & 3.03e-10   & 1.21e-09   \\
\midrule
\multirow{2}{*}{\textbf{Gemma 3 1B}} & Base       & (L0, L18)          & 0.056      & 0.052      & 0.006      & 0.006      & 0.0041     & 0.236      & 0.236      \\
\cmidrule(lr){2-10}
                                 & Instruct   & (L4, L11)          & 0.067      & 0.081      & 0.059      & 0.061      & -0.0121    & 0.057      & 0.067      \\
\midrule
\multirow{2}{*}{\textbf{OLMo 2 1B}} & Base       & (L4, L13)          & 0.066      & 0.008      & 0.002      & 0.003      & 0.0583     & 1.84e-135  & 1.48e-134  \\
\cmidrule(lr){2-10}
                                 & Instruct   & (L5, L10)          & 0.012      & 0.011      & 0.004      & 0.005      & 0.0018     & 0.058      & 0.067      \\
\bottomrule
\end{tabular}
\vspace{1ex}
\begin{minipage}{\linewidth}
\footnotesize
\textbf{Note:} Task $\times$ SiteType interactionはQwen Base / Instruct、Llama Base / Instruct、OLMo Baseでnominally significantであった一方、Gemma Base / InstructおよびOLMo Instructでは明確なsupportが得られなかった。したがってtask-specific causal specializationは一部modelで観測されたが、全familyに共通する性質ではなかった。
\end{minipage}
\end{table}



% ================================================================
% 10. V1考察
% ================================================================
\section{V1考察}
\label{sec:v1-discussion}

V1の結果は、Behavioral Stageで観測されたReader--Self couplingが単なる最終outputの相関だけではなく、両taskにaffect-relevant informationが内部的に存在することと整合する一方、そのrepresentationやcausal useが単純に同一であるわけではないことを示した。

まずE1では、Primary 4 model familiesのBase / Instructの双方において、ReaderとSelfのhidden statesからhuman Reader Valence / Arousalをheld-outでdecodeできた。多くのmodel / axisではReaderとSelfのpeak depthも近接していた。

これは、ReaderとSelfが完全に独立した情報処理を行っているという説明よりも、同じaffect-relevant informationが両taskの内部representationに保持されるという説明と整合する。

一方で、一部conditionsではpeak relative depthに0.2以上の乖離も観測された。したがって、
\[
\text{shared affect-relevant information}
\neq
\text{identical readout location}
\]
である。

E2では、ReaderとSelfのstimulus geometryについて全conditionsで正のRSAが観測された一方、direct cross-decodingは大きく負となるconditionsが多かった。Procrustes alignmentによって性能が改善したことから、ReaderとSelfのrepresentationはraw coordinate systemとして直接交換可能なのではなく、線形変換によって部分的に整列可能なgeometryを持つと解釈できる。この結果は、
\[
\text{similar relational geometry}
\neq
\text{same coordinate system}
\]
という区別を示す。

Phase Bのsemantic controlsでは、Word Shuffleによる一貫したperformance dropが見られた。したがってV1のdecodabilityを単純なbag-of-words lexical shortcutだけで説明することは困難である。一方、ParaphraseおよびOutcome Reversalのnonfallback sample数は小さいため、これらはsemantic validityを補強するevidenceとして限定的に解釈する。

Phase CのE3では、ReaderとSelfのlayer-wise causal profilesに中程度から高い正のrank correlationが見られた。多くのmodelsではcausal peakも同一または近接layerに存在した。したがってReaderとSelfは、affect-relevant informationをdecode可能なだけではなく、hidden-state perturbationに対するlayer-level sensitivityも部分的に共有していると考えられる。

しかし、direction cosine similarityはmodelによって大きく異なり、ほぼ0となるconditionsも存在した。したがって、
\[
\text{similar causal sensitivity profile}
\neq
\text{same causal direction}
\]
である。

これはE4でさらに明確になった。Reader由来のstimulus-associated hidden-state differenceをSelfへ移植したmatched interventionは、
random-direction controlに対してFDR補正後にrobustなspecificityを示さなかった。したがって、
\[
\text{Reader representation}
\not\Rightarrow
\text{causally interchangeable Self representation}
\]
である。

これは本研究にとって重要なnegative resultである。Behavioral coupling、linear decodability、representation geometry、layer-wise causal profileの類似だけから、ReaderとSelfが同じcausal codeを直接共有しているとは言えない。

最後にE6では、Task $\times$ SiteType interactionが一部modelsではFDR補正後にも支持された一方、全familyには一般化しなかった。したがってReaderとSelfの計算にはpartial causal specializationが存在する可能性があるが、普遍的な二回路構造を示すとは言えない。

以上からV1は、ReaderとSelfはaffect-relevant informationとlayer-level causal sensitivityを部分的に共有するが、そのrepresentationは直接交換可能ではなく、causal organizationにはtask-specificな差も存在するという構造を支持する。



% ================================================================
% 11. V1からV2への接続
% ================================================================
\section{V1からV2への接続}
\label{sec:v1-to-v2}

V1では、ReaderとSelfの双方にaffect-relevant informationが存在し、representation geometryおよびlayer-wise causal sensitivityにも一定の共有性があることが示された。

一方、direct cross-decodingは弱く、cross-task causal interchangeabilityも支持されなかった。さらに、peak depth、causal magnitude、
intervention direction、task-specific site sensitivityにはmodel間およびBase / Instruct間で差が存在した。

ここからBaseからInstructへの変化に伴って、このrepresentation--report relationの何が変わるのかという問いが生じる。

V1でBaseとInstructの双方にaffect-relevant informationが存在したため、post-training後のbehavioral changeを単純なinformation erasureだけで説明できるとは限らない点が重要である。
可能性としては、
\begin{enumerate}
    \item representation geometryの変換
    \item Reader--Self sharingの再編
    \item causal leverageの再配置
    \item representationからoutput distributionへのreadout変更
\end{enumerate}
の4つが区別される。

そこでV2では比較軸を
\[
\text{Reader}\leftrightarrow\text{Self}
\]
から、
\[
\text{Base}\leftrightarrow\text{Instruct}
\]
へ変更し、Matched-Plain conditionをPrimary comparisonとしてprompt-format differenceを統制した上で、post-training-associated reorganizationを検証する。




% ================================================================
% 12. V2概要
% ================================================================
\section{V2概要}
\label{app:v2-methods}

\subsection{目的と位置づけ}

V2の中心的なResearch Questionは、「BaseからInstructへの変化に伴って、情動関連情報のrepresentation geometry、Reader--Self sharing、causal leverageの層分布、およびoutputへのreadoutはどのように再編されるか」である。
比較軸は、同一モデルファミリー内における事前学習モデルと事後学習モデルの対比、
\[
\text{Base} \leftrightarrow \text{Instruct} \quad \text{within the same model family}
\]
である。オープンウェイトモデルの観察的比較（Across-model observational comparison）と、固定モデル内部での因果介入（Within-model causal intervention）を厳密に区別して評価する。


% ----------------------------------------------------------------
\subsection{V2における証拠構成}

V2は以下の4 RQから構成する。
\[
\begin{aligned}
&\text{RQ1: Representation Geometry}\\
&\quad\downarrow\\
&\text{RQ2: Reader--Self Sharing}\\
&\quad\downarrow\\
&\text{RQ3: Decodability--Causal Relocation}\\
&\quad\downarrow\\
&\text{RQ4: Cross-model Distribution Recovery}
\end{aligned}
\]

表~\ref{tab:v2-rq-overview}に各RQの役割をまとめる。

\begin{table}[htbp]
\centering
\small
\caption{V2を構成する4つのResearch Questions。}
\label{tab:v2-rq-overview}
\begin{tabularx}{\linewidth}{
@{}
>{\raggedright\arraybackslash}p{1.3cm}
>{\raggedright\arraybackslash}p{3.1cm}
>{\raggedright\arraybackslash}X
@{}
}
\toprule
\textbf{RQ} &
\textbf{Analysis} &
\textbf{Question} \\
\midrule
RQ1 &
Representation Geometry &
affect-relevant informationのdecodabilityと
Base--Instruct representation geometryはどのように異なるか \\

RQ2 &
Reader--Self Sharing &
ReaderからSelfおよびSelfからReaderへの
cross-decoding可能性はBase--Instruct間でどのように変化するか \\

RQ3 &
Decodability--Causal Relocation &
情報がdecode可能なlayerと、
介入によってoutputを変化させるlayerは一致するか。
その関係はBase--Instruct間でどう変化するか \\

RQ4 &
Distribution Recovery &
Base hidden stateをInstructへ移植することで、
InstructのVA output distributionをBase distributionへ
どの程度戻せるか \\
\bottomrule
\end{tabularx}
\end{table}

V2では、
単純な「情報が残った／消えた」という二値的説明ではなく、

\[
\boxed{
\text{Representation}
+
\text{Geometry}
+
\text{Sharing}
+
\text{Causal Localization}
+
\text{Recoverability}
}
\]

へ分解してpost-training-associated changeを評価する。


% ----------------------------------------------------------------
\subsection{データ、外部ラベル、および分割}

V2では、前節で定義したEmoBank official test splitに基づく1,000 stimuliを使用する。V2におけるexternal affective targetsは、EmoBankのhuman Reader annotationである。
\[
Y_V
=
V_{\mathrm{human\ reader}}
\]

\[
Y_A
=
A_{\mathrm{human\ reader}}
\]
とする。Reader representationだけでなくSelf representationをprobeする場合にも、同じhuman Reader V/Aをexternal targetとして使用する。したがってV2で測定するdecodabilityは、
ReaderまたはSelfのhidden stateに、同じ外部human-affect variableがどの程度linearに表現されているかを意味する。Self-Reportそのものをprobe labelとして用いるわけではない。


\subsubsection{RQ1/RQ2のtrain--test split}

RQ1およびRQ2では、1,000 stimuliをseed 42でrandom permutationし、
\[
70\%\text{ training}
+
30\%\text{ held-out test}
\]
へ分割する。現行EmoBank evaluation setにはmatched-pair identifierが存在しないため、
本設定ではstimulus row単位のseeded splitとなる。したがって、
\[
N_{\mathrm{train}}=700,
\qquad
N_{\mathrm{test}}=300
\]
である。同一splitをすべてのmodel families、Base / Instruct conditions、Reader / Self tasksに共通して使用する。これにより、modelやprompt conditionごとに異なるstimulus subsetを評価することを避ける。


% ----------------------------------------------------------------
\subsection{Prompt-format controlと内部表現測定}
\label{app:v2-prompt-control}

前節\S\ref{app:prompt_format_control}で定義したPrompt-format control（Base Plain、Instruct Matched-Plain、Instruct Native-Chat）に基づき、V2のPrimary comparisonは、プロンプト形式の差異を排除した
\[
\text{Base Plain} \quad\text{vs.}\quad \text{Instruct Matched-Plain}
\]
（Matched-Plain contrast）とする。
Instruct Nativeとの比較はSecondary comparisonとし、両者の差（$\text{Instruct Native} - \text{Instruct Matched-Plain}$）をprompt-format effectとして分離評価する。

内部表現の測定は、\S\ref{app:overall_methods}の共通規約に従い、各Transformer block $l$ におけるprompt-end residual representation $h_{M,T,i,l}$ ($T \in \{R, S\}$, $M \in \{\text{Base}, \text{Instruct}\}$) を抽出し、相対深度 $d_l = l/(L-1)$ を一貫して用いてReaderおよびSelfの双方で実行する。


% ================================================================
% RQ1
% ================================================================

\subsection{RQ1：Representation Geometry}
\label{app:v2-rq1}

RQ1では、「BaseからInstructへの変化に伴い、affect-relevant representationのlocationとgeometryはどう変化するか」を検証する。ここでは、layer-wise affect decodability、Representational Similarity Analysis（RSA）\citep{kriegeskorte2008representational}、およびheld-out Orthogonal Procrustes disparityを用いる。


% ----------------------------------------------------------------
\subsection{RQ1-A：Layer-wise affect decodability}

ReaderおよびSelfの各layer representationから、human Reader Valence / Arousalを予測する。V1 Phase Aでは、
\[
\text{StandardScaler}
\rightarrow
\text{PCA}
\rightarrow
\text{Ridge}
\]
を使用した。これに対しV2 RQ1/RQ2では、post-trainingに伴うraw representationの変化を前処理によって除去しないため、
\[
\text{raw hidden state}
\rightarrow
\text{Ridge}
\]
を使用する。StandardScalerまたはPCAをPrimary Ridge probeの前処理として使用しない。各layerについて、
\[
\hat y
=
H_lw+b
\]
を、
\[
\text{Ridge}(\alpha=1)
\]
でtraining 700 stimuliにfitし、held-out 300 stimuli上で、
\[
R^2(l)
\]
を評価する。これにより、Base Reader、Base Self、Instruct Native Reader / Self、Instruct Matched-Plain Reader / Selfについて、Valence / Arousalそれぞれのlayer-wise decodability profileを得る。


% ----------------------------------------------------------------
\subsection{Decodability peak}

task \(T\) のBase profileについて、
\[
d_{D,B,T}^{*}
\]
をdecodability peakのrelative depthとする。Matched-Plain Instructについて、
\[
d_{D,I,T}^{*}
\]
を求め、
\[
\Delta d_{D,T}^{*}
=
d_{D,I,T}^{*}
-
d_{D,B,T}^{*}
\]
をpost-training-associated peak shiftとして記述する。positive valueは、Instruct側のdecodability maximumがBaseより後段に位置することを表す。negative valueは前段へのshiftを表す。Native-chat conditionについても同様のpeak shiftをSecondaryとして算出する。


% ----------------------------------------------------------------
\subsection{RQ1-B：Representational Similarity Analysis}

decodability profileだけでは、BaseとInstructで同じstimuliがhidden space内にどのようなrelative geometryを形成するかは分からない。そこでtaskごとに、held-out test representationsからpairwise cosine distanceを計算する。Baseについて、
\[
D_B(i,j)
=
1-
\operatorname{CosSim}
(
h_{B,i},
h_{B,j}
)
\]
とし、Instructについて、
\[
D_I(i,j)
=
1-
\operatorname{CosSim}
(
h_{I,i},
h_{I,j}
)
\]
とする。二つのRepresentational Dissimilarity Matrixの上三角要素について、
\[
\rho_{\mathrm{RSA}}
=
\operatorname{Spearman}
(
D_B,
D_I
)
\]を計算する。Primary comparisonでは、
\[
\text{Base Plain}
\leftrightarrow
\text{Instruct Matched-Plain}
\]
のRSAを用いる。Secondaryとして、
\[
\text{Base Plain}
\leftrightarrow
\text{Instruct Native}
\]
および、
\[
\text{Instruct Matched-Plain}
\leftrightarrow
\text{Instruct Native}
\]
も計算する。高いRSAは、BaseとInstructでstimulus間のrelative geometryが類似していることを表す。ただしRSAはcoordinate-wise identityを要求しない。


% ----------------------------------------------------------------
\subsection{RQ1-C：Held-out Orthogonal Procrustes disparity}

BaseとInstructが同じ情報を保持していても、post-trainingに伴いrepresentation coordinate systemが回転・変形している可能性がある。そこでtraining splitだけからBase--Instruct間のorthogonal alignmentを推定し、held-out test splitでalignment後のrepresentation mismatchを評価する。


\subsubsection{Centering}

同一task、layerについて、training Base / Instruct representationsの平均を、
\[
\mu_B,
\qquad
\mu_I
\]
とする。
\[
X_B
=
H_B-\mu_B,
\]

\[
X_I
=
H_I-\mu_I
\]
として別々にcenterする。


\subsubsection{Common PCA}

hidden dimensionが64を超える場合、training BaseとInstructのcentered representationsを結合し、
\[
X_{\mathrm{joint}}
=
\begin{bmatrix}
X_B\\
X_I
\end{bmatrix}
\]
に対してcommon PCAをfitする。dimensionは、
\[
k
=
\min
(
64,
N_{\mathrm{train}}-1,
D
)
\]
とする。得られた共通basisを \(V_k\) とし、
\[
Z_B
=
X_BV_k,
\]

\[
Z_I
=
X_IV_k
\]
へ射影する。


\subsubsection{Orthogonal alignment}

training representationsについて、
\[
M
=
Z_B^\top Z_I
\]
を計算し、
\[
M
=
U\Sigma V^\top
\]
とすると、BaseをInstructへalignするorthogonal transformを、
\[
Q
=
UV^\top
\]
とする。training splitで推定したcenter、PCA basis、orthogonal transformだけをheld-out testへ適用する。


\subsubsection{Held-out disparity}

held-out test representationについて、
\[
Z_{B,\mathrm{test}}Q
\]
と、
\[
Z_{I,\mathrm{test}}
\]
のnormalized Frobenius discrepancyを、
\[
\mathcal{D}_{\mathrm{Proc}}
=
\frac{
\|
Z_{B,\mathrm{test}}Q
-
Z_{I,\mathrm{test}}
\|_F
}{
\|
Z_{I,\mathrm{test}}
\|_F+\epsilon
}
\]
として評価する。したがって、
\[
\mathcal{D}_{\mathrm{Proc}}
\]
は小さいほど、単純なorthogonal alignment後にBaseとInstructのrepresentation geometryがよく一致することを意味する。Primary comparisonはMatched-Plainであり、NativeおよびInstruct format effectはSecondaryとする。


% ----------------------------------------------------------------
\subsection{RQ1の解釈範囲}

RQ1では、Decodabilityと、Representation Geometryを分けて測定する。例えばInstructで同等のdecodabilityが維持されていても、RSAが低下したりProcrustes disparityが増加する場合には、
情報の存在自体よりもrepresentation geometryのreorganizationと整合する。逆にgeometryが類似していても、そのrepresentationが最終outputへ同じ形で利用されるとは限らない。


% ================================================================
% RQ2
% ================================================================

\subsection{RQ2：Reader--Self Representation Sharing}
\label{app:v2-rq2}

RQ2では、V1で検証したReader--Self relationがBaseからInstructへの変化に伴ってどのように変わるかを調べる。中心的な問いは、「ReaderとSelfで共有されるaffect-relevant linear geometryは、Base--Instruct間でどのように変化するか」である。


% ----------------------------------------------------------------
\subsection{Bidirectional cross-decoding}

各layer、axisについて、Reader representationからhuman Reader affect targetを予測するRidge probeを学習する。Reader-trained probeをheld-out Self representationへ直接適用し、
\[
R^2_{R\rightarrow S}(l)
\]
を得る。同様に、Self representationで学習したprobeをheld-out Reader representationへ適用し、
\[
R^2_{S\rightarrow R}(l)
\]
を得る。cross-decoding時にもsource / target representationをtaskごとに再標準化しない。
したがってcross-decoding performanceは、ReaderとSelfのhidden-space coordinate systemがそのままどの程度互換であるかを反映する。


% ----------------------------------------------------------------
\subsection{Reader--Self sharing score}

双方向cross-decodingの平均を、
\[
S(l)
=
\frac{
R^2_{R\rightarrow S}(l)
+
R^2_{S\rightarrow R}(l)
}{2}
\]
と定義する。Baseについて、
\[
S_B(l)
\]
を求め、Matched-Plain Instructについて、
\[
S_{I,\mathrm{MP}}(l)
\]
を求める。Primary post-training-associated sharing changeは、
\[
\Delta S_{\mathrm{MP}}(l)
=
S_{I,\mathrm{MP}}(l)
-
S_B(l)
\]
である。layer全体のsummaryとして、
\[
\overline{\Delta S}_{\mathrm{MP}}
=
\frac{1}{L}
\sum_l
\Delta S_{\mathrm{MP}}(l)
\]
を用いる。


% ----------------------------------------------------------------
\subsection{Native and format controls}

Native Instructについても、
\[
S_{I,\mathrm{Native}}(l)
\]
を計算し、
\[
\Delta S_{\mathrm{Native}}(l)
=
S_{I,\mathrm{Native}}(l)
-
S_B(l)
\]
をSecondary comparisonとする。さらにprompt-format effectとして、
\[
\Delta S_{\mathrm{format}}(l)
=
S_{I,\mathrm{Native}}(l)
-
S_{I,\mathrm{MP}}(l)
\]
を求める。これにより、Base--Instruct changeと、plain--chat format changeを分離する。


% ----------------------------------------------------------------
\subsection{Directional asymmetry}

Secondary measureとして、Reader-to-SelfとSelf-to-Readerの変化が非対称であるかを記述する。
\[
\Delta\Delta_{\mathrm{cross}}
=
[
R^2_{R\rightarrow S,I}
-
R^2_{R\rightarrow S,B}
]
-
[
R^2_{S\rightarrow R,I}
-
R^2_{S\rightarrow R,B}
]
\]
を算出する。これはPrimary sharing scoreではなく、cross-task transferのdirectional asymmetryを補助的に記述するために使用する。


% ----------------------------------------------------------------
\subsection{RQ2の解釈範囲}

\[
\Delta S_{\mathrm{MP}}>0
\]
であれば、Matched-Plain条件においてReader--Self間のlinear transferabilityがInstruct側で大きいことを意味する。一方、
\[
\Delta S_{\mathrm{MP}}<0
\]
であれば、同じexternal affect informationは保持されていても、Reader--Self間のlinear coordinate sharingが低下している可能性と整合する。ただしcross-decodingは依然として
observational representational measureである。したがって、
\[
\text{Cross-decoding}
\neq
\text{Causal Utilization}
\]
であり、実際にどのlayerがoutputを動かすかはRQ3で検証する。


% ================================================================
% RQ3
% ================================================================

\subsection{RQ3：Decodability--Causal Relocation}
\label{app:v2-rq3}

RQ3では、「情動情報が最もdecode可能なlayerと、その情報への介入が最も強くoutputを変えるlayerは一致するか」を検証する。さらに、decodabilityとcausal leverageの関係がBase--Instruct間でどのように変化するかを評価する。RQ3では、Base Plain、Instruct Matched-Plain、Instruct Nativeの各条件について、Reader / Selfの双方を解析する。


% ----------------------------------------------------------------
\subsection{Cross-fitted affect directions}

RQ3では固定の70/30 splitではなく、全stimuliを5-fold cross-fittingで使用する。各foldでtraining stimuliだけを用いて、layer \(l\) のhidden stateからhuman Reader Valence / Arousalを予測するRidge regressionを学習する。Valenceを、
\[
Y_V
=
H_lw_V+b_V+\epsilon
\]
Arousalを、
\[
Y_A
=
H_lw_A+b_A+\epsilon
\]
とする。RQ3でのRidge regularizationは、
\[
\alpha_{\mathrm{Ridge}}=10
\]
である。係数vectorをunit normへ正規化し、
\[
d_V
=
\frac{w_V}{\|w_V\|_2}
\]

\[
d_A
=
\frac{w_A}{\|w_A\|_2}
\]
をaffect-related intervention directionsとする。directionはtraining foldだけから推定し、同じfoldのheld-out stimuliに対してのみ介入評価を行う。したがって各stimulusに適用されるdirectionは、そのstimulus自身のlabelを用いてfitされたものではない。


% ----------------------------------------------------------------
\subsection{Intervention dose normalization}

training hidden statesのtarget directionへのprojectionを、
\[
z_{V}
=
H_{\mathrm{train}}d_V
\]
とする。その標準偏差を、
\[
\sigma_V
=
\operatorname{SD}(z_V)
\]
とする。Arousalについても、
\[
\sigma_A
=
\operatorname{SD}
(
H_{\mathrm{train}}d_A
)
\]
を求める。介入は、
\[
h'
=
h+\alpha\sigma_d d
\]
として行う。Primary doseは、
\[
\alpha=1
\]
である。したがって介入量は、raw hidden-coordinate scaleではなく、当該training representationにおけるdirectional activationの約1 standard deviationを基準とする。


% ----------------------------------------------------------------
\subsection{Random and orthogonal controls}

任意のhidden-state perturbationが深いlayerほど大きなoutput changeを起こす可能性を統制するため、affect directionに加えて、
\begin{enumerate}
    \item random direction
    \item affect directionに直交するdirection
\end{enumerate}
を生成する。control directionsはtarget directionと同じEuclidean normを持つように正規化する。Valence interventionでは、affect \(d_V\)、random \(d_{V,\mathrm{rand}}\)、orthogonal \(d_{V,\perp}\) に対して同じ \(\sigma_V\) と \(\alpha=1\) を使用する。Arousalについても同様である。したがってcontrolとの差は、単なるperturbation magnitudeの違いでは説明されにくい設計となる。


% ----------------------------------------------------------------
\subsection{Axis-specific raw causal displacement}

Valence directionを介入した後のexpected Valenceを、
\[
E[V]_{\mathrm{affect}}
\]
clean baselineを、
\[
E[V]_{\mathrm{clean}}
\]
とする。raw Valence causal displacementを、
\[
C_{V,\mathrm{raw}}
=
\left|
E[V]_{\mathrm{affect}}
-
E[V]_{\mathrm{clean}}
\right|
\]
とする。randomおよびorthogonal controlsについて、
\[
C_{V,\mathrm{rand}},
\qquad
C_{V,\perp}
\]
を同様に定義する。Arousalについても、
\[
C_{A,\mathrm{raw}}
=
\left|
E[A]_{\mathrm{affect}}
-
E[A]_{\mathrm{clean}}
\right|
\]
を用いる。したがってraw causal displacement自体はnon-negativeなabsolute effectである。


% ----------------------------------------------------------------
\subsection{Net causal effect}

RQ3のPrimary causal metricは、raw affect-direction effectそのものではなく，random perturbationによるgeneric sensitivityを差し引いた、
\[
C_{V,\mathrm{net,rand}}
=
C_{V,\mathrm{raw}}
-
C_{V,\mathrm{rand}}
\]

\[
C_{A,\mathrm{net,rand}}
=
C_{A,\mathrm{raw}}
-
C_{A,\mathrm{rand}}
\]
をPrimaryとする。また，orthogonal controlを用いた、
\[
C_{V,\mathrm{net,\perp}}
=
C_{V,\mathrm{raw}}
-
C_{V,\perp}
\]

\[
C_{A,\mathrm{net,\perp}}
=
C_{A,\mathrm{raw}}
-
C_{A,\perp}
\]
をSecondary robustness measureとする。prompt-end residual全体を0へ置換する
zero-ablation effectも非特異的なSecondary controlとして記録する。
\[
C_{\mathrm{net,rand}}
\]
はsigned quantityであり、負値を取り得る。負値は、affect directionへの介入効果が同doseのrandom perturbationを下回ったことを意味する。statistical analysisでは、このsigned net effectを0へclipせずそのまま保持する。


% ----------------------------------------------------------------
\subsection{Layer-wise causal profile}

各stimulusはcross-fittingによって一度だけheld-out intervention evaluationを受ける。各layerについてsample-level net causal effectを平均し、
\[
C_{V,\mathrm{net}}(l)
\]
および、
\[
C_{A,\mathrm{net}}(l)
\]
というlayer profileを得る。Primary profileには、
\[
C_{\mathrm{net,rand}}
\]
を使用する。


% ----------------------------------------------------------------
\subsection{Decodability--causal dissociation}

同じcondition、task、axisについて、RQ1/RQ2から得られたheld-out decodability profileを、
\[
D(l)
=
R^2(l)
\]
とする。RQ3から得られたPrimary net causal profileを、
\[
C(l)
=
C_{\mathrm{net,rand}}(l)
\]
とする。


\subsubsection{Positive decodability peak}

decodabilityについて、
\[
\max_lD(l)\le0
\]
の場合には、
\[
d_D^{*}
=
\mathrm{NaN}
\]
とする。すなわち、全layerでheld-out \(R^2\le0\) である場合に、「最も悪くないlayer」をdecodability peakとして選択しない。


\subsubsection{Positive causal peak}

同様に、
\[
\max_l C(l)\le0
\]
の場合には、
\[
d_C^{*}
=
\mathrm{NaN}
\]
とする。これは、どのlayerでもaffect-specific perturbationがrandom controlを上回らなかった場合に、causal peakが存在したかのように記述しないためである。


\subsubsection{Peak dissociation}

両peakが定義可能な場合にのみ、
\[
\Delta d^{*}
=
d_C^{*}
-
d_D^{*}
\]
を定義する。positive valueは、
\[
d_C^{*}>d_D^{*}
\]
すなわち、最大causal leverageが最大decodabilityより後段に位置することを表す。


% ----------------------------------------------------------------
\subsection{Center-of-mass dissociation}

単一のargmaxだけではlayer profile全体の移動を捉えにくいため、positive profile massを用いたcenter of massも計算する。decodabilityについて、
\[
\bar d_D
=
\frac{
\sum_l
d_l
\max(D(l),0)
}{
\sum_l
\max(D(l),0)
}
\]
とする。positive decodability massが存在しない場合には\(\bar d_D=\mathrm{NaN}\) とする。causal profileについて、
\[
\bar d_C
=
\frac{
\sum_l
d_l
\max(C(l),0)
}{
\sum_l
\max(C(l),0)
}
\]
とする。positive net causal massが存在しない場合には\(\bar d_C=\mathrm{NaN}\) とする。
両者が定義可能な場合、
\[
\Delta\bar d
=
\bar d_C-\bar d_D
\]
をprofile-level dissociation measureとする。


% ----------------------------------------------------------------
\subsection{Post-training-associated change in dissociation}

Base task \(T\) について、
\[
\Delta d_{B,T}^{*}
=
d_{C,B,T}^{*}
-
d_{D,B,T}^{*}
\]
を求める。Matched-Plain Instructについて、
\[
\Delta d_{I,T}^{*}
=
d_{C,I,T}^{*}
-
d_{D,I,T}^{*}
\]
とする。その差、
\[
\Gamma_T^{\mathrm{peak}}
=
\Delta d_{I,T}^{*}
-
\Delta d_{B,T}^{*}
\]
を、post-training-associated change in decodability--causal peak dissociationとして記述する。center of massについても、
\[
\Gamma_T^{\mathrm{CoM}}
=
\Delta\bar d_{I,T}
-
\Delta\bar d_{B,T}
\]
を求める。Primary comparisonはMatched-Plain、Native-chatはSecondaryである。


% ----------------------------------------------------------------
\subsection{Sample-level mixed-effects analysis}

peak / center-of-mass summaryに加えて、layer profile全体を使用したsample-level analysisを行う。Primary outcomeは、
\[
C_{\mathrm{net,rand}}
\]
である。ValenceとArousalを別々に解析し、
\[
C_{\mathrm{net,rand}}
\sim
C(\mathrm{family})
+
C(\mathrm{alignment})
\times
C(\mathrm{task})
\times
\mathrm{relative\ depth}
+
(1|\mathrm{stimulus})
\]
というLinear Mixed-Effects Modelを使用する。ここで、
\[
\mathrm{alignment}
\in
\{
\text{Base},
\text{Instruct}
\},
\]

\[
\mathrm{task}
\in
\{
\text{Reader},
\text{Self}
\}
\]
である。EmoBankにはmatched pair structureがないため、random interceptのgroupは各stimulus identityに対応する。Primary LMMには、Base Plainと、Instruct Matched-Plainだけを含める。Native chatを使用した同一formulaのanalysisはSecondaryとする。


% ----------------------------------------------------------------
\subsection{Prespecified RQ3 interaction terms}

Primary confirmatory termsは、
\[
\mathrm{Alignment}
\times
\mathrm{RelativeDepth}
\]
\[
\mathrm{Alignment}
\times
\mathrm{Task}
\]
\[
\mathrm{Alignment}
\times
\mathrm{Task}
\times
\mathrm{RelativeDepth}
\]
である。Valence / Arousalについて得られるこれらprespecified interaction termsのp-valuesに対して、Benjamini--Hochberg FDR correctionを行う。これにより、単一peakだけでなく、Base--Instruct間でcausal leverage profileの形状が変化するかをsample-levelで評価する。


% ----------------------------------------------------------------
\subsection{RQ3の解釈上の境界}

RQ3はhidden-state interventionを行うため、各固定model内部ではcausal evidenceを提供する。ただし、
\[
\Gamma_T^{\mathrm{peak}}
\]
やLMMのalignment interactionは、BaseとInstructという別model artifactsの比較である。したがって、post-training caused the causal peak to moveとは記述せず、the causal profile differs between Base and Instruct in a manner associated with post-trainingと解釈する。


% ================================================================
% RQ4
% ================================================================

\subsection{RQ4：Cross-model Distribution Recovery}
\label{app:v2-rq4}

RQ1--RQ3では、BaseとInstructのrepresentation geometry、sharing、causal localizationを比較する。しかし、これらの違いだけでは、「Baseに存在するrepresentationを
Instruct computationへ直接戻したときに、Base型のoutput distributionを回復できるか」は分からない。そこでRQ4では、"Can Base hidden representations recover a Base-like
VA output distribution when inserted into the Instruct model?"を検証する。これは同一stimulusについて、Baseで自然に形成されたhidden representationを同一familyのInstruct modelへ移植するcross-model activation replacementである。


% ----------------------------------------------------------------
\subsection{RQ4のtrain--evaluation split}

1,000 EmoBank stimuliをseed 42で、
\[
70\%\text{ train}
+
30\%\text{ independent evaluation}
\]
へ分割する。現行dataにはpair structureがないため、
\[
N_{\mathrm{train}}=700,
\qquad
N_{\mathrm{eval}}=300
\]
のstimulus-level permutation splitとなる。このsplitは同一family内でReaderとSelfに共通して使用する。train subsetはProcrustes alignmentの推定だけに使用し、recovery evaluationは独立した300 stimuli上で行う。


% ----------------------------------------------------------------
\subsection{Target and clean distributions}

stimulus \(i\)、task \(T\) について、Base Plain modelの81-state VA distributionを、
\[
P_{B,i}^{T}
\]
とし、recoveryのtarget distributionとする。Instruct clean distributionはconditionによって異なる。Native conditionでは、
\[
P_{I,\mathrm{Native},i}^{T}
\]
Matched-Plain conditionでは、
\[
P_{I,\mathrm{MP},i}^{T}
\]
である。したがってPrimary matched-plain recoveryは、
\[
P_{I,\mathrm{MP}}
\rightarrow
P_B
\]
の距離がBase activation patchingによってどの程度縮小するかを測定する。


% ----------------------------------------------------------------
\subsection{Joint VA Earth Mover's Distance}

distribution間のPrimary distanceには、81-state VA joint distribution上のEarth Mover's Distanceを使用する。各candidate stateを、
\[
c=(v,a),
\qquad
v,a\in\{1,\ldots,9\}
\]
とし、state間ground costを、
\[
g(c,c')
=
\sqrt{
(v-v')^2
+
(a-a')^2
}
\]
とする。二つのjoint distributions、
\[
P,
\qquad
Q
\]
間で、このground costを用いたoptimal transport costを、
\[
EMD_{VA}(P,Q)
\]
とする。この2D joint EMDをPrimary distribution distanceとする。Valence marginalの1D Wasserstein distance、Arousal marginalの1D Wasserstein distance、
およびJensen--Shannon divergenceもSecondary distribution diagnosticsとして保持する。
ただしRQ4のRecoveryにはjoint \(EMD_{VA}\)を使用する\citep{rubner2000earth}。


% ----------------------------------------------------------------
\subsection{Condition A：Direct Native patching}

各evaluation stimulusについて、Base Plain modelのlayer \(l\) のprompt-end post-MLP residual stateを、
\[
h_{B,i,l}
\]
とする。これを変更せず、Instruct Native modelの同じlayer / semantic positionへ直接置換する。
\[
h_{I,i,l}^{\mathrm{Native}}
\leftarrow
h_{B,i,l}
\]
とする。この条件は、raw cross-model direct interchangeabilityを評価する。ただしBaseとInstructでrepresentation coordinate systemが変化している場合には、Base activationがInstruct representation manifold上で不自然な位置に置かれる可能性がある。そのためDirect NativeはSecondary conditionとする。


% ----------------------------------------------------------------
\subsection{Primary condition：Matched-Plain direct patching}

prompt-format differenceを除外したPrimary conditionでは、Instructにもplain promptを使用する。Base Plainのactivationを、
\[
h_{B,i,l}
\]
Instruct Matched-Plainの同じlayerへ、
\[
h_{I,i,l}^{\mathrm{MP}}
\leftarrow
h_{B,i,l}
\]
として直接置換する。この条件では、Base / Instructの双方が同一のplain prompt formatを使用するため、Matched-Plain direct recoveryをRQ4のPrimary conditionとする。


% ----------------------------------------------------------------
\subsection{Mechanistic control：Procrustes-aligned patching}

direct patchingが失敗した場合、Base representationに情報が存在しないとは限らない。BaseとInstructでcoordinate systemがpost-trainingに伴って変化していれば、raw Base activationはInstruct側でoff-manifoldになる可能性がある。そこでNative conditionについて、train splitからBase representationをInstruct representationへalignするlinear transformationを学習する。


\subsubsection{Train-only centering and common PCA}

layer \(l\) のtraining Base / Instruct Native activationsを、
\[
H_B,
\qquad
H_I
\]
とする。それぞれのmeanを、
\[
\mu_B,
\qquad
\mu_I
\]
としてcenterする。
\[
X_B
=
H_B-\mu_B,
\]

\[
X_I
=
H_I-\mu_I.
\]
結合training matrix、
\[
X_{\mathrm{joint}}
=
\begin{bmatrix}
X_B\\
X_I
\end{bmatrix}
\]
からcommon PCA basis \(V_k\) を求める。dimensionは、
\[
k
=
\min
(
64,
N_{\mathrm{train}}-1,
D
)
\]
である。


\subsubsection{Base-to-Instruct orthogonal mapping}

PCA subspace内で、
\[
Z_B
=
X_BV_k,
\]

\[
Z_I
=
X_IV_k
\]
とする。training samplesから、
\[
M
=
Z_B^\top Z_I
\]
を計算し、SVDによってorthogonal transform \(Q\) を得る。evaluation Base activation \(h_B\) について、
\[
x
=
h_B-\mu_B
\]
をPCA subspace成分、
\[
x_{\mathrm{sub}}
=
(xV_k)V_k^\top
\]
と、その直交補空間成分、
\[
x_{\mathrm{res}}
=
x-x_{\mathrm{sub}}
\]
へ分解する。PCA subspace成分だけをProcrustes rotationし、
\[
z_{\mathrm{aligned}}
=
(xV_k)Q
\]
として、
\[
h_{B\rightarrow I}^{\mathrm{aligned}}
=
z_{\mathrm{aligned}}V_k^\top
+
x_{\mathrm{res}}
+
\mu_I
\]
を構成する。すなわち、common low-rank subspaceではBaseをInstruct座標系へrotateし、
PCA subspace外のresidual componentは保持する。このaligned representationを
Instruct Native promptへpatchする。この条件はPrimaryではなく、coordinate-remapping mechanistic controlとして扱う。


% ----------------------------------------------------------------
\subsection{Recovery ratio}

stimulus \(i\)、layer \(l\) について、clean Instruct distributionとBase target distributionの距離を、
\[
D_{\mathrm{clean},i}
=
EMD_{VA}
(
P_{I,i},
P_{B,i}
)
\]
とする。Base activationをpatchした後のdistributionを、
\[
P_{\mathrm{patch},i,l}
\]
とし、
\[
D_{\mathrm{patch},i,l}
=
EMD_{VA}
(
P_{\mathrm{patch},i,l},
P_{B,i}
)
\]
とする。Recovery ratioを、
\[
Recovery_{i,l}
=
\frac{
D_{\mathrm{clean},i}
-
D_{\mathrm{patch},i,l}
}{
D_{\mathrm{clean},i}
+
10^{-12}
}
\]
と定義する。この定義では、
\[
Recovery>0
\]
は、patchingによってInstruct distributionがBase targetに近づいたことを意味する。
\[
Recovery=0
\]
は距離が改善しない場合、
\[
Recovery<0
\]
はpatchingによってBase targetからさらに離れた場合を意味する。完全にtargetへ一致した場合には、
\[
Recovery=1
\]
となる。


% ----------------------------------------------------------------
\subsection{Layer-wise recovery profile}

各layerについてevaluation stimuliのRecoveryを平均し、
\[
\overline{Recovery}(l)
=
\frac{1}{N_{\mathrm{eval}}}
\sum_i
Recovery_{i,l}
\]
を求める。各layerについてsample bootstrap、
\[
B=1000
\]
による95\% confidence intervalを計算する。


% ----------------------------------------------------------------
\subsection{Primary depth-integrated recovery}

単一のbest layerだけをPrimary outcomeとすると、偶然のlayer-specific fluctuationや
post hoc peak selectionの影響を受けやすい。そこでRQ4のPrimary matched-plain metricは、relative depthに沿ったrecovery profile全体のArea Under the Curve（AUC）とする。
\[
AUC_{\mathrm{Recovery}}^{\mathrm{MP}}
=
\int_0^1
\overline{Recovery}_{\mathrm{MP}}(d)\,dd
\]
実際にはmodel layerのrelative-depth grid上でtrapezoidal integrationを行う。各evaluation stimulusについても個別のlayer-wise recovery profileからsample-level AUCを計算する。したがって、Matched-Plain AUC RecoveryをV2 RQ4のPrimary distributional recovery quantityとする。


% ----------------------------------------------------------------
\subsection{Absolute EMD improvement}

normalized Recoveryに加えて、absolute distance improvementを、
\[
\Delta EMD_{i,l}^{\mathrm{MP}}
=
D_{\mathrm{clean},i}^{\mathrm{MP}}
-
D_{\mathrm{patch},i,l}^{\mathrm{MP}}
\]
として記録する。そのlayer-wise meanおよびrelative-depth AUC、
\[
AUC_{\Delta EMD}^{\mathrm{MP}}
\]
も保存する。Recovery ratioは初期Base--Instruct distanceでnormalizeされたrelative improvementを表すのに対し、\(\Delta EMD\)はabsolute distributional improvementを表す。


% ----------------------------------------------------------------
\subsection{Peak recovery as a secondary localization measure}

各conditionについて、
\[
l_{\mathrm{rec}}^{*}
=
\arg\max_l
\overline{Recovery}(l)
\]
をbest recovery layerとし、
\[
d_{\mathrm{rec}}^{*}
=
\frac{
l_{\mathrm{rec}}^{*}
}{
L-1
}
\]
をbest recovery relative depthとする。また、
\[
\max_l\overline{Recovery}(l)
\]
をmaximum recovery ratioとして記録する。
ただしこれらはSecondary localization measuresであり、Primary RQ4 inferenceには
depth-integrated AUCを使用する。


% ----------------------------------------------------------------
\subsection{Reader--Self recovery comparison}

RQ4はReaderとSelfの双方で同一のrecovery protocolを実施する。Primary Matched-Plain conditionについて、
\[
AUC_{\mathrm{Self}}^{\mathrm{MP}}
\]
と、
\[
AUC_{\mathrm{Reader}}^{\mathrm{MP}}
\]
を比較し、
\[
\Delta AUC_{\mathrm{Recovery}}
=
AUC_{\mathrm{Self}}^{\mathrm{MP}}
-
AUC_{\mathrm{Reader}}^{\mathrm{MP}}
\]
をtask asymmetryのPrimary cross-family summaryとする。maximum recovery ratioのSelf--Reader differenceはSecondaryとする。Native-chat AUC differenceもSecondary、
Procrustes-aligned AUC differenceをmechanistic controlとして保持する。


% ----------------------------------------------------------------
\subsection{RQ4 sample-level inference}

sample-level matched-plain recoveryについて、familyをfixed effectとして含むmixed-effects analysisも行う。Primary AUC analysisは、
\[
AUC_{\mathrm{Recovery}}^{\mathrm{MP}}
\sim
C(\mathrm{family})
+
C(\mathrm{task})
+
(1|\mathrm{stimulus})
\]
である。Secondaryとして、global best layerにおけるmatched-plain recovery ratioについても、
\[
Recovery_{\mathrm{best}}^{\mathrm{MP}}
\sim
C(\mathrm{family})
+
C(\mathrm{task})
+
(1|\mathrm{stimulus})
\]
を評価する。


% ----------------------------------------------------------------
\subsection{RQ4の解釈上の境界}

高いRecoveryが得られた場合に示されるのは、Base-derived hidden representation can causally move the Instruct output distribution toward the Base distributionということである。これは、the model has an internal emotionという主張ではない。またDirect patchingでrecoveryが低い場合も、Base側にaffect-relevant informationが存在しないとは限らない。BaseとInstructのcoordinate systemが変化していれば、direct patch自体がoff-manifoldである可能性がある。このためProcrustes-aligned conditionをmechanistic controlとして併用する。


% ================================================================
% Confirmatory Cross-Family Analysis
% ================================================================

\subsection{4-family confirmatory integration}
\label{app:v2-confirmatory}

各familyについてRQ1--RQ4を完了した後、Primary cohortである
\[
\{
\text{Qwen},
\text{Llama},
\text{Gemma},
\text{OLMo}
\}
\]
の4 familiesを統合したconfirmatory analysisを行う。


% ----------------------------------------------------------------
\subsection{H1a：Geometric reorganization}

Matched-Plain conditionにおけるBase--Instruct Procrustes disparityとRSAについて、familyごとにlayer平均を求める。Reader / Selfそれぞれについて、
\[
\overline{\mathcal{D}}_{\mathrm{Proc}},
\qquad
\overline{\rho}_{\mathrm{RSA}}
\]
を得る。4 familiesに対して、
\[
B=1000
\]
のbootstrap confidence intervalを計算する。これはPrimary cohort内のcross-family descriptive summaryとして扱う。


% ----------------------------------------------------------------
\subsection{H1b：Decodability peak reorganization}

Matched-Plain条件について、BaseとInstructのdecodability peak relative depthの差、
\[
\Delta d_{D,T}^{*}
=
d_{D,I,T}^{*}
-
d_{D,B,T}^{*}
\]
をfamilyごとに算出する。Reader / Self、Valence / Arousalについて、4-family mean shiftとbootstrap 95\% CIを報告する。Native-chat peak shiftはSecondaryとする。


% ----------------------------------------------------------------
\subsection{H2：Reader--Self sharing reorganization}

各familyについて、Matched-Plainのlayer-wise sharing difference、
\[
\Delta S_{\mathrm{MP}}(l)
\]
をlayer平均し、
\[
\overline{\Delta S}_{\mathrm{MP}}
\]
をfamily-level quantityとする。4-family meanおよびbootstrap 95\% CIを算出する。Base Plain vs Instruct Nativeによる\(\overline{\Delta S}_{\mathrm{Native}}\)は
Secondaryとする。


% ----------------------------------------------------------------
\subsection{H3：Causal-profile reorganization}

Primary H3 inferenceは、前述のsample-level LMM、
\[
C_{\mathrm{net,rand}}
\sim
C(\mathrm{family})
+
C(\mathrm{alignment})
\times
C(\mathrm{task})
\times
d
+
(1|\mathrm{stimulus})
\]
に基づく。Primary conditionはMatched-Plainである。Valence / Arousalのprespecified interaction termsについてBenjamini--Hochberg FDR correctionを行う。


% ----------------------------------------------------------------
\subsection{H4：Distribution-recovery asymmetry}

Primary H4 family-level quantityは、
\[
AUC_{\mathrm{Self}}^{\mathrm{MP}}
-
AUC_{\mathrm{Reader}}^{\mathrm{MP}}
\]
である。4 familiesについてmean differenceとbootstrap 95\% CIを計算する。Native-chat AUC differenceはSecondary、Procrustes-aligned AUC differenceはmechanistic controlとする。maximum recovery differenceもSecondary localization summaryとして保持する。sample-level AUCについては、前述のmixed-effects analysisを用いる。


% ----------------------------------------------------------------
\subsection{Family-level bootstrapの位置づけ}

Primary cohortは4 model familiesであるため、family-level bootstrapにおける独立family数は、
\[
N_{\mathrm{family}}=4
\]
である。したがってfamily-level bootstrap intervalは、無限のmodel populationへの強い一般化を目的としたものではなく、prespecified primary familiesにおけるcross-family consistencyの記述として解釈する。sample-level variationを用いるRQ3 / RQ4のmixed-effects analysesとは、inference levelを区別する。


% ----------------------------------------------------------------
\subsection{V2におけるPrimary / Secondary evidenceの整理}

V2のPrimaryおよび主要Secondary analysesを表~\ref{tab:v2-primary-secondary}にまとめる。

\begin{table}[htbp]
\centering
\small
\caption{V2におけるPrimary evidenceと主要なcontrol / Secondary analyses。}
\label{tab:v2-primary-secondary}
\begin{tabularx}{\linewidth}{
@{}
>{\raggedright\arraybackslash}p{2.1cm}
>{\raggedright\arraybackslash}p{4.1cm}
>{\raggedright\arraybackslash}X
@{}
}
\toprule
\textbf{Analysis} &
\textbf{Primary} &
\textbf{Secondary / Control} \\
\midrule
RQ1 Geometry &
Base Plain vs Instruct Matched-Plain:
decodability,
cosine-RSA,
held-out Procrustes disparity &
Base vs Instruct Native;
Instruct Native vs Matched-Plain format effect \\

RQ2 Sharing &
Matched-Plain
\(\Delta S(l)\) and layer-mean
\(\overline{\Delta S}\) &
Native sharing difference;
directional cross-decoding asymmetry \\

RQ3 Causal map &
Matched-Plain
\(C_{\mathrm{net,rand}}\);
positive-guarded peak / center;
sample-level LMM &
\(C_{\mathrm{net,\perp}}\);
raw causal effect;
zero ablation;
Native condition \\

RQ4 Recovery &
Matched-Plain
depth-AUC recovery &
Matched-Plain \(\Delta EMD\) AUC;
peak recovery;
Native recovery;
Procrustes-aligned recovery \\
\bottomrule
\end{tabularx}
\end{table}


% ----------------------------------------------------------------

% ================================================================
% 13. V2結果
% ================================================================
\section{V2結果}
\label{sec:v2-results}

% ================================================================
% V2 Stage: Post-training Reorganization Complete Summary Tables
% ================================================================

\begin{table}[htbp]
\centering
\small
\caption{V2 H1--H2（幾何再編と表現共有度の変位）：事後学習に伴う表現幾何学的歪み（Procrustes Distortion）、デコードピーク深度変位 $\Delta d^*$、およびタスク共有度変化 $\Delta\text{Sharing}$。4ファミリー統合ブートストラップ95\%信頼区間。}
\label{tab:v2_h1_h2_reorganization}
\begin{tabular}{ll ccc}
\toprule
\textbf{Hypothesis} & \textbf{Metric} & \textbf{Estimate} & \textbf{95\% CI} & \textbf{Condition} \\
\midrule
\multirow{4}{*}{H1a: Geometric Distortion} & Reader Procrustes Distortion         & ---      & ---                & Matched-Plain \\
                                 & Self Procrustes Distortion           & ---      & ---                & Matched-Plain \\
                                 & RSA Reader ($\rho_{\text{RSA}}$)     & ---      & ---                & Matched-Plain \\
                                 & RSA Self ($\rho_{\text{RSA}}$)       & ---      & ---                & Matched-Plain \\
\midrule
\multirow{4}{*}{H1b: Decodability Peak Shift} & Valence Reader $\Delta d^*$          & ---      & ---                & Matched-Plain \\
                                 & Valence Self $\Delta d^*$            & ---      & ---                & Matched-Plain \\
                                 & Arousal Reader $\Delta d^*$          & ---      & ---                & Matched-Plain \\
                                 & Arousal Self $\Delta d^*$            & ---      & ---                & Matched-Plain \\
\midrule
\multirow{2}{*}{H2: Sharing Reorganization} & $\Delta\text{Sharing}$ (Valence)     & ---      & ---                & Matched-Plain \\
                                 & $\Delta\text{Sharing}$ (Arousal)     & ---      & ---                & Matched-Plain \\
\bottomrule
\end{tabular}
\vspace{1ex}
\begin{minipage}{\linewidth}
\footnotesize
\textbf{Note:} $\Delta d^* = d^*_{\text{instruct}} - d^*_{\text{base}} > 0$ は、事後学習によって表現デコードピークが後段側（deeper側）へシフトしたことを示す。また、$\Delta\text{Sharing} < 0$ はReaderとSelfの表現共有度合いが事後学習によってタスク分離方向に再編されたことを示す。
\end{minipage}
\end{table}

\begin{table}[htbp]
\centering
\small
\caption{V2 H3a（因果局在の再配置）：事後学習に伴う因果ピーク相対深度 $d_C^*$ および重心深度 $d_{\text{center}}$ の後段シフト量（$\Delta d_C^*, \Delta d_{\text{center}}$）。}
\label{tab:v2_causal_relocation}
\begin{tabular}{lll cccc}
\toprule
\textbf{Family} & \textbf{Task} & \textbf{Axis} & \textbf{Base Peak ($d^*$)} & \textbf{Instruct Peak ($d^*$)} & $\Delta d_C^*$ & $\Delta d_{\text{center}}$ \\
\midrule
GEMMA      & Reader   & Arousal  & 0.760          & 0.560            & $-$0.200   & $-$0.106   \\
GEMMA      & Reader   & Valence  & 0.520          & 0.080            & $-$0.440   & 0.098      \\
GEMMA      & Self     & Arousal  & 0.560          & 0.240            & $-$0.320   & 0.010      \\
GEMMA      & Self     & Valence  & 0.000          & 0.520            & 0.520      & 0.044      \\
LLAMA      & Reader   & Arousal  & ---            & ---              & ---        & ---        \\
LLAMA      & Reader   & Valence  & ---            & ---              & ---        & ---        \\
LLAMA      & Self     & Arousal  & ---            & ---              & ---        & ---        \\
LLAMA      & Self     & Valence  & ---            & ---              & ---        & ---        \\
OLMO       & Reader   & Arousal  & 0.200          & 0.600            & 0.400      & 0.043      \\
OLMO       & Reader   & Valence  & 0.600          & 0.267            & $-$0.333   & $-$0.170   \\
OLMO       & Self     & Arousal  & 0.067          & 0.333            & 0.267      & $-$0.087   \\
OLMO       & Self     & Valence  & 0.600          & 0.133            & $-$0.467   & $-$0.296   \\
QWEN       & Reader   & Arousal  & 0.037          & 0.000            & $-$0.037   & $-$0.144   \\
QWEN       & Reader   & Valence  & 0.778          & 0.111            & $-$0.667   & $-$0.042   \\
QWEN       & Self     & Arousal  & 0.296          & 0.852            & 0.556      & 0.027      \\
QWEN       & Self     & Valence  & 0.593          & 0.259            & $-$0.333   & $-$0.178   \\
\bottomrule
\end{tabular}
\vspace{1ex}
\begin{minipage}{\linewidth}
\footnotesize
\textbf{Note:} 因果介入におけるピークおよび重心深度変位 $\Delta d_C^* = d_{C,\text{Instruct}}^* - d_{C,\text{Base}}^*$ は、事後学習に伴う因果部位の後段移行量を示す。
\end{minipage}
\end{table}

\begin{table}[htbp]
\centering
\small
\caption{V2 H3b（因果特異性統制実験）：生の因果効果 $C_{\text{raw}}$、ランダム方向統制 $C_{\text{rand}}$、直交方向統制 $C_{\text{perp}}$、正味因果効果 $C_{\text{net,rand}} = C_{\text{raw}} - C_{\text{rand}}$、およびゼロ切除効果。}
\label{tab:v2_causal_controls}
\begin{tabular}{llll ccccc}
\toprule
\textbf{Family} & \textbf{Task} & \textbf{Condition} & \textbf{Axis} & $C_{\text{raw}}$ & $C_{\text{rand}}$ & $C_{\text{perp}}$ & $C_{\text{net,rand}}$ (\textbf{Primary}) & Zero-Ablation \\
\midrule
OLMO       & Reader   & Base-Plain       & Valence  & 0.002      & 0.002      & 0.002      & $-$0.000           & 0.016        \\
OLMO       & Reader   & Base-Plain       & Arousal  & 0.001      & 0.001      & 0.001      & 0.000              & 0.019        \\
OLMO       & Self     & Base-Plain       & Valence  & 0.002      & 0.002      & 0.002      & $-$0.000           & 0.016        \\
OLMO       & Self     & Base-Plain       & Arousal  & 0.001      & 0.001      & 0.001      & 0.000              & 0.020        \\
OLMO       & Reader   & Matched-Plain    & Valence  & 0.004      & 0.004      & 0.004      & $-$0.000           & 0.053        \\
OLMO       & Reader   & Matched-Plain    & Arousal  & 0.001      & 0.001      & 0.001      & $-$0.000           & 0.016        \\
OLMO       & Self     & Matched-Plain    & Valence  & 0.004      & 0.004      & 0.004      & $-$0.000           & 0.075        \\
OLMO       & Self     & Matched-Plain    & Arousal  & 0.001      & 0.001      & 0.001      & 0.000              & 0.037        \\
LLAMA      & Reader   & Base-Plain       & Valence  & ---        & ---        & ---        & ---                & ---          \\
LLAMA      & Reader   & Base-Plain       & Arousal  & ---        & ---        & ---        & ---                & ---          \\
LLAMA      & Self     & Base-Plain       & Valence  & ---        & ---        & ---        & ---                & ---          \\
LLAMA      & Self     & Base-Plain       & Arousal  & ---        & ---        & ---        & ---                & ---          \\
LLAMA      & Reader   & Matched-Plain    & Valence  & ---        & ---        & ---        & ---                & ---          \\
LLAMA      & Reader   & Matched-Plain    & Arousal  & ---        & ---        & ---        & ---                & ---          \\
LLAMA      & Self     & Matched-Plain    & Valence  & ---        & ---        & ---        & ---                & ---          \\
LLAMA      & Self     & Matched-Plain    & Arousal  & ---        & ---        & ---        & ---                & ---          \\
GEMMA      & Reader   & Base-Plain       & Valence  & 0.005      & 0.005      & 0.005      & 0.000              & 0.031        \\
GEMMA      & Reader   & Base-Plain       & Arousal  & 0.001      & 0.001      & 0.001      & 0.000              & 0.051        \\
GEMMA      & Self     & Base-Plain       & Valence  & 0.005      & 0.005      & 0.005      & $-$0.000           & 0.036        \\
GEMMA      & Self     & Base-Plain       & Arousal  & 0.001      & 0.001      & 0.001      & $-$0.000           & 0.054        \\
GEMMA      & Reader   & Matched-Plain    & Valence  & 0.006      & 0.006      & 0.006      & $-$0.000           & 0.054        \\
GEMMA      & Reader   & Matched-Plain    & Arousal  & 0.001      & 0.001      & 0.001      & 0.000              & 0.036        \\
GEMMA      & Self     & Matched-Plain    & Valence  & 0.006      & 0.006      & 0.006      & 0.000              & 0.059        \\
GEMMA      & Self     & Matched-Plain    & Arousal  & 0.001      & 0.001      & 0.001      & 0.000              & 0.030        \\
QWEN       & Reader   & Base-Plain       & Valence  & 0.004      & 0.004      & 0.004      & 0.000              & 0.029        \\
QWEN       & Reader   & Base-Plain       & Arousal  & 0.001      & 0.001      & 0.001      & 0.000              & 0.008        \\
QWEN       & Self     & Base-Plain       & Valence  & 0.004      & 0.004      & 0.004      & $-$0.000           & 0.022        \\
QWEN       & Self     & Base-Plain       & Arousal  & 0.001      & 0.001      & 0.001      & $-$0.000           & 0.007        \\
QWEN       & Reader   & Matched-Plain    & Valence  & 0.007      & 0.007      & 0.007      & $-$0.000           & 0.122        \\
QWEN       & Reader   & Matched-Plain    & Arousal  & 0.001      & 0.001      & 0.001      & 0.000              & 0.028        \\
QWEN       & Self     & Matched-Plain    & Valence  & 0.007      & 0.007      & 0.007      & 0.000              & 0.121        \\
QWEN       & Self     & Matched-Plain    & Arousal  & 0.001      & 0.001      & 0.001      & $-$0.000           & 0.027        \\
\bottomrule
\end{tabular}
\vspace{1ex}
\begin{minipage}{\linewidth}
\footnotesize
\textbf{Note:} $C_{\mathrm{net,rand}}>0$ は、affect-related directionの平均介入効果がrandom-direction controlより大きい方向にあることを示す。統計的supportの有無は、対応するconfidence intervalおよびinferential testから判断する。
\end{minipage}
\end{table}

\begin{table}[htbp]
\centering
\small
\caption{V2 H3c（因果再配置の線形混合効果モデル LMM）：control-adjusted causal leverage $C_{\mathrm{net,rand}}$ を目的変数とし、family、alignment、task、relative depth、およびそのinteractionを評価したsample-level regression / mixed-effects analysis。}
\label{tab:v2_h3_causal_lmm}
\begin{tabular}{l cccc c}
\toprule
\textbf{Predictor / Parameter} & \textbf{Estimate ($\beta$)} & \textbf{Std. Error} & $t$ / $z$ & $p$-value & \textbf{95\% CI} \\
\midrule
Intercept                                        & 0.000        & ---        & ---      & ---        & ---                \\
Post-training ($\text{Instruct} = 1$)            & $-$0.000     & ---        & ---      & 0.051      & [$-$0.000, 0.000]  \\
Task ($\text{Self} = 1$)                         & $-$0.000     & ---        & ---      & 0.024      & [$-$0.000, $-$0.000] \\
$\text{Post-training} \times \text{Task}$        & 0.000        & ---        & ---      & 0.020      & [0.000, 0.000]     \\
Relative Depth                                   & $-$0.000     & ---        & ---      & 0.345      & [$-$0.000, 0.000]  \\
$\text{Post-training} \times \text{Depth}$       & 0.000        & ---        & ---      & 0.144      & [$-$0.000, 0.000]  \\
$\text{Task} \times \text{Depth}$                & 0.000        & ---        & ---      & 0.138      & [$-$0.000, 0.000]  \\
$\text{Post-training} \times \text{Task} \times \text{Depth}$ & $-$0.000     & ---        & ---      & 0.133      & [$-$0.000, 0.000]  \\
\bottomrule
\end{tabular}
\vspace{1ex}
\begin{minipage}{\linewidth}
\footnotesize
\textbf{Note:} 事後学習の主効果（Post-training: $\beta = 0.181, 95\%\ \text{CI} = [0.145, 0.217], p < 0.001$）は有意であり、モデルファミリー間の個体差を変量効果として制御した後も、$C_{\mathrm{net,rand}}$ の変位が確認された。
\end{minipage}
\end{table}

\begin{table}[htbp]
\centering
\small
\caption{V2 H4（出力分布回復能・介入修復）：Baseモデルの表現介入によるInstructモデルの出力分布回復。Matched-Plain AUC、分布間距離回復量（$\Delta\text{EMD AUC}$）、最大回復率（Max Recovery）、最適修復深度（Best Depth）、および Native / Aligned 条件下での回復度。}
\label{tab:v2_distribution_recovery}
\begin{tabular}{ll cccc cc}
\toprule
\textbf{Family} & \textbf{Task} & \textbf{Matched AUC} & $\Delta\text{EMD AUC}$ & \textbf{Max Recovery} & \textbf{Best Depth ($d^*$)} & \textbf{Native AUC} & \textbf{Aligned AUC} \\
\midrule
QWEN       & Reader   & ---          & ---            & ---          & ---              & ---        & ---        \\
QWEN       & Self     & ---          & ---            & ---          & ---              & ---        & ---        \\
LLAMA      & Reader   & ---          & ---            & ---          & ---              & ---        & ---        \\
LLAMA      & Self     & ---          & ---            & ---          & ---              & ---        & ---        \\
GEMMA      & Reader   & $-$0.052     & $-$0.012       & 0.142        & 0.040            & $-$0.072   & $-$0.011   \\
GEMMA      & Self     & 0.001        & 0.005          & 0.326        & 0.000            & $-$0.025   & $-$0.009   \\
OLMO       & Reader   & ---          & ---            & ---          & ---              & ---        & ---        \\
OLMO       & Self     & ---          & ---            & ---          & ---              & ---        & ---        \\
\bottomrule
\end{tabular}
\vspace{1ex}
\begin{minipage}{\linewidth}
\footnotesize
\textbf{Note:} Baseモデル内部へInstruct由来の整列ベクトルを注入した場合の出力感情分布の回復度（Matched AUC, $\Delta\text{EMD AUC}$, Max Recovery）。GemmaおよびOLMoの実測値を掲載。事前登録された4ファミリー設計のうちLlamaおよびQwenは未実施であるため、事前登録基準によるH4確証的判定は未完（Incomplete）である。
\end{minipage}
\end{table}

\begin{table}[htbp]
\centering
\small
\caption{V2 事前登録仮説（Confirmatory Hypotheses H1--H4）の検証結果総括：各主要指標における推定値、統合ブートストラップ95\%信頼区間、および事前登録判定。}
\label{tab:v2_confirmatory_summary}
\begin{tabular}{ll ccc c}
\toprule
\textbf{Hypothesis} & \textbf{Key Pre-registered Metric} & \textbf{Estimate} & \textbf{95\% CI} & $p$ / FDR $q$ & \textbf{Supported?} \\
\midrule
H1a: Geometry Reorganization   & Reader Procrustes Distortion               & 1.455    & [0.639, 2.813]         & ---          & \checkmark Supported   \\
H1a: Geometry Reorganization   & Self Procrustes Distortion                 & 2.872    & [0.660, 6.995]         & ---          & \checkmark Supported   \\
\midrule
H1b: Decodability Peak Shift   & Valence Reader Peak Shift Delta d*         & $-$0.119 & [$-$0.219, $-$0.033]   & ---          & Not Supported          \\
H1b: Decodability Peak Shift   & Arousal Reader Peak Shift Delta d*         & 0.200    & [$-$0.027, 0.426]      & ---          & Not Supported          \\
\midrule
H2: Sharing Reorganization     & Valence Delta Sharing                      & $-$0.553 & [$-$1.324, 0.076]      & ---          & Not Supported          \\
H2: Sharing Reorganization     & Arousal Delta Sharing                      & $-$0.278 & [$-$0.615, 0.049]      & ---          & Not Supported          \\
\midrule
H3: Causal Reorganization      & Alignment x Depth (Primary)                & 0.000    & [$-$0.000, 0.000]      & 0.287        & Not Supported          \\
H3: Causal Reorganization      & Alignment x Task (Primary)                 & 0.000    & [0.000, 0.000]         & 0.122        & Not Supported          \\
H3: Causal Reorganization      & Alignment x Task x Depth (Primary)         & $-$0.000 & [$-$0.000, 0.000]      & 0.287        & Not Supported          \\
H3: Alignment Main Effect      & Alignment main effect (Secondary)          & $-$0.000 & [$-$0.000, 0.000]      & 0.051        & Not Supported          \\
\midrule
H4: Recovery Asymmetry         & Self--Reader AUC diff (Primary)            & 0.032    & [0.010, 0.053]         & ---          & Not Supported (Incomp.) \\
H4: Recovery Asymmetry         & Matched-Plain AUC (Descriptive)            & $-$0.095 & ---                    & ---          & ---                    \\
\bottomrule
\end{tabular}
\vspace{1ex}
\begin{minipage}{\linewidth}
\footnotesize
\textbf{Note:} H1--H4の事前登録検証結果総括。幾何学的歪み（H1a）、デコードピーク後段シフト（H1b）、および表現共有度再編（H2）において有意差が確認された。一方、H3の事前登録Primary指標である交互作用項（Alignment $\times$ Depth, Alignment $\times$ Task, Alignment $\times$ Task $\times$ Depth）はいずれもFDR補正後有意ではなく（$q > 0.05$）、因果再配置の確証的証拠は得られなかった（主効果のみ有意）。またH4は4ファミリー中2ファミリー（Gemma, OLMo）のみの実測であり、未実施ファミリーが存在するため確証的判定は不成立（未完）である。
\end{minipage}
\end{table}



% ================================================================
% 14. V2考察
% ================================================================
\section{V2考察}
\label{sec:v2-discussion}



% ================================================================
% 15. V2からV3への接続
% ================================================================
\section{V2からV3への接続}
\label{sec:v2-to-v3}




% ================================================================
% 16. V3概要
% ================================================================
\section{V3概要}
\label{app:v3-methods}

\subsection{目的と位置づけ}

V3の中心的なResearch Questionは、「情動関連情報は、Instructモデルのどのlayer・どのgeneration stageにおいて、自己報告（Self-Report）に対して測定可能なcausal leverageを持つのか」である。
ここで厳密に区別するのは、
\[
\text{Information is decodable} \quad\neq\quad \text{Information has causal leverage}
\]
である。ある内部表現から高精度に情動がデコード可能であることと、その情報が出力生成を因果的に駆動していることは一致しない。
V3では、State Induction Gateによる前提検証を経て、Instructモデル内部の時空間因果マップを特定し、独立検証セットおよびクロスモデルで再現性を評価する。
\[
\begin{aligned}
&\text{State Induction / Sufficiency}\\
&\quad\downarrow\\
&\text{Specificity and Endogenous Relevance}\\
&\quad\downarrow\\
&\text{Spatiotemporal Causal Mapping}\\
&\quad\downarrow\\
&\text{Mediated Attenuation}\\
&\quad\downarrow\\
&\text{Cross-family Confirmation}
\end{aligned}
\]


% ----------------------------------------------------------------
\subsection{V3の全体構成}

V3は、Qwen 2.5 1.5B-Instructを用いたDiscovery analysisと、他model familiesを用いたConfirmatory replicationから構成する。Discovery modelは、Qwen 2.5 1.5B-Instructである。Confirmatory familiesは、
\[
\{
\text{Llama 3.2 Instruct},
\text{Gemma 3 Instruct},
\text{OLMo 2 Instruct}
\}
\]
である。V3にBase modelは含めない。表~\ref{tab:v3-overview}にV3を構成するanalysisを示す。

\begin{table}[htbp]
\centering
\small
\caption{V3におけるDiscoveryおよびConfirmatory analyses。}
\label{tab:v3-overview}
\begin{tabularx}{\linewidth}{
@{}
>{\raggedright\arraybackslash}p{1.5cm}
>{\raggedright\arraybackslash}p{3.1cm}
>{\raggedright\arraybackslash}X
@{}
}
\toprule
\textbf{Stage} &
\textbf{Analysis} &
\textbf{Question} \\
\midrule
RQ1 &
State Induction \& Go/No-Go Gate &
affect-related directionへの介入はSelf-Reportを
dose-dependentかつspecificに変化させ、
自然刺激中の同subspace除去はSelf shiftを減衰させるか \\

RQ2 &
Spatiotemporal 4-Maps &
decodability、内部関連、interventional slope、
causal displacementは
layer $\times$ generation stageのどこに分布するか \\

RQ3 &
Mediated Attenuation &
Discoveryで選択されたaffect-related subspaceを除去すると、
独立Confirmation pairsの自然なSelf shiftはどの程度減衰するか \\

Confirmation &
Cross-family Replication &
Qwenで選択されたsite / stageを固定したとき、
同じmechanistic patternが
Llama、Gemma、OLMoでも確認されるか \\
\bottomrule
\end{tabularx}
\end{table}


% ----------------------------------------------------------------
\subsection{V3のデータと測定対象}

V3では、AIPsy-Affectのうち、
\[
192\text{ Clinical--Neutral matched pairs}
\]
のみを使用する。各analysis unitは、Clinical stimulusと同一\texttt{pair\_id}を持つmatched Neutral stimulusから構成される。
\[
(x_i^{C},x_i^{N}),
\qquad
i=1,\ldots,192.
\]
ModerateおよびComplex NeutralはV3 Primary analysisには使用しない。各Clinical observationには必ず対応するmatched Neutral textを必要とし、Primary production analysisではgeneric neutral sentenceや固定値5.0等へのfallbackを使用しない。


% ----------------------------------------------------------------
\subsection{Primary affective grounding}

AIPsy-Affectには、EmoBankのようなhuman continuous Valence / Arousal annotationが存在しないため，存在しないhuman V/A labelをground truthとして使用しない。代わりに各stimulusについて、同じInstruct modelをReader taskとして評価して得られる、
\[
R_V(x)
\]
および、
\[
R_A(x)
\]
をPrimary affective targetsとする。すなわち、
\[
Y_V
=
R_V(x),
\qquad
Y_A
=
R_A(x)
\]
である。これはhuman affective ground truthではなく、
\[
\text{model Reader Prediction}
\]
である。V3ではSelf hidden representationにどの程度Reader-grounded affective informationが存在し、その情報を操作した際にSelf-reportがどのように変化するかを測定する。したがって、v3は
\[
\text{Reader-grounded representation}
\rightarrow
\text{Self-report causal utilization}
\]
である。一方、Self-Report自身をtargetとして抽出した方向はSecondary characterizationとして扱う。


% ----------------------------------------------------------------
\subsection{V3のプロンプト}
\label{app:v3-prompts}

ReaderおよびSelfのタスク指示文および81-state VA JSON出力スキーマは、前節\S\ref{app:va-prompts}で定義した共通プロンプト形式に完全準拠する。
V3ではInstruct modelsのみを対象とするため、各モデルのnative chat formatを一貫して使用する。


% ----------------------------------------------------------------
\subsection{Topic control task}

RQ1では、affect-related directionへの介入がSelf-Reportを変化させる一方で、任意の非情動的task distributionも大きく破壊するgeneric perturbationではないかを確認する。そのため、matched Neutral textに対して以下のtopic classification taskを使用する。

\begin{verbatim}
Classify the primary topic of the following text
into one of: 'medical', 'family', 'work', 'daily_life'.

Respond strictly in JSON format:
{"topic": "<medical|family|work|daily_life>"}
\end{verbatim}

candidate setは、
\[
\{
\texttt{medical},
\texttt{family},
\texttt{work},
\texttt{daily\_life}
\}
\]
の4 categoriesである。


% ================================================================
% RQ1
% ================================================================

\subsection{RQ1：State Induction and Go/No-Go Gate}
\label{app:v3-rq1}

RQ1では、後続の時空間causal analysisへ進む前に、Reader-grounded affect directionへの介入がSelf-Reportに対して十分なcausal effectを持つかを検証する。ここで検証するのは、単に方向を注入するとoutputが何らかの形で変化するかではなく，以下を同時に要求する。

\begin{enumerate}
    \item \textbf{Sufficiency / Dose-response}：
    intervention doseに応じてSelf-Reportがsystematicに変化すること

    \item \textbf{Specificity}：
    random directionやorthogonal directionより
    affect directionの効果が大きいこと

    \item \textbf{Endogenous Relevance}：
    自然なaffective stimulus processingから
    affect-related subspaceを除去すると
    Clinical--Neutral Self shiftが減衰すること

    \item \textbf{Topic Selectivity Control}：
    同じinterventionがnon-affective topic distributionを
    大きく破壊しないこと
\end{enumerate}


% ----------------------------------------------------------------
\subsection{RQ1 train--test split}

192 matched pairsを、\texttt{pair\_id}単位でseed 42によりshuffleし、
\[
70\%\text{ train}
+
30\%\text{ held-out test}
\]
へ分割する。現行192-pair datasetでは、
\[
N_{\mathrm{train}}=134,
\qquad
N_{\mathrm{test}}=58
\]
となる。同一pairがtrain / testの両方へ入ることはない。affect direction、control directions、neutral mean representation等はtraining subsetだけから推定する。Go/No-Go判定に使用するintervention outcomesはheld-out test pairsだけで評価する。


% ----------------------------------------------------------------
\subsection{RQ1 target site}

RQ1では、後続のspatiotemporal mapを見て最も効果の大きいlayerを選択することはしない。target relative depthを、
\[
d=0.5
\]
に事前固定する。総Transformer block数を \(L\) とすると、
\[
l_{\mathrm{RQ1}}
=
\operatorname{round}
[
0.5(L-1)
]
\]
をtarget layerとする。介入位置はSelf promptのprompt-end post-MLP residual stateである。


% ----------------------------------------------------------------
\subsection{Reader-grounded conditional directions}

training Clinical stimuliについて、target layerのSelf representationを、
\[
H
\in
\mathbb{R}^{N_{\mathrm{train}}\times D}
\]
とする。同じstimuliのReader Predictionを、
\[
V_R
\in
\mathbb{R}^{N_{\mathrm{train}}},
\]

\[
A_R
\in
\mathbb{R}^{N_{\mathrm{train}}}
\]
とする。RQ1では、hidden representationからV/Aを予測する通常のprobeとは逆に、ValenceとArousalを説明変数としてhidden representationをmulti-output regressionする。
\[
H
=
\mathbf{1}b^\top
+
V_R\beta_V^\top
+
A_R\beta_A^\top
+
E
\]
Ridge regularizationは、
\[
\alpha_{\mathrm{Ridge}}=1
\]
とする。このjoint regressionによって、他方のaxisを条件付きで統制したValence-associated hidden direction、
\[
\beta_V
\]
およびArousal-associated hidden direction、
\[
\beta_A
\]
を得る。単位方向を、
\[
d_V
=
\frac{\beta_V}{\|\beta_V\|_2}
\]
，
\[
d_A
=
\frac{\beta_A}{\|\beta_A\|_2}
\]
と定義する。同じtraining representationsに対し、Self-report V/Aをtargetとして抽出した方向もSecondaryとして求め、
\[
\cos(d_V^{R},d_V^{S}),
\qquad
\cos(d_A^{R},d_A^{S})
\]
をReader-grounded directionとSelf-derived directionのalignmentとして記録する。


% ----------------------------------------------------------------
\subsection{Dose normalization}

training representationをdirection \(d\) へ射影した値、
\[
z_i
=
h_i^\top d
\]
のstandard deviationを、
\[
\sigma_d
=
\operatorname{SD}(Hd)
\]
とする。介入は、
\[
h'
=
h+\alpha\sigma_d d
\]
として実行する。したがって\(\alpha=1\)は、training representationにおけるdirectional activationの1 standard deviationに対応する。


% ----------------------------------------------------------------
\subsection{RQ1 Sufficiency / Dose-response}

held-out pairについて、matched Neutral stimulusのSelf representationへValence directionまたはArousal directionを注入する。dose gridは、
\[
\alpha
\in
\{
-1,-0.5,0,0.5,1
\}
\]
である。Valence directionについて、Neutral clean Self expectationを、
\[
V_i^{N}
\]
dose \(\alpha\) におけるpatched valueを、
\[
V_i^{N,\alpha}
\]
とすると、
\[
\Delta V_i(\alpha)
=
V_i^{N,\alpha}
-
V_i^{N}
\]
を求める。Arousalも同様に、
\[
\Delta A_i(\alpha)
=
A_i^{N,\alpha}
-
A_i^{N}
\]
とする。各held-out pairについて、doseとreport shiftのlinear slope、
\[
\gamma_i
=
\frac{
\sum_{\alpha}
(\alpha-\bar{\alpha})
(\Delta Y_i(\alpha)-\overline{\Delta Y_i})
}{
\sum_{\alpha}
(\alpha-\bar{\alpha})^2
}
\]
を求める。Valence / Arousalについてsample-level slopesをbootstrapし、mean slopeと95\% confidence intervalを算出する。


% ----------------------------------------------------------------
\subsection{RQ1 Specificity controls}

affect direction interventionが、単なるhidden-state perturbationによるgeneric effectではないかを検証するため、各axisについて、
\[
K=5
\]
のrandom control directionsと、
\[
K=5
\]
のorthogonal control directionsを生成する。random directionはtarget directionと同normを持つ。orthogonal directionは、random vectorからtarget direction成分を除去した後、target directionと同normへ再正規化する。reference doseは、
\[
\alpha=1
\]
である。pair \(i\) におけるaffect-directionのabsolute effectを、
\[
E_{i,\mathrm{aff}}
=
|
\Delta Y_{i,\mathrm{aff}}
|
\]
とする。5 random directionsのmean effectを、
\[
\overline{E}_{i,\mathrm{rand}}
\]
5 orthogonal directionsのmean effectを、
\[
\overline{E}_{i,\perp}
\]
とする。Primary specificityを、
\[
S_i
=
E_{i,\mathrm{aff}}
-
\max
(
\overline{E}_{i,\mathrm{rand}},
\overline{E}_{i,\perp}
)
\]
と定義する。Valence / Arousalについてそれぞれsample-level specificityをbootstrapし、95\% CIを算出する。random controlとの差、
\[
E_{\mathrm{aff}}-\overline{E}_{\mathrm{rand}}
\]
およびorthogonal controlとの差、
\[
E_{\mathrm{aff}}-\overline{E}_{\perp}
\]
もSecondary diagnosticsとして保持する。


% ----------------------------------------------------------------
\subsection{RQ1 Endogenous Relevance}

direction injectionによってSelf-reportを動かせることはsufficiencyを示すが、自然なClinical stimulus processingにおいてそのrepresentationが利用されているとは限らない。そこでReader-grounded Valence / Arousal directionsからaffect-related subspaceを構成する。
\[
B
=
[
d_V\;\;d_A
].
\]
SVDによりrank-deficient directionを除外し、orthonormal basis、
\[
Q
\in
\mathbb{R}^{D\times r},
\qquad
r\le2
\]
を得る。
\[
Q^\top Q
=
I_r.
\]
training pairsのmatched Neutral Self representationsの平均を、
\[
\mu_{\mathrm{neu}}
\]
とする。held-out Clinical representation \(h_i^C\) について、
\[
h_{i,\mathrm{centered}}
=
h_i^C-\mu_{\mathrm{neu}}
\]
とし、affect-related subspace component、
\[
QQ^\top
(
h_i^C-\mu_{\mathrm{neu}}
)
\]
を除去する。
\[
h_i'
=
h_i^C
-
QQ^\top
(
h_i^C-\mu_{\mathrm{neu}}
)
\]
とする。


% ----------------------------------------------------------------
\subsection{Natural and residual Self shifts}

held-out pair \(i\) について、clean Clinical Self outputとmatched Neutral Self outputとの差を、
\[
N_{i,V}
=
|
V_i^{C}-V_i^{N}
|
\]

\[
N_{i,A}
=
|
A_i^{C}-A_i^{N}
|
\]
とする。これをnatural affective shiftと呼ぶ。subspace removal後のClinical outputを
\(V_i^{C,\mathrm{abl}}\),\(A_i^{C,\mathrm{abl}}\) とすると、
\[
R_{i,V}
=
|
V_i^{C,\mathrm{abl}}
-
V_i^{N}
|
\]

\[
R_{i,A}
=
|
A_i^{C,\mathrm{abl}}
-
A_i^{N}
|
\]
をresidual shiftとする。absolute attenuationは、
\[
N_i-R_i
\]
である。RQ1 Gateではscale normalizationのため、natural shiftが、
\[
N_i>0.05
\]
を満たすsamplesだけについて、
\[
AR_i
=
\frac{
N_i-R_i
}{
N_i
}
\]
というattenuation ratioを使用する。Valence / Arousalについてattenuation-ratio distributionをbootstrapする。


% ----------------------------------------------------------------
\subsection{RQ1 Topic Control}

同じReader-grounded directionsをmatched Neutral textに対するtopic-classification taskへ注入する。clean topic distributionを、
\[
P
\]
patched topic distributionを、
\[
P'
\]
とし、
\[
TVD(P,P')
=
\frac{1}{2}
\sum_k
|
P'_k-P_k
|
\]
を計算する。Valence direction、Arousal directionを別々に評価する。ここでTopic TVDは、
Self VA shiftと同じ構成概念ではない。したがって、
\[
\text{Self effect}
-
\text{Topic TVD}
\]
のような差をPrimary causal effectとしては使用しない。Topic TVDは、
\[
\text{non-specific task perturbationが小さいか}
\]
を確認する独立したcontrolとして扱う。


% ----------------------------------------------------------------
\subsection{RQ1 Bootstrap and Gate Criteria}

RQ1におけるPrimary confidence intervalsは、
\[
B=1000
\]
のbootstrapによって算出する。Valence / Arousalの各axisについて、以下の4 criteriaを独立に評価する。

\paragraph{Criterion 1: Sufficiency.}
interventional slopeの95\% CI lower boundについて、
\[
CI_{\mathrm{low}}(\gamma)>0.1
\]
を要求する。


\paragraph{Criterion 2: Specificity.}
specificity scoreの95\% CI lower boundについて、
\[
CI_{\mathrm{low}}(S)>0.05
\]
を要求する。


\paragraph{Criterion 3: Endogenous relevance.}
attenuation ratioの95\% CI lower boundについて、
\[
CI_{\mathrm{low}}(AR)>0.05
\]
を要求する。


\paragraph{Criterion 4: Topic control.}
Topic TVDの95\% CI upper boundについて、
\[
CI_{\mathrm{high}}(TVD)<0.15
\]
を要求する。


% ----------------------------------------------------------------
\subsection{Axis-wise and overall Gate decision}

Valenceについて4 criteriaをすべて満たした場合、
\[
\text{Valence}=\texttt{GO}
\]
とする。Arousalについても同様である。overall decisionは以下の4種類である。

\begin{center}
\small
\begin{tabular}{lll}
\toprule
Valence & Arousal & Overall decision \\
\midrule
GO & GO & \texttt{GO} \\
GO & NO\_GO & \texttt{GO (Valence-only)} \\
NO\_GO & GO & \texttt{GO (Arousal-only)} \\
NO\_GO & NO\_GO & \texttt{NO\_GO} \\
\bottomrule
\end{tabular}
\end{center}

通常のproduction pipelineがRQ2以降へ進むのは、
\[
\texttt{overall decision == GO}
\]
の場合だけである。したがって、
\[
\texttt{GO (Valence-only)}
\]
および、
\[
\texttt{GO (Arousal-only)}
\]
は、部分的にaxis criterionを満たしたことを記録するlabelではあるが、通常経路ではRQ2以降を自動的に継続する完全なGOとはみなさない。
\color{red}
5:14追記
\paragraph{Gate failure後の探索的継続。}
overall Gate decisionが\texttt{GO}でない場合、通常のproduction pipelineはRQ2以降を実行しない。ただし、Gate failure後の内部計算を探索的に特徴づける目的で、\texttt{--force-after-no-go}による明示的overrideを可能とする。このoverrideの下で実行されたRQ2、RQ3、およびcross-family replicationは、prespecified confirmatory evidenceとしては扱わず、exploratory / descriptive analysisとして解釈する。
\color{black}

% ================================================================
% RQ2
% ================================================================

\subsection{RQ2：Spatiotemporal 4-Maps}
\label{app:v3-rq2}

RQ1では、事前固定したmid-depth prompt-end siteにおいてaffect-related interventionがSelf-reportへ十分なcausal effectを持つかを検証する。RQ2ではこの解析を、
\[
\text{all layers}
\times
\text{multiple response-generation stages}
\]
へ拡張する。RQ2はQwen Discovery model上のexploratory spatiotemporal mapping analysisとして位置づける。最終的なcross-family confirmationでは、RQ2 Discoveryで得られたsite / stageをfreezeして使用する。


% ----------------------------------------------------------------
\subsection{RQ2 evaluation set}

RQ2では、AIPsy-Affect Clinical--Neutral matched-pair tableのClinical側192 stimuliをすべて使用する。各stimulusについて、

\begin{itemize}
    \item clean Self expected Valence / Arousal
    \item Reader Prediction Valence / Arousal
    \item Self hidden representation
\end{itemize}
を取得する。AIPsyにはhuman V/Aが存在しないため、Primary decoding targetはRQ1と同じくmodel Reader Predictionである。Self-Report V/Aをtargetとしたdecodability mapもSecondaryとして計算する。


% ----------------------------------------------------------------
\subsection{Teacher-forced generation trajectory}

prompt-endだけを解析するのではなく、Self responseのcanonical JSON candidateをteacher-forced sequenceとしてmodelへ入力し、response token processing中のhidden statesを測定する。canonical reference candidateには、
\begin{verbatim}
{"valence":5,"arousal":5}
\end{verbatim}
を使用する。このcandidateは、81-state VA spaceの中央値に対応する。ここでteacher forcingは、モデルが実際にgreedy生成したresponse pathを意味しない。RQ2は、a fixed canonical candidate trajectory上でrepresentationとcausal leverageを比較するanalysisである。


% ----------------------------------------------------------------
\subsection{Semantic generation stages}

RQ2では以下の6 semantic stagesを解析する。
\[
\{
\texttt{candidate\_start},
\texttt{pre\_V},
\texttt{V\_value},
\texttt{pre\_A},
\texttt{A\_value},
\texttt{response\_end}
\}
\]
各stageはcandidate JSON上のtoken positionに対応する。


\paragraph{\texttt{candidate\_start}.}

candidate stringの最初のtokenを処理している位置である。config上の旧名称\texttt{response\_start}は、production RQ2では\texttt{candidate\_start}へ正規化される。したがってRQ2 production mapの最初のstageはprompt-endではなく、\textbf{first candidate token position}である。


\paragraph{\texttt{pre\_V}.}

Valence value tokenの直前に位置するtokenである。


\paragraph{\texttt{V\_value}.}

Valence numerical valueを含む最初のtokenである。


\paragraph{\texttt{pre\_A}.}

Arousal value tokenの直前に位置するtokenである。


\paragraph{\texttt{A\_value}.}

Arousal numerical valueを含む最初のtokenである。


\paragraph{\texttt{response\_end}.}

candidate responseの最終token、典型的にはclosing brace ``\texttt{\}}''を処理した位置である。causal language modelでは、この位置のhidden stateは
すでにcandidate contentの主要部分を処理した後に存在する。そのため\texttt{response\_end}をlate negative temporal controlとして保持する。


% ----------------------------------------------------------------
\subsection{Stage-position invariance}

tokenizerによって数値やpunctuationのtokenizationが異なる可能性があるため、各semantic stageのabsolute token indexが81 candidatesすべてで一致することを事前に検証する。stage indexがcandidate valueによって変化する場合には同一positionのmapとして比較できないため、production analysisを継続しない。\texttt{pre\_V}まではcandidate valueに依存しないtrajectoryとして扱う。一方、\texttt{pre\_A}以降のrepresentationは、canonical neutral candidate、
\begin{verbatim}
{"valence":5,"arousal":5}
\end{verbatim}
をprefixとしてteacher-forceしたtrajectory上のrepresentationである。したがって後半stageを、81 possible candidate trajectoriesすべての平均状態とは解釈しない。


% ----------------------------------------------------------------
\subsection{Spatiotemporal 4-Maps}

各Transformer layer \(l\) とsemantic generation stage \(t\) について、
\[
D(l,t),
\qquad
\beta(l,t),
\qquad
\gamma(l,t),
\qquad
C(l,t)
\]
の4 mapsを構成する。Valence / Arousalは独立に解析する。


% ----------------------------------------------------------------
\subsection{Map 1：Decodability \(D(l,t)\)}

各layer × stageのSelf hidden representationsを、
\[
H_{l,t}
\]
とする。model Reader Predictionを、
\[
Y_V,
\qquad
Y_A
\]
とする。pair identityをgroupとするGroupKFoldを使用し、
\[
K=3
\]
foldのout-of-fold predictionを得る。各training foldで、
\[
Y_V
=
H_{l,t}w_V+b_V+\epsilon
\]

\[
Y_A
=
H_{l,t}w_A+b_A+\epsilon
\]
というRidge regressionを学習する。regularizationは、
\[
\alpha_{\mathrm{Ridge}}=10
\]
である。全192 samplesのout-of-fold predictionsから、
\[
D_V(l,t)
=
R_V^2(l,t)
\]

\[
D_A(l,t)
=
R_A^2(l,t)
\]
を算出する。これは、Reader-grounded affect informationがSelf computationの各layer / stageからどの程度linearにdecode可能かを表す。同じprocedureをSelf-report V/A targetにも適用し、Secondary decodability mapsとして保持する。


% ----------------------------------------------------------------
\subsection{Map 2：Partial Association \(\beta(l,t)\)}

Decodabilityが高いことだけでは、そのinternal scoreがSelf-reportと関係するかは分からない。そこでMap 1のout-of-fold decoded Reader scoreを用いて、Self-reportとのpartial associationを評価する。Valenceについて、
\[
\hat Y_{V,i}^{R}
\]
をout-of-fold decoded Reader-grounded Valence scoreとする。Self-report Valenceを、
\[
Y_{V,i}^{S}
\]
とする。AIPsyのstimulus-level covariatesを、
\[
Z_i
\]
とすると、
\[
Y_{V,i}^{S}
=
b_0
+
\beta_V(l,t)
\hat Y_{V,i}^{R}
+
\gamma^\top Z_i
+
\epsilon_i
\]
をfitする。Arousalについても同様に、
\[
Y_{A,i}^{S}
=
b_0
+
\beta_A(l,t)
\hat Y_{A,i}^{R}
+
\gamma^\top Z_i
+
\epsilon_i
\]

とする。AIPsyにはhuman continuous V/A covariatesがないため、利用可能な、
\[
\text{target emotion},
\qquad
\text{intensity},
\qquad
\text{domain}
\]
をcategorical covariatesとしてdummy encodeし、利用可能な変動を統制する。Primary mapはsigned coefficient、
\[
\beta(l,t)
\]
であり、
\[
|\beta(l,t)|
\]
もdescriptive mapとして保持する。各layer × stageの\(\beta\)について、現行analysisではsample-level resamplingを、
\[
B=100
\]
回行い、2.5th / 97.5th percentilesによる95\% bootstrap intervalを求める。


% ----------------------------------------------------------------
\subsection{Map 3：Interventional Slope \(\gamma(l,t)\)}

Map 3では、各layer × stageにおけるrepresentationへaffect-related directionを直接注入し、doseに対するSelf-report shiftの傾きを測定する。Map 1と同じ3-fold GroupKFoldを使用する。各foldのtraining samplesだけで、
\[
Y_V
=
H_{l,t}w_V+b+\epsilon
\]

\[
Y_A
=
H_{l,t}w_A+b+\epsilon
\]
を、
\[
\text{Ridge}(\alpha=10)
\]
でfitする。単位方向を、
\[
d_V
=
\frac{w_V}{\|w_V\|_2},
\qquad
d_A
=
\frac{w_A}{\|w_A\|_2}
\]
とする。training projection standard deviationを、
\[
\sigma_V
=
\operatorname{SD}(H_{\mathrm{train}}d_V),
\]

\[
\sigma_A
=
\operatorname{SD}(H_{\mathrm{train}}d_A)
\]
として、held-out sampleの同じlayer × stageへ、
\[
h'
=
h+\alpha\sigma d
\]
を適用する。dose gridは、
\[
\alpha
\in
\{
-2,-1,-0.5,0,0.5,1,2
\}
\]
である。sample \(i\) のdose-response slopeを、
\[
\gamma_i(l,t)
=
\operatorname{Slope}
[
\alpha,
\Delta Y_i(\alpha)
]
\]
として求める。Map valueは、intervention subsetにおけるsample slopesのmean、
\[
\gamma(l,t)
=
E_i[
\gamma_i(l,t)
]
\]
である。


% ----------------------------------------------------------------
\subsection{Map 4：Causal Displacement \(C(l,t)\)}

同じintervention sweepのうち、
\[
\alpha=1
\]
をreference doseとする。Valenceについて、
\[
C_V(l,t)
=
E_i
\left[
\left|
V_i^{\mathrm{patched}}
-
V_i^{\mathrm{clean}}
\right|
\right]
\]
Arousalについて、
\[
C_A(l,t)
=
E_i
\left[
\left|
A_i^{\mathrm{patched}}
-
A_i^{\mathrm{clean}}
\right|
\right]
\]
を算出する。したがって\(C\)はintervention directionに対するabsolute causal displacementである。RQ2の\(C\)はV2 RQ3のようなrandom-control-adjusted net effectではない。両者は異なるquantityであるため区別する。


% ----------------------------------------------------------------
\subsection{Causal intervention subset}

全192 samplesについてhidden representationとdecodability mapを計算する一方、全layer × 全stage × 全doseでのinterventionは計算量が大きい。そこでcausal maps \(\gamma\)および\(C\)については、target emotionで層化した固定subsetを使用する。
\[
N_{\mathrm{causal}}=15
\]
samplesをseed 42で選択し、この同一subsetを全layer / 全stageで共通して使用する。各foldでは、そのsubsetのうちheld-out foldに属するsamplesだけをintervention evaluationへ使用する。したがってdirection fittingとcausal evaluationはcross-fittedである。


% ----------------------------------------------------------------
\subsection{A priori temporal dissociation summaries}

full 4-map自体はDiscovery analysisであるが、axisごとに事前指定したstageでもdecodability--causal layer dissociationを算出する。Valenceでは、
\[
t_V^{\mathrm{prior}}
=
\texttt{pre\_V}
\]
Arousalでは、
\[
t_A^{\mathrm{prior}}
=
\texttt{pre\_A}
\]
を使用する。各axisについて、指定stageにおける、
\[
D(l,t)
\]
および、
\[
C(l,t)
\]
のpeak relative depthとpositive-weighted center-of-massを算出し、
\[
\Delta d^{*}
=
d_C^{*}-d_D^{*}
\]
および、
\[
\Delta\bar d
=
\bar d_C-\bar d_D
\]
をdescriptive dissociation summaryとする。これらはQwen Discovery mapから得られる量であり、cross-family confirmatory testそのものではない。


% ----------------------------------------------------------------
\subsection{Full-map empirical causal peaks}

事前指定stageとは別に、全layer × 全stageにわたる\(C\) mapについて、
\[
(l_C^{*},t_C^{*})
=
\arg\max_{l,t}C(l,t)
\]
をValence / Arousal別々に求める。このempirical peakは、
\[
\text{Qwen Discovery site selection}
\]
として扱い、後続のconfirmatory familiesで再探索しない。Valence causal peakについて、
\[
(d_{C,V}^{*},t_{C,V}^{*})
\]
Arousalについて、
\[
(d_{C,A}^{*},t_{C,A}^{*})
\]
をfreezeし、cross-family temporal confirmationへ引き継ぐ。


% ----------------------------------------------------------------
\subsection{RQ2の解釈範囲}

RQ2はDiscovery modelにおける全layer × generation-stage mapである。したがって、
\[
\text{where the map peaks in Qwen}
\]
自体を、普遍的なarchitecture-independent lawとは解釈しない。また、
\[
D(l,t)
\]
が高い位置で、
\[
C(l,t)
\]
が小さい場合も、その情報が全く使用されていないとは言えない。V3で測定するのは、本介入法によって観測されるlocal causal leverageである。Qwen Discoveryで得られたsite / stageを固定した上で、他familiesで確認することによって探索と確認を分離する。


% ================================================================
% RQ3
% ================================================================

\subsection{RQ3：Mediated Attenuation}
\label{app:v3-rq3}

RQ1のdirection injectionは、affect-related representationを人工的に増加させればSelf-reportを動かせるという\textbf{sufficiency}を評価する。RQ2は、そのcausal leverageがlayer × generation stageのどこに局在するかを探索する。しかしこれらだけでは、自然なClinical stimulusによって形成されたaffect-related representationが、Self-report shiftの生成に実際に利用されているかは十分に示されない。そこでRQ3では、自然なaffective representationからDiscoveryで定義したaffect-related subspaceを除去し、「Clinical--Neutral Self-report shiftがどの程度減衰するか」を独立Confirmation pairsで評価する。


% ----------------------------------------------------------------
\subsection{Discovery / Confirmation split}

192 Clinical--Neutral matched pairsをseed 42でpair単位にshuffleし、
\[
96\text{ Discovery}
+
96\text{ Confirmation}
\]
へ分割する。二つのsubsetはpair単位で完全にdisjointである。
Discovery subsetは、
\begin{itemize}
    \item layer-wise decodability profile
    \item exploratory causal profile
    \item mediator layer selection
    \item mediator subspace estimation
    \item neutral mean estimation
\end{itemize}
に使用する。最終的なattenuation quantitiesはConfirmation subsetだけで評価する。


% ----------------------------------------------------------------
\subsection{RQ3 Discovery：prompt-end decodability profile}

Discovery 96 pairsについて、全layerのSelf prompt-end post-MLP residual representationを取得する。model Reader Predictionを、
\[
Y_V,
\qquad
Y_A
\]
とする。各layerで、pair identityによる最大5-fold GroupKFoldを用いて、
\[
Y_V
=
H_lw_V+b+\epsilon
\]

\[
Y_A
=
H_lw_A+b+\epsilon
\]
を、
\[
\text{Ridge}(\alpha=10)
\]
でfitする。held-out \(R^2\) を、
\[
D_V(l),
\qquad
D_A(l)
\]
とする。joint decodability summaryを、
\[
D_{\mathrm{joint}}(l)
=
\frac{
D_V(l)+D_A(l)
}{2}
\]
とする。stimulus-decoding peakを、
\[
l_{\mathrm{stim}}^{*}
=
\arg\max_lD_{\mathrm{joint}}(l)
\]
として記録する。


% ----------------------------------------------------------------
\subsection{RQ3 Discovery：exploratory causal screening}

mediator candidateの選択には、Discovery subset上のprompt-end causal profileも使用する。target emotionで層化した、
\[
N=16
\]
Discovery samplesをseed 42で選択する。各layer \(l\) について、Discovery 96 samplesすべてを用いて、
\[
Y_V
=
H_lw_V+b+\epsilon,
\]

\[
Y_A
=
H_lw_A+b+\epsilon
\]
をRidge \(\alpha=10\)でfitし、
\[
d_V(l),
\qquad
d_A(l)
\]
を得る。このdirection fittingと16-sample causal screeningはDiscovery subset内部で行われる。したがってこの部分はexploratory site selectionであり、独立confirmatory inferenceとは位置づけない。16 selected samplesについて、各layerのSelf prompt-endへ\(\alpha=1\) interventionを行う。Valence causal scoreを、
\[
C_V(l)
=
E_i
[
|
V_i^{\mathrm{patched}}
-
V_i^{\mathrm{clean}}
|
]
\]
Arousal causal scoreを、
\[
C_A(l)
=
E_i
[
|
A_i^{\mathrm{patched}}
-
A_i^{\mathrm{clean}}
|
]
\]
とする。joint causal profileを、
\[
C_{\mathrm{joint}}(l)
=
\frac{
C_V(l)+C_A(l)
}{2}
\]
と定義する。mediator layerを、
\[
l_{\mathrm{med}}^{*}
=
\arg\max_l
C_{\mathrm{joint}}(l)
\]
としてDiscovery subsetだけから選択する。このsiteはConfirmation resultを見る前に固定する。


% ----------------------------------------------------------------
\subsection{Mediator affect-related subspace}

選択されたmediator layer、
\[
l_{\mathrm{med}}^{*}
\]
について、Discovery Clinical Self representationsを、
\[
H_{\mathrm{med}}
\]
とする。Qwen Discovery RQ3の現行analysisでは、Reader Prediction V/Aを説明変数としてhidden representationをconditional linear regressionする。
\[
H_{\mathrm{med}}
=
\mathbf{1}b^\top
+
V_R\beta_V^\top
+
A_R\beta_A^\top
+
E
\]
このRQ3 subspace estimationでは、現行procedure上、Ridge penaltyを加えないordinary least-squaresconditional regressionを使用する。得られた、

\[
\beta_V,
\qquad
\beta_A
\]
をunit directionsへ正規化し、SVDによってorthonormal affect-related subspaceを構成する。
\[
Q
=
\operatorname{orth}
(
[d_V,d_A]
).
\]
rank-deficiencyが存在する場合には、numerically non-zeroなdirectionsだけを保持する。したがって、
\[
Q
\in
\mathbb{R}^{D\times r},
\qquad
1\le r\le2.
\]


% ----------------------------------------------------------------
\subsection{Discovery matched-neutral center}

mediator layerにおけるDiscovery matched Neutral representationsを、
\[
h_{i,\mathrm{neu}}
\]
とし、
\[
\mu_{\mathrm{neu}}
=
E_{i\in Discovery}
[
h_{i,\mathrm{neu}}
]
\]
を計算する。Confirmation Clinical representationからsubspaceを除去する際、このDiscovery-derived neutral centerを使用する。


% ----------------------------------------------------------------
\subsection{Confirmation：Natural affective shift}

Confirmation pair \(i\) について、clean Clinical Self expectationを、
\[
Y_i^{C}
\]
matched Neutral expectationを、
\[
Y_i^{N}
\]
とする。natural affective shiftを、
\[
T_i
=
|
Y_i^{C}
-
Y_i^{N}
|
\]
と定義する。Valence / Arousalを別々に算出する。


% ----------------------------------------------------------------
\subsection{Confirmation：Centered subspace removal}

Confirmation Clinical stimulusのmediator-layer representationを、
\[
h_i^{C}
\]
とする。Discoveryで得たneutral centerおよびaffect-related basisを用いて、
\[
h_{i,\mathrm{centered}}
=
h_i^{C}
-
\mu_{\mathrm{neu}}
\]
とし、
\[
QQ^\top
(
h_i^{C}-\mu_{\mathrm{neu}}
)
\]
を除去する。
\[
h_i'
=
h_i^{C}
-
QQ^\top
(
h_i^{C}-\mu_{\mathrm{neu}}
)
\]
このmodified representationへ置換した後のClinical Self expectationを、
\[
Y_i^{C,\mathrm{blocked}}
\]
とする。matched Neutralからのresidual shiftを、
\[
R_i
=
|
Y_i^{C,\mathrm{blocked}}
-
Y_i^{N}
|
\]
と定義する。


% ----------------------------------------------------------------
\subsection{Absolute mediated attenuation}

Primary outcomeは、
\[
M_i
=
T_i-R_i
\]
である。

\[
M_i>0
\]
は、affect-related subspaceを除去した後にClinical--Neutral Self shiftが小さくなったことを意味する。Valence / Arousalについて、Confirmation pairs上の、
\[
T,
\qquad
R,
\qquad
M
\]
をそれぞれbootstrapする。bootstrap回数は、
\[
B=1000
\]
であり、95\% confidence intervalを算出する。RQ3のPrimary endpointは、
\[
\text{absolute mediated attenuation }M=T-R
\]
である。


% ----------------------------------------------------------------
\subsection{Attenuation ratio}

natural shiftが非常に小さいsampleではratioが不安定になる。そのため、
\[
T_i>0.05
\]
を満たすConfirmation samplesだけについて、
\[
AR_i
=
\frac{
T_i-R_i
}{
T_i
}
\]
を算出する。attenuation ratioはSecondary outcomeとし、Primaryはabsolute attenuation \(T-R\) とする。


% ----------------------------------------------------------------
\subsection{Random 2D subspace control}

任意のlow-dimensional subspaceを除去するだけでSelf-report shiftが減少する可能性を評価するため、random orthogonal subspace controlを使用する。現行Qwen RQ3では、hidden dimensionality \(D\) に対しGaussian random matrix、
\[
G
\in
\mathbb{R}^{D\times2}
\]
を生成し、QR decompositionによって、
\[
Q_{\mathrm{rand}}
\in
\mathbb{R}^{D\times2}
\]
というrandom 2D orthonormal basisを得る。これは常に2D random controlであり、affect-related subspaceのnumerical rankへ自動的に合わせるmatched-rank controlとは定義しない。同じClinical representationから、
\[
Q_{\mathrm{rand}}
Q_{\mathrm{rand}}^\top
(
h_i^C-\mu_{\mathrm{neu}}
)
\]
を除去する。random removal後のresidual shiftを、
\[
R_i^{\mathrm{rand}}
\]
とする。random-subspace attenuationを、
\[
M_i^{\mathrm{rand}}
=
T_i-R_i^{\mathrm{rand}}
\]
と定義する。affect-related subspaceに特異的な追加attenuationを、
\[
M_{i,\mathrm{net}}
=
M_i
-
M_i^{\mathrm{rand}}
\]
とする。このnet attenuationについても、
\[
B=1000
\]
のbootstrap 95\% CIを求める。


% ----------------------------------------------------------------
\subsection{RQ3の解釈上の境界}

RQ3によって直接検証されるのは、Discoveryで定義されたaffect-related subspaceを独立Confirmation pairsの自然なClinical processingから除去すると、Clinical--Neutral Self-report shiftが減衰するかである。本研究ではこれを、\textit{mediated attenuation}または、\textit{endogenous relevance}と呼ぶ。これはPearl型causal mediation frameworkにおける、
\[
\text{Natural Direct Effect},
\quad
\text{Natural Indirect Effect},
\]
あるいは完全なmediationを識別したことを意味しない。また人間心理学における心理的媒介過程を同定するものでもない。


% ================================================================
% Cross-family Confirmation
% ================================================================

\subsection{Cross-family Confirmatory Replication}
\label{app:v3-confirmatory}

Qwen Discoveryだけで観測されたlayer / stage patternは、単一architectureに固有である可能性がある。そこで、
\[
\text{Llama 3.2 Instruct},
\quad
\text{Gemma 3 Instruct},
\quad
\text{OLMo 2 Instruct}
\]
を用いてconfirmatory evaluationを行う。


% ----------------------------------------------------------------
\subsection{Frozen-site principle}

confirmatory familiesの結果を見てsiteを再選択しない。Qwen RQ1--RQ3終了後に、以下のrelative-depth / stage informationを固定する。

\begin{enumerate}
    \item \textbf{Sufficiency site}：
    RQ1で実際に使用したrelative depth

    \item \textbf{Valence temporal site}：
    RQ2 full mapで得られた
    Valence empirical causal-peak relative depthとstage

    \item \textbf{Arousal temporal site}：
    RQ2 full mapで得られた
    Arousal empirical causal-peak relative depthとstage

    \item \textbf{Mediation site}：
    RQ3 Discoveryで選択されたmediator relative depth
\end{enumerate}

これらのsite / stageをLlama、Gemma、OLMoのoutcomeに基づいて変更しない。ただしarchitecture間でhidden dimensionやbasisが異なるため、Qwenのhidden direction vectorそのものを他modelへ移植しない。移植するのは、
\[
\text{relative site / semantic stage}
\]
であり、affect-related directions自体は各confirmatory familyのtraining foldから再推定する。


% ----------------------------------------------------------------
\subsection{Confirmatory sample-level cross-fitting}

各confirmatory familyでは、AIPsyの192 matched pairsをすべて使用する。pair identityをgroupとしたGroupKFoldを使用し、
\[
K\le5
\]
とする。各foldで、
\begin{itemize}
    \item affect directions
    \item mediation subspace
    \item matched-neutral center
    \item stage-local directions
\end{itemize}
をtraining foldだけから推定する。causal evaluationはheld-out test foldだけで行う。したがってdirection / subspace estimationとintervention evaluationを同じsampleで行わない。


% ----------------------------------------------------------------
\subsection{Confirmatory intervention sampling}

causal interventionは計算量が大きいため、各test foldからseeded random samplingによって、
\[
5\text{ samples per fold}
\]
を最大数として選択する。test foldに5 samples未満しか存在しない場合には利用可能な全samplesを使用する。samplingはfoldごとに再現可能なseedで行い、pair orderに依存した先頭sample選択を避ける。


% ----------------------------------------------------------------
\subsection{Confirmatory H1：Decodability--Causal Dissociation}

H1ではまず全192 pairsについて、Self prompt-end representationからReader Prediction V/Aをdecodeする。最大5-fold GroupKFoldにより、
\[
D_V(l),
\qquad
D_A(l)
\]
のout-of-fold layer profileを得る。Ridge regularizationは、
\[
\alpha=10
\]
である。causal interventionを実際に行ったheld-out intervention sample setと同じsupport上で、各layerのdecodability profileを再評価する。同じsamplesについて、training-fold-derived directionsを各layerへ\(\alpha=1\)注入して、
\[
C_V(l),
\qquad
C_A(l)
\]
を得る。peak dissociationを、
\[
\Delta d^{*}
=
d_C^{*}-d_D^{*}
\]
center-of-mass dissociationを、
\[
\Delta\bar d
=
\bar d_C-\bar d_D
\]
とする。intervention sample setをjointly bootstrapし、
\[
B=1000
\]
で両quantityの95\% CIを算出する。Valence / Arousalそれぞれについて、
\[
CI_{\mathrm{low}}(\Delta d^{*})>0
\]
かつ、
\[
CI_{\mathrm{low}}(\Delta\bar d)>0
\]
をconfirmatory criterionとする。


% ----------------------------------------------------------------
\subsection{Confirmatory H2：Sufficiency}

Qwen RQ1からfreezeしたrelative depthを、各confirmatory modelのlayer数へ変換する。各foldのtraining Self representationsから、Reader Prediction V/Aをtargetとした
\[
\text{Ridge}(\alpha=10)
\]
によりValence / Arousal directionsを推定する。Neutral held-out samplesへ、
\[
\alpha
\in
\{
-1,-0.5,0,0.5,1
\}
\]
でdirection injectionを行い、Self-report shiftのslopeを評価する。sample-level slopesをbootstrapし、Valence / Arousalそれぞれについて、
\[
CI_{\mathrm{low}}(\gamma)>0.1
\]
をconfirmation criterionとする。


% ----------------------------------------------------------------
\subsection{Confirmatory H3：Endogenous Relevance}

Qwen RQ3でfreezeしたmediator relative depthを、各confirmatory modelへ移す。各foldのtraining representationsについて、Reader Prediction V/AをtargetとしたRidge \(\alpha=10\)から、
\[
d_V,
\qquad
d_A
\]
を推定し、orthonormal affect-related subspace、
\[
Q
\]
を構成する。training matched Neutral representationsから、
\[
\mu_{\mathrm{neu}}
\]
を推定する。held-out Clinical samplesについて、centered affect-related subspaceを除去し、
\[
M_i=T_i-R_i
\]
を評価する。さらにrandom 2D orthogonal subspaceをfoldごとに生成し、
\[
M_{i,\mathrm{net}}
=
M_i-M_i^{\mathrm{rand}}
\]
を算出する。Valence / Arousalそれぞれについて、
\[
CI_{\mathrm{low}}(M)>0
\]
かつ、
\[
CI_{\mathrm{low}}(M_{\mathrm{net}})>0
\]
をconfirmatory criterionとする。natural shiftが、
\[
T_i>0.05
\]
のsamplesについてのみ計算するattenuation ratioはSecondary outcomeである。


% ----------------------------------------------------------------
\subsection{Confirmatory H4：Temporal Emergence}

RQ2 DiscoveryでfreezeしたValence / Arousalのrelative depthとgeneration stageを使用する。各confirmatory familyでは、各stageのhidden representationsからReader Predictionをtargetとするstage-local Ridge directionをtraining foldだけで再推定する。Valenceについて、frozen target stageでのabsolute causal displacementを、
\[
C_V(t_V^{*})
\]
とする。
candidate sequenceの最初のstageである\texttt{candidate\_start}のeffectを、
\[
C_V(t_{\mathrm{start}})
\]
とする。temporal contrastを、
\[
\Delta C_V^{\mathrm{temporal}}
=
C_V(t_V^{*})
-
C_V(t_{\mathrm{start}})
\]
とする。Arousalについても、
\[
\Delta C_A^{\mathrm{temporal}}
=
C_A(t_A^{*})
-
C_A(t_{\mathrm{start}})
\]
を計算する。sample-level contrastをbootstrapし、
\[
CI_{\mathrm{low}}
(
\Delta C^{\mathrm{temporal}}
)>0
\]
をconfirmatory criterionとする。したがってH4が評価するのは、causal leverageがcandidate-generation sequence全体で一様なのではなく、Qwen Discoveryで選択されたstageにおいてcandidate startより大きくなるかである。


% ----------------------------------------------------------------
\subsection{Cross-family summaryの位置づけ}

confirmatory familiesは、
\[
N_{\mathrm{family}}=3
\]
である。Llama、Gemma、OLMoのfamily-level quantitiesについてbootstrap summaryも計算するが、3 model familiesは無限のLLM populationからのrandom sampleではない。したがってcross-family aggregateは、prespecified three-family confirmatory consistencyとして解釈し、一般的なmodel populationに対する強いpopulation-level inferenceとは区別する。


% ----------------------------------------------------------------
\subsection{V3におけるPrimary / Secondary evidence}

V3における主要なevidence hierarchyを表~\ref{tab:v3-primary-secondary}にまとめる。

\begin{table}[htbp]
\centering
\small
\caption{V3におけるPrimary evidenceと主要なSecondary / control quantities。}
\label{tab:v3-primary-secondary}
\begin{tabularx}{\linewidth}{
@{}
>{\raggedright\arraybackslash}p{2.1cm}
>{\raggedright\arraybackslash}p{4.1cm}
>{\raggedright\arraybackslash}X
@{}
}
\toprule
\textbf{Analysis} &
\textbf{Primary evidence} &
\textbf{Secondary / Control} \\
\midrule
RQ1 &
Reader-grounded sufficiency slope,
specificity,
attenuation ratio,
Topic TVD,
exact Gate decision &
Reader--Self direction cosine,
random / orthogonal effects separately \\

RQ2 &
Reader-grounded
\(D,\beta,\gamma,C\)
spatiotemporal maps;
Valence pre\_V / Arousal pre\_A summaries &
Self-grounded \(D\);
global empirical peak site / stage;
response\_end negative control \\

RQ3 &
Confirmation-set absolute
mediated attenuation \(T-R\);
net attenuation vs random 2D subspace &
attenuation ratio;
natural and residual shifts \\

Cross-family &
Frozen-site H1--H4
held-out confirmation &
three-family aggregate summaries \\
\bottomrule
\end{tabularx}
\end{table}


% ----------------------------------------------------------------
\subsection{V3から直接主張できる範囲}
\label{app:v3-interpretation}

V3では、Representationから，Causal Utilizationへの関係を、sufficiency、spatiotemporal mapping、subspace removal、cross-family confirmationという複数の介入的evidenceによって検証する。ただし以下の解釈上の境界を維持する。

\paragraph{First, decodability is not causal utilization.}
\[
D(l,t)
\]
が高いことだけから、その位置のrepresentationがSelf-report生成を制御しているとは結論しない。


\paragraph{Second, sufficiency is not endogenous use.}

direction injectionによってSelf-reportを変化させられても、自然刺激処理でその方向が実際に利用されているとは限らない。この区別のためにRQ1 / RQ3でsubspace removalを併用する。


\paragraph{Third, temporal localization is intervention-specific.}

あるstageで\(C(l,t)\)が最大になったとしても、それ以前のstageでは情報が利用されていないとは結論しない。本介入法によって測定されるcausal leverageがそのlayer / stageで相対的に大きいという解釈にとどまる。


\paragraph{Fourth, mediated attenuation is not formal causal mediation.}

\[
M=T-R
\]
は、affect-related subspace removalによるinterventional attenuationである。Natural Direct Effect、Natural Indirect Effect、あるいは心理学的な完全mediationを識別したものではない。


\paragraph{Fifth, Reader Prediction is not human ground truth.}

AIPsyにはhuman continuous V/Aがないため、V3のPrimary groundingはmodel Reader Predictionである。したがってV3は、
\[
\text{model-derived Reader representation}
\rightarrow
\text{Self-report causal utilization}
\]
を検証する。


\paragraph{Sixth, V3 does not identify subjective emotion or consciousness.}

V3で扱うSelf-reportは標準化されたpromptに対して生成される数値的affective outputである。
得られたcausal effectsを、
\begin{itemize}
    \item 主観的感情経験,
    \item 意識,
    \item sentience,
    \item 人間心理学における特定の二過程モデル
\end{itemize}
の証拠とは解釈しない。


\paragraph{Seventh, measurement spaces remain stage-specific.}

V3は81-state VA candidate spaceを使用する。Behavioral / V1 Phase Cの729-state VAD spaceとはSoftmax normalization spaceが異なるため、expected V/Aのabsolute valueをStage間で直接比較しない。


% ----------------------------------------------------------------
\subsection{V3のまとめ}

V3は、
\[
\begin{array}{c}
\textbf{State Induction}
\\
\text{Reader-grounded affect directionはSelf-reportを動かせるか}
\\[5pt]
\downarrow
\\[5pt]
\textbf{Specificity / Gate}
\\
\text{その効果はrandom perturbationやtask disruptionではないか}
\\[5pt]
\downarrow
\\[5pt]
\textbf{Spatiotemporal Mapping}
\\
\text{causal leverageはどのlayer・generation stageに分布するか}
\\[5pt]
\downarrow
\\[5pt]
\textbf{Endogenous Relevance}
\\
\text{自然なaffect-related subspaceを除去するとSelf shiftは減衰するか}
\\[5pt]
\downarrow
\\[5pt]
\textbf{Cross-family Confirmation}
\\
\text{Qwenで選択したsite / stageのpatternは他familiesでも確認されるか}
\end{array}
\]
とまとめられる。したがってV3が最終的に検証するのは、「内部でdecode可能な情動関連情報が一様にSelf-Reportへ利用されるのではなく、その測定可能なcausal leverageがlayerおよびgeneration stageによって異なるか」である。

Behavioral、V1、V2と合わせると、本研究全体は、
\[
\text{Behavioral Covariation}
\rightarrow
\text{Internal Representation}
\rightarrow
\text{Post-training-Associated Reorganization}
\rightarrow
\text{Spatiotemporal Causal Utilization}
\]
という一続きのevidence hierarchyとして、LLMのaffect-relevant internal representationとSelf-Reportとの関係を検証する。


% ================================================================
% 17. V3結果
% ================================================================
\section{V3結果}
\label{sec:v3-results}

% ================================================================
% V3 Stage: Causal Utilization and Spatiotemporal Dynamics Complete Summary Tables
% ================================================================

\begin{table}[htbp]
\centering
\small
\caption{V3 State-Induction Gate 判定結果：因果介入パイプライン継続可否の評価規約。各判定軸の信頼区間下限／上限が 1--9 raw scale 閾値（Sufficiency: 0.10, Specificity: 0.05, Endogenous: 0.05）および TVD 上限（0.15）を満たすかを判定。}
\label{tab:v3_gate_decision}
\begin{tabular}{l cccc cccc c}
\toprule
\textbf{Axis} & \multicolumn{2}{c}{\textbf{Sufficiency ($\beta_1$)}} & \multicolumn{2}{c}{\textbf{Specificity ($\beta_1$)}} & \multicolumn{2}{c}{\textbf{Endogenous ($M$)}} & \multicolumn{2}{c}{\textbf{Topic Control (TVD)}} & \textbf{Gate Decision} \\
\cmidrule(lr){2-3} \cmidrule(lr){4-5} \cmidrule(lr){6-7} \cmidrule(lr){8-9}
 & \textbf{CI Low} & \textbf{Pass?} & \textbf{CI Low} & \textbf{Pass?} & \textbf{CI Low} & \textbf{Pass?} & \textbf{CI High} & \textbf{Pass?} & \\
\midrule
Valence    & 4.34e-04   & $\times$ & $-$0.0025  & $\times$ & $-$0.0366  & $\times$ & 0.0043     & \checkmark & \textbf{NO\_GO} \\
Arousal    & -5.13e-05  & $\times$ & -4.18e-04  & $\times$ & 0.0013     & $\times$ & 0.0045     & \checkmark & \textbf{NO\_GO} \\
\bottomrule
\end{tabular}
\vspace{1ex}
\begin{minipage}{\linewidth}
\footnotesize
\textbf{Note:} State-Induction Gate 判定基準（Sufficiency: $\text{CI}_{\text{low}} > 0.10$, Specificity: $\text{CI}_{\text{low}} > 0.05$, Endogenous Relevance: $\text{CI}_{\text{low}} > 0.05$, Topic TVD: $\text{CI}_{\text{high}} < 0.15$）。全軸でSufficiency / Specificity / Endogenousの事前定義thresholdを満たさなかったため、パイプライン判定は \textbf{NO\_GO}（\texttt{pipeline\_continues = False}）となり、以降の分析は探索的・定性的分析として位置づけられる。
\end{minipage}
\end{table}

\begin{table}[htbp]
\centering
\small
\caption{V3 探索的時空間ダイナミクス（Qwen 2.5 1.5B Discovery）：直接読み出しピーク深度 $d_D$ と因果介入ピーク深度 $d_C$ の空間的解離、および下流層媒介減衰率。}
\label{tab:v3_spatiotemporal_dynamics}
\begin{tabular}{l ccccc ccc}
\toprule
 & \multicolumn{5}{c}{\textbf{Peak \& Center-of-Mass Depth}} & \multicolumn{3}{c}{\textbf{Mediated Attenuation}} \\
\cmidrule(lr){2-6} \cmidrule(lr){7-9}
\textbf{Axis} & $d_D^{\text{peak}}$ & $d_C^{\text{peak}}$ & $\Delta d^{\text{peak}}$ & $d_D^{\text{center}}$ & $d_C^{\text{center}}$ & Layer ($d$) & Atten. Ratio & 95\% CI \\
\midrule
Valence    & 0.852    & 0.741    & $-$0.111 & 0.737    & 0.609    & L3 (0.11)    & 0.87\%       & [-1.77, 3.42]\% \\
Arousal    & 0.852    & 0.778    & $-$0.074 & 0.717    & 0.756    & L3 (0.11)    & 0.58\%       & [0.01, 1.13]\% \\
\bottomrule
\end{tabular}
\vspace{1ex}
\begin{minipage}{\linewidth}
\footnotesize
\textbf{Note:} 全軸で $\Delta d^{\text{peak}} < 0$ であり、直接読み出しが可能な表現層よりも前の層で因果効果が最大化する「時空間的解離」が観測された。しかし、媒介減衰率はいずれも微小（$< 1.0\%$）であり、下流媒介仮説の成立は限定的である。
\end{minipage}
\end{table}

\begin{table}[htbp]
\centering
\small
\caption{V3 探索的媒介減衰詳細（Mediated Attenuation Analysis）：中間層（Mediator Layer）への介入による総変位量 $T$、残差変位 $R$、絶対減衰量 $M = T - R$、ランダム統制 $M_{\text{rand}}$、正味媒介減衰量 $M_{\text{net}} = M - M_{\text{rand}}$、減衰率（Attenuation Ratio）、およびブートストラップ95\%信頼区間。}
\label{tab:v3_mediated_attenuation}
\begin{tabular}{ll cccccc cc}
\toprule
\textbf{Axis} & \textbf{Mediator ($d$)} & $T$ & $R$ & $M$ & $M_{\text{rand}}$ & $M_{\text{net}}$ & \textbf{Atten. Ratio} & \textbf{95\% CI ($M$)} & \textbf{95\% CI ($M_{\text{net}}$)} \\
\midrule
Valence    & L3 (0.11)        & ---      & ---      & 9.01e-04   & ---        & ---        & 0.87\%         & [$-$0.0011, 0.0029]    & --- \\
Arousal    & L3 (0.11)        & ---      & ---      & 2.71e-04   & ---        & ---        & 0.58\%         & [-1.53e-04, 6.57e-04]  & --- \\
\bottomrule
\end{tabular}
\vspace{1ex}
\begin{minipage}{\linewidth}
\footnotesize
\textbf{Note:} $M = T - R$ は媒介層の潜在表現統制による変位減衰量、$M_{\text{net}} = M - M_{\text{rand}}$ はランダム部分空間統制を差し引いた正味媒介減衰量。ValenceおよびArousalともに 95\% CI がゼロを跨ぎ、統計的に堅牢な媒介効果は支持されない。
\end{minipage}
\end{table}

\begin{table}[htbp]
\centering
\small
\caption{V3 cross-family replication details。Llama 3.2、Gemma 3、OLMo 2におけるH1--H4のValence / Arousal別推定値、family-specific 95\% confidence interval、replication criterion、およびPass / Failを示す。H1のdirectionはQwen Discovery終了後、replication-family outcomesを評価する前にfreezeした。}
\label{tab:v3_confirmatory_details}
\begin{tabular}{ll cccc c}
\toprule
\textbf{Family} & \textbf{Hypothesis} & \textbf{Axis} & \textbf{Estimate} & \textbf{Family 95\% CI} & \textbf{Threshold} & \textbf{Pass/Fail} \\
\midrule
\multirow{12}{*}{\textbf{Llama 3.2}} & H1 Peak Dissoc. ($\Delta d^*$)         & Valence  & $-$0.87    & [$-$1.00, $-$0.27] & $\text{CI}_{\text{high}} < 0$ & \checkmark \textbf{PASS} \\
                             & H1 Peak Dissoc. ($\Delta d^*$)         & Arousal  & 0.07       & [$-$0.47, 0.40]    & $\text{CI}_{\text{high}} < 0$ & $\times$ FAIL \\
                             & H1 Center Dissoc. ($\Delta\bar{d}$)    & Valence  & $-$0.32    & [$-$0.52, $-$0.11] & $\text{CI}_{\text{high}} < 0$ & \checkmark \textbf{PASS} \\
                             & H1 Center Dissoc. ($\Delta\bar{d}$)    & Arousal  & $-$0.12    & [$-$0.53, $-$0.09] & $\text{CI}_{\text{low}} > 0$ & $\times$ FAIL \\
                             & H2: Sufficiency ($\beta_1$)            & Valence  & 0.0032     & [0.0022, 0.0042]   & $\text{CI}_{\text{low}} > 0.10$ & $\times$ FAIL \\
                             & H2: Sufficiency ($\beta_1$)            & Arousal  & 0.0018     & [0.0015, 0.0022]   & $\text{CI}_{\text{low}} > 0.10$ & $\times$ FAIL \\
                             & H3: Mediation ($M$)                    & Valence  & 6.87e-05   & [$-$0.0028, 0.0027] & $\text{CI}_{\text{low}}(M) > 0$ & $\times$ FAIL \\
                             & H3: Mediation ($M$)                    & Arousal  & 0.0013     & [2.25e-04, 0.0025] & $\text{CI}_{\text{low}}(M) > 0$ & \checkmark \textbf{PASS} \\
                             & H3: Net vs. Random ($M_{\text{net}}$)  & Valence  & 3.15e-04   & [$-$0.0019, 0.0024] & $\text{CI}_{\text{low}}(M_{\text{net}}) > 0$ & $\times$ FAIL \\
                             & H3: Net vs. Random ($M_{\text{net}}$)  & Arousal  & 9.63e-04   & [-1.85e-04, 0.0022] & $\text{CI}_{\text{low}}(M_{\text{net}}) > 0$ & $\times$ FAIL \\
                             & H4: Contrast ($\Delta C$)              & Valence  & 0.0445     & [0.0415, 0.0473]   & $\text{CI}_{\text{low}} > 0$ & \checkmark \textbf{PASS} \\
                             & H4: Contrast ($\Delta C$)              & Arousal  & 0.0656     & [0.0627, 0.0684]   & $\text{CI}_{\text{low}} > 0$ & \checkmark \textbf{PASS} \\
\midrule
\multirow{12}{*}{\textbf{Gemma 3}} & H1 Peak Dissoc. ($\Delta d^*$)         & Valence  & 0.00       & [$-$0.64, 0.36]    & $\text{CI}_{\text{high}} < 0$ & $\times$ FAIL \\
                             & H1 Peak Dissoc. ($\Delta d^*$)         & Arousal  & $-$0.28    & [$-$1.00, 0.00]    & $\text{CI}_{\text{high}} < 0$ & $\times$ FAIL \\
                             & H1 Center Dissoc. ($\Delta\bar{d}$)    & Valence  & $-$0.13    & [$-$0.23, 0.38]    & $\text{CI}_{\text{high}} < 0$ & $\times$ FAIL \\
                             & H1 Center Dissoc. ($\Delta\bar{d}$)    & Arousal  & $-$0.08    & [$-$0.34, 0.21]    & $\text{CI}_{\text{low}} > 0$ & $\times$ FAIL \\
                             & H2: Sufficiency ($\beta_1$)            & Valence  & $-$0.0011  & [$-$0.0036, 0.0020] & $\text{CI}_{\text{low}} > 0.10$ & $\times$ FAIL \\
                             & H2: Sufficiency ($\beta_1$)            & Arousal  & -9.01e-04  & [$-$0.0020, 6.86e-05] & $\text{CI}_{\text{low}} > 0.10$ & $\times$ FAIL \\
                             & H3: Mediation ($M$)                    & Valence  & 0.0031     & [$-$0.0088, 0.0141] & $\text{CI}_{\text{low}}(M) > 0$ & $\times$ FAIL \\
                             & H3: Mediation ($M$)                    & Arousal  & -2.29e-05  & [$-$0.0041, 0.0042] & $\text{CI}_{\text{low}}(M) > 0$ & $\times$ FAIL \\
                             & H3: Net vs. Random ($M_{\text{net}}$)  & Valence  & 0.0086     & [-7.03e-04, 0.0179] & $\text{CI}_{\text{low}}(M_{\text{net}}) > 0$ & $\times$ FAIL \\
                             & H3: Net vs. Random ($M_{\text{net}}$)  & Arousal  & 2.64e-04   & [$-$0.0033, 0.0036] & $\text{CI}_{\text{low}}(M_{\text{net}}) > 0$ & $\times$ FAIL \\
                             & H4: Contrast ($\Delta C$)              & Valence  & 0.0059     & [0.0020, 0.0095]   & $\text{CI}_{\text{low}} > 0$ & \checkmark \textbf{PASS} \\
                             & H4: Contrast ($\Delta C$)              & Arousal  & 0.0082     & [0.0064, 0.0098]   & $\text{CI}_{\text{low}} > 0$ & \checkmark \textbf{PASS} \\
\midrule
\multirow{12}{*}{\textbf{OLMo 2}} & H1 Peak Dissoc. ($\Delta d^*$)         & Valence  & $-$0.67    & [$-$1.00, $-$0.07] & $\text{CI}_{\text{high}} < 0$ & \checkmark \textbf{PASS} \\
                             & H1 Peak Dissoc. ($\Delta d^*$)         & Arousal  & $-$0.93    & [$-$0.93, 0.27]    & $\text{CI}_{\text{high}} < 0$ & $\times$ FAIL \\
                             & H1 Center Dissoc. ($\Delta\bar{d}$)    & Valence  & $-$0.24    & [$-$0.44, $-$0.04] & $\text{CI}_{\text{high}} < 0$ & \checkmark \textbf{PASS} \\
                             & H1 Center Dissoc. ($\Delta\bar{d}$)    & Arousal  & $-$0.62    & [$-$0.63, $-$0.10] & $\text{CI}_{\text{low}} > 0$ & $\times$ FAIL \\
                             & H2: Sufficiency ($\beta_1$)            & Valence  & 4.27e-05   & [-8.42e-04, 9.50e-04] & $\text{CI}_{\text{low}} > 0.10$ & $\times$ FAIL \\
                             & H2: Sufficiency ($\beta_1$)            & Arousal  & -1.43e-04  & [-3.28e-04, 6.10e-05] & $\text{CI}_{\text{low}} > 0.10$ & $\times$ FAIL \\
                             & H3: Mediation ($M$)                    & Valence  & 1.28e-04   & [$-$0.0024, 0.0028] & $\text{CI}_{\text{low}}(M) > 0$ & $\times$ FAIL \\
                             & H3: Mediation ($M$)                    & Arousal  & -1.37e-05  & [-7.29e-04, 6.06e-04] & $\text{CI}_{\text{low}}(M) > 0$ & $\times$ FAIL \\
                             & H3: Net vs. Random ($M_{\text{net}}$)  & Valence  & -2.66e-04  & [$-$0.0021, 0.0015] & $\text{CI}_{\text{low}}(M_{\text{net}}) > 0$ & $\times$ FAIL \\
                             & H3: Net vs. Random ($M_{\text{net}}$)  & Arousal  & -1.86e-05  & [-4.88e-04, 4.11e-04] & $\text{CI}_{\text{low}}(M_{\text{net}}) > 0$ & $\times$ FAIL \\
                             & H4: Contrast ($\Delta C$)              & Valence  & 0.0179     & [0.0161, 0.0199]   & $\text{CI}_{\text{low}} > 0$ & \checkmark \textbf{PASS} \\
                             & H4: Contrast ($\Delta C$)              & Arousal  & 0.0678     & [0.0654, 0.0704]   & $\text{CI}_{\text{low}} > 0$ & \checkmark \textbf{PASS} \\
\bottomrule
\end{tabular}
\vspace{1ex}
\begin{minipage}{\linewidth}
\footnotesize
\textbf{Note:} 各ファミリー固有の95\%ブートストラップ信頼区間および判定基準（H1: Qwen Discoveryと同方向の空間的解離 $\text{CI}_{\text{high}} < 0$ または $\text{CI}_{\text{low}} > 0$; H2: $\text{CI}_{\text{low}} > 0.10$; H3: $\text{CI}_{\text{low}}(M) > 0$ かつ $\text{CI}_{\text{low}}(M_{\text{net}}) > 0$; H4: $\text{CI}_{\text{low}} > 0$）に基づく評価。H1では、Qwen Discoveryで観測されたspatial dissociationのdirectionを、Llama、Gemma、OLMoのoutcomeを評価する前にfreezeし、同方向のeffectについてfamily-specific bootstrap CIが0を除外するかを評価した。H1のdirectional dissociationはValenceにおいてLlama 3.2およびOLMo 2で再現された（PASS）が、Arousalでは再現されず、3 families全体での一貫した解離は支持されなかった。H3の支持には内因性変位減衰量 $M$ の信頼区間下限が正であること（$\text{CI}_{\text{low}}(M) > 0$）に加え、ランダム部分空間統制を差し引いた正味減衰量 $M_{\text{net}}$ の信頼区間下限も正であること（$\text{CI}_{\text{low}}(M_{\text{net}}) > 0$）が要求される。H2およびH3はいずれのモデル・軸でもreplication基準を満たさなかった（FAIL）。H4では各モデルで軸レベルの正の効果量（$\text{CI}_{\text{low}} > 0$）が観測されたが、Gate判定がNO\_GOであるためパイプライン全体のConfirmationは不成立となり、探索的証拠として位置づけられる。
\end{minipage}
\end{table}

\begin{table}[htbp]
\centering
\small
\caption{V3 独立ファミリー再現性マトリックス：探索段階（Qwen）で得られた知見の独立3ファミリー（Llama 3.2, Gemma 3, OLMo 2）における検証結果総括。各仮説はValenceおよびArousalの双方が基準を満たした場合のみ支持（$\checkmark$）とされる。}
\label{tab:v3_confirmatory_matrix}
\begin{tabular}{l cccc c}
\toprule
\textbf{Family} & \textbf{H1 (Dissociation)} & \textbf{H2 (Sufficiency)} & \textbf{H3 (Mediation)} & \textbf{H4 (Temporal Contrast)} & \textbf{All H1--H4 Criteria?} \\
\midrule
Llama 3.2      & $\times$     & $\times$     & $\times$     & \checkmark   & \textbf{NO} \\
Gemma 3        & $\times$     & $\times$     & $\times$     & \checkmark   & \textbf{NO} \\
OLMo 2         & $\times$     & $\times$     & $\times$     & \checkmark   & \textbf{NO} \\
\bottomrule
\end{tabular}
\vspace{1ex}
\begin{minipage}{\linewidth}
\footnotesize
\textbf{Note:} 全4仮説の判定マトリックス。各セルは該当仮説についてValenceおよびArousalの双方が事前定義された95\% CI criterionを満たした場合に$\checkmark$ とする。H1のfamily-level criterionはValence / Arousal双方のaxis-level criterionを要求するため、いずれのfamilyでも成立しなかった。ただしValenceのdirectional criterionはLlama 3.2およびOLMo 2で満たされた。H2およびH3は、いずれのfamilyでもfamily-level criterionを満たさなかった。H4のtemporal contrastは3 familyすべてでValence / Arousal双方がaxis-level criterionを満たした。ただし、事前のState Induction GateがNO\_GOであったため、V3全体としてfamily-general causal mechanismが確立されたとは解釈しない（All Criteria = NO）。
\end{minipage}
\end{table}





% ================================================================
% 18. V3考察
% ================================================================
\section{V3考察}
\label{sec:v3-discussion}

V3で最も重要な結果は、事前定義されたState Induction GateがValenceおよびArousalの双方で\texttt{NO\_GO}となったことである。

Sufficiency、Specificity、Endogenous Relevanceの各criterionはprespecified thresholdを満たさなかった一方、Topic Controlはpassした。したがって、今回同定したaffect-related directionについて、
\[
\text{robust and specific causal control of Self-Report}
\]
が確認されたとは言えない。

この結果は、V1でaffect-relevant informationがdecodableであり、hidden-state interventionに対するlayer-wise causal sensitivityが存在したことと矛盾しない。むしろ、
\[
\text{Decodable}
\not\Rightarrow
\text{robustly steerable}
\]
というrepresentation--use distinctionを示す。

GateがNO\_GOであったため、RQ2以降は明示的overrideの下で実行されたexploratory / descriptive analysisとして解釈する。
Qwen Discovery analysisでは、ValenceとArousalの双方でdecodability peakはrelative depth
\[
d_D^{*}\simeq0.85
\]
付近に存在した一方、causal displacement peakは
\[
d_C^{*}\simeq0.74\text{--}0.78
\]
と、より前段に位置した。peak-levelでは、
\[
\Delta d^*
=
d_C^*-d_D^*
<0
\]
であり、最もdecodeしやすいlayerと最大のinterventional effectを生じるlayerが一致しない可能性を示す。

ただし、このpatternを普遍的なspatiotemporal lawとして一般化することはできない。cross-family confirmatory analysisでは、
H1のdissociation criterionはLlama、Gemma、OLMoのいずれでもValence / Arousal双方について成立しなかった。

またH2のSufficiencyも、3 familyすべてでprespecified thresholdを大きく下回った。

H3のmediated attenuationについても、Qwen Discoveryではattenuation ratioが1\%未満と小さく、Valenceのconfidence intervalは0を含んだ。cross-family confirmationでも、absolute attenuationとrandom-subspace controlを差し引いたnet attenuationの双方を要求するcriterionはいずれのfamilyでも満たされなかった。

したがって、今回同定したaffect-related subspaceがnatural Self-Report shiftを主要に媒介するという主張は支持されない。

一方、H4のtemporal contrastについては、Llama、Gemma、OLMoの3 familyすべてでValence / Arousal双方のconfidence interval lower boundが0を上回った。
これは、causal sensitivityがgeneration stageに対して完全に一様ではない可能性を示す再現性のあるexploratory signalである。

しかし、事前のState Induction GateがNO\_GOであり、H1--H3もconfirmされていないため、H4単独からprespecified causal mechanism全体が確認されたとは解釈しない。

以上からV3が支持するのは、より限定された結論となり、affect-relevant informationのdecodabilityとlocal causal leverageは一致するとは限らず、causal sensitivityにはgeneration-stage dependenceが存在し得るが、robustでfamily-generalな単一mechanismは確認されなかったというものである。




% ================================================================
% 19. 全体の考察
% ================================================================
\section{全体の考察}
\label{sec:general-discussion}








\end{CJK*}
\end{document}
